% !TEX program = lualatex
% CCAR-F 英単語帳（公開版）— gen_vocab.py の自動生成物。手で直さない。
\documentclass[a4paper,10pt]{ltjsarticle}
\usepackage[top=14mm,bottom=16mm,left=14mm,right=14mm]{geometry}

\usepackage{luatexja-fontspec}
\setmainfont{TeX Gyre Pagella}
\setsansfont{TeX Gyre Heros}
\setmainjfont{Noto Serif CJK JP}
\setsansjfont{Noto Sans CJK JP}
% 発音記号用（IPA 拡張を全部持っているフォント。TeX Live 同梱の DejaVu を使う）
\newfontfamily\ipafont{DejaVu Sans}[Scale=0.94]
% θ などギリシャ文字を和文扱いにしない（既定では range 2 が JAchar になり、発音記号の中で字間が開く）
\ltjsetparameter{jacharrange={-2}}

\usepackage{xcolor}
\usepackage{graphicx}
\usepackage{longtable}
\usepackage{booktabs}
\usepackage{array}
\usepackage{enumitem}
\usepackage{tcolorbox}
\usepackage{fancyhdr}

\definecolor{vhead}{HTML}{1F3A5F}
\definecolor{vrule}{HTML}{9AA7B5}
\definecolor{vkey}{HTML}{B3261E}
\definecolor{vnote}{HTML}{8A6D1F}
\definecolor{vband}{HTML}{EDF1F6}
\definecolor{vgray}{HTML}{5A5A5A}

% 見出し
\usepackage{needspace}
\usepackage{titlesec}
\titleformat{\section}{\sffamily\large\bfseries\color{vhead}}{}{0pt}{}
\titlespacing*{\section}{0pt}{1.6em}{0.6em}
\titleformat{\subsection}{\sffamily\normalsize\bfseries\color{vhead}}{}{0pt}{}
\titlespacing*{\subsection}{0pt}{1.0em}{0.4em}
% titlesec が設定した見出し処理に、hyperref のリンク先処理を適用する。
\usepackage[hidelinks]{hyperref}

\pagestyle{fancy}
\fancyhf{}
\fancyhead[L]{\footnotesize\sffamily\color{vgray}CCAR-F 英単語帳}
\fancyhead[R]{\footnotesize\sffamily\color{vgray}\thepage}
\renewcommand{\headrulewidth}{0.2pt}

\setcounter{secnumdepth}{0}
\setlength{\parindent}{0pt}
\setlength{\parskip}{0.4\baselineskip}
\setlist{itemsep=0.15em, parsep=0pt, topsep=0.25em, leftmargin=1.2em}

% 英語見出し語と、その下に置く発音記号
\newcommand{\en}[1]{\textsf{\textbf{#1}}}
\newcommand{\ipa}{\ipafont\color{vgray}}
\newcommand{\pos}[1]{{\sffamily\mdseries\scriptsize\color{vhead}#1}}

% 語彙表
\newlength{\vocabrowspace}
\newenvironment{vtab}{%
  \footnotesize
  % 表の中では baselineskip が変わるため、本文1行分を先に保存する。
  \setlength{\vocabrowspace}{\baselineskip}
  \setlength{\tabcolsep}{4pt}
  \renewcommand{\arraystretch}{1.18}
  \begin{longtable}{@{}>{\raggedright}p{0.265\textwidth}>{\raggedright}p{0.205\textwidth}>{\raggedright\arraybackslash}p{0.475\textwidth}@{}}
  \toprule
  \sffamily\bfseries 英語 & \sffamily\bfseries 品詞・意味 & \sffamily\bfseries 使い方・オリジナル例文 \\
  \midrule
  \endfirsthead
  \toprule
  \sffamily\bfseries 英語 & \sffamily\bfseries 品詞・意味 & \sffamily\bfseries 使い方・オリジナル例文 \\
  \midrule
  \endhead
  \bottomrule
  \endfoot
}{%
  \end{longtable}
}

\newtcolorbox{vbx}[1]{colback=vband,colframe=vrule,boxrule=0.3pt,arc=1pt,
  left=5pt,right=5pt,top=4pt,bottom=4pt,
  fonttitle=\sffamily\bfseries\small,title={#1}}


% Stress glyph ink centers are aligned over the vowel nucleus; offsets are em.
\newcommand{\ipastress}[5]{\mbox{\rlap{\kern#4em\raisebox{#5em}{{\color{#1}\bfseries#2}}}#3}}
\newcommand{\ipastrut}{\rule{0pt}{1.45em}}
\begin{document}
\begin{center}
{\sffamily\bfseries\Large\color{vhead} CCAR-F 英単語帳 — 問題文を読み切る}\\[0.4em]
{\footnotesize\sffamily\color{vgray} 公開版 2026-10-11 ｜ 620項目 ｜ 品詞・オリジナル例文・和訳・強勢付きIPA}
\end{center}

\begin{vbx}{発音記号と強勢の読み方}
\small
主強勢 {\ipa {\color{vkey}\bfseries ˈ}}（赤）と副強勢 {\ipa {\color{vhead}\bfseries ˌ}}（青）を、対応する母音の真上に重ねる。印は、その母音を含む音節の強さを示す。発音は原則アメリカ英語の目安。UKは英国発音。\par
二重母音は2文字のまとまりの中央に置く。標準IPAの音節前の印を上へ移した、この単語帳の学習用表記。\par
\en{require}\quad {\ipa \ipastrut /rɪkw\ipastress{vkey}{ˈ}{aɪ}{0.340750}{0.192500}ər/}\quad 後半の二重母音の真上に主強勢。\par
\en{optional}\quad {\ipa \ipastrut /\ipastress{vkey}{ˈ}{ɑ}{0.146500}{0.192500}pʃənəl/}\quad 最初の母音の真上に主強勢。\par
\en{validation}\quad {\ipa \ipastrut /v\ipastress{vhead}{ˌ}{æ}{0.341750}{0.701000}ləd\ipastress{vkey}{ˈ}{eɪ}{0.339250}{0.192500}ʃən/}\quad 最初の母音に副強勢、後半の二重母音に主強勢。\par
単音節語では通常、強勢記号を省く。句のIPAは語の強勢を示し、文全体の抑揚は固定しない。長い表現は主要語のみ、コード識別子は原文の綴りを示す。
\par{\scriptsize 参考: \href{https://dictionary.cambridge.org/help/phonetics.html}{Cambridge Dictionary: Pronunciation symbols}}
\end{vbx}

{\footnotesize 用例・和訳・使い方の説明は、この公開版のために作成した学習用の内容です。本書は非公式の学習資料です。}

\begingroup\setlength{\parskip}{0pt}
\tableofcontents
\endgroup
\clearpage
\Needspace{14\baselineskip}

\section{1. 条件・対比・例外を読む}
\begin{vtab}
\en{rather than}\newline{\ipa \ipastrut /r\ipastress{vkey}{ˈ}{æ}{0.341750}{0.192500}ðər ðæn/} & \pos{前置詞（複数語）}\newline …ではなく・…よりも & 二つの方法を比べ、一方を優先して選ぶときに使う。\newline 用例: We chose a short meeting rather than another email.\newline 訳: 私たちは、もう一通メールを送るよりも短い会議を選んだ。 \\
\addlinespace[\vocabrowspace]
\en{yet}\newline{\ipa \ipastrut /jɛt/} & \pos{接続詞}\newline それでも・しかし & 前の内容から予想される結果と違う事実を続ける。\newline 用例: The office is small, yet everyone has enough space.\newline 訳: 事務所は小さいが、全員に十分なスペースがある。 \\
\addlinespace[\vocabrowspace]
\en{while}\newline{\ipa \ipastrut /waɪl/} & \pos{接続詞}\newline …している間に・一方で & ある動作が続く時間や、二つの事実の対比を表す。\newline 用例: I reviewed the agenda while my colleague prepared the room.\newline 訳: 同僚が部屋を準備している間、私は議題を確認した。\newline 用例: The desktop version is free, while the mobile version costs money.\newline 訳: デスクトップ版は無料だが、モバイル版は有料だ。 \\
\addlinespace[\vocabrowspace]
\en{even though}\newline{\ipa \ipastrut /\ipastress{vkey}{ˈ}{i}{-0.013750}{0.392500}vən ðoʊ/} & \pos{接続表現}\newline …であるにもかかわらず & 事実として成り立っている条件に反する結果を述べる。\newline 用例: She attended the workshop even though her schedule was full.\newline 訳: 予定が詰まっていたにもかかわらず、彼女は研修に参加した。 \\
\addlinespace[\vocabrowspace]
\en{before}\newline{\ipa \ipastrut /bɪf\ipastress{vkey}{ˈ}{ɔ}{0.125500}{0.192500}r/} & \pos{接続詞}\newline …する前に・…より前に & 二つの出来事のうち、先に行うことを示す。\newline 用例: Please save your work before you close the application.\newline 訳: アプリケーションを閉じる前に、作業内容を保存してください。 \\
\addlinespace[\vocabrowspace]
\en{whether}\newline{\ipa \ipastrut /w\ipastress{vkey}{ˈ}{ɛ}{0.115750}{0.193000}ðər/} & \pos{接続詞}\newline …かどうか & まだ決まっていない選択や、確かめる必要のある事柄を表す。\newline 用例: We need to decide whether to extend the survey.\newline 訳: 調査を延長するかどうか決める必要がある。 \\
\addlinespace[\vocabrowspace]
\en{instead of}\newline{\ipa \ipastrut /ɪnst\ipastress{vkey}{ˈ}{ɛ}{0.115750}{0.193000}d əv/} & \pos{前置詞（複数語）}\newline …の代わりに・…ではなく & あるものを別のもので置き換える選択を示す。\newline 用例: Please send a calendar invitation instead of a text message.\newline 訳: テキストメッセージの代わりに、カレンダーの招待を送ってください。 \\
\addlinespace[\vocabrowspace]
\en{at least}\newline{\ipa \ipastrut /æt list/} & \pos{副詞句}\newline 少なくとも & 必要な量や確保できる数について、下限を示す。\newline 用例: The workshop needs at least six people to proceed.\newline 訳: 研修を実施するには、少なくとも六人が必要だ。 \\
\addlinespace[\vocabrowspace]
\en{since}\newline{\ipa \ipastrut /sɪns/} & \pos{接続詞・前置詞・副詞}\newline …なので・…以来／その後・それ以来 & 理由を示すほか、ある時点からの継続や、その後の変化を表す。\newline 用例: Since the road is closed, we will take the train.\newline 訳: 道路が閉鎖されているので、私たちは電車で行く。\newline 用例: The library has offered evening classes since last autumn.\newline 訳: 図書館は昨年の秋から夜間講座を開いている。\newline 用例: We visited the studio once and have since stayed in touch.\newline 訳: 私たちはその工房を一度訪れ、その後も連絡を取り合っている。 \\
\addlinespace[\vocabrowspace]
\en{up to}\newline{\ipa \ipastrut /ʌp tu/} & \pos{前置詞（複数語）}\newline …まで & 許される量や扱う期間が、どこまで届くかを示す。\newline 用例: This small elevator can carry up to eight people.\newline 訳: この小型エレベーターには八人まで乗れる。 \\
\addlinespace[\vocabrowspace]
\en{within}\newline{\ipa \ipastrut /wɪð\ipastress{vkey}{ˈ}{ɪ}{0.032500}{0.179500}n/} & \pos{前置詞}\newline …の内側で・…の範囲内で & 決まった境界や期限、能力の範囲に収まることを表す。\newline 用例: Please return the borrowed equipment within seven working days.\newline 訳: 借りた機材は七営業日以内に返却してください。\newline 用例: The proposed purchase is well within our remaining budget.\newline 訳: 提案された購入品の金額は、残りの予算に余裕をもって収まる。 \\
\addlinespace[\vocabrowspace]
\en{based on}\newline{\ipa \ipastrut /beɪst ɑn/} & \pos{過去分詞句（他動詞由来）}\newline …に基づいて・…を根拠に & 結論や計画を支える根拠が何かを示す。\newline 用例: The opening hours were revised based on visitor feedback.\newline 訳: 来場者の意見に基づいて、開館時間が見直された。 \\
\addlinespace[\vocabrowspace]
\en{without}\newline{\ipa \ipastrut /wɪð\ipastress{vkey}{ˈ}{aʊ}{0.465750}{0.192500}t/} & \pos{前置詞}\newline …なしで・…せずに & 何かを伴わない状態や、ある動作を行わないことを表す。\newline 用例: You can preview the document without downloading a copy.\newline 訳: コピーをダウンロードせずに、文書をプレビューできる。 \\
\addlinespace[\vocabrowspace]
\en{because}\newline{\ipa \ipastrut /bɪk\ipastress{vkey}{ˈ}{ɔ}{0.125500}{0.192500}z/} & \pos{接続詞}\newline …だから・…なので & 行動や判断が生じた理由を具体的に述べる。\newline 用例: We moved the meeting because the room was unavailable.\newline 訳: 部屋が使えなかったため、会議を移動した。 \\
\addlinespace[\vocabrowspace]
\en{only when}\newline{\ipa \ipastrut /\ipastress{vkey}{ˈ}{oʊ}{0.462750}{0.192500}nli wɛn/} & \pos{接続表現}\newline …のときに限り & 動作が起こる時点や状況を、一つの条件に限定する。\newline 用例: The light turns green only when the battery is full.\newline 訳: バッテリーが満充電になったときに限り、ランプが緑になる。 \\
\addlinespace[\vocabrowspace]
\en{roughly a third}\newline{\ipa \ipastrut /r\ipastress{vkey}{ˈ}{ʌ}{0.142750}{0.179500}fli ə θɜrd/} & \pos{数量表現}\newline およそ3分の1 & 全体を三等分した一つ分に近い、おおよその量を示す。\newline 用例: Roughly a third of the workshop participants arrived by bicycle.\newline 訳: 研修参加者のおよそ三分の一が自転車で来た。 \\
\addlinespace[\vocabrowspace]
\en{at each stage}\newline{\ipa \ipastrut /æt itʃ steɪdʒ/} & \pos{前置詞句}\newline 各段階で・一つひとつの段階で & 一連の過程の中で、どの段階でも行うことを示す。\newline 用例: The designer saved a sketch at each stage of development.\newline 訳: デザイナーは制作の各段階でスケッチを保存した。 \\
\addlinespace[\vocabrowspace]
\en{by itself}\newline{\ipa \ipastrut /baɪ ɪts\ipastress{vkey}{ˈ}{ɛ}{0.115750}{0.193000}lf/} & \pos{前置詞句}\newline それだけで・それ自体で & ほかの要素を加えず、そのものだけで成立することを表す。\newline 用例: This small panel can power the sensor by itself.\newline 訳: この小さなパネルだけでセンサーに電力を供給できる。 \\
\addlinespace[\vocabrowspace]
\en{whether ... or ...}\newline{\ipa \ipastrut /w\ipastress{vkey}{ˈ}{ɛ}{0.115750}{0.193000}ðər ... ɔr .../} & \pos{接続表現}\newline …であっても…であっても & 二つの可能性を挙げ、どちらにも当てはまることや選択を表す。\newline 用例: The exhibition remains open whether it rains or snows.\newline 訳: 雨でも雪でも、展覧会は開いている。\newline 用例: We have not decided whether to rent or buy equipment.\newline 訳: 機材を借りるか買うか、私たちはまだ決めていない。 \\
\addlinespace[\vocabrowspace]
\en{despite}\newline{\ipa \ipastrut /dɪsp\ipastress{vkey}{ˈ}{aɪ}{0.340750}{0.192500}t/} & \pos{前置詞}\newline …にもかかわらず & 何かが妨げになりそうでも、実際には別の結果になることを表す。\newline 用例: The outdoor event attracted many visitors despite the cold weather.\newline 訳: 寒い天気にもかかわらず、屋外イベントには多くの来場者が集まった。 \\
\addlinespace[\vocabrowspace]
\en{regardless of}\newline{\ipa \ipastrut /rɪɡ\ipastress{vkey}{ˈ}{ɑ}{0.146500}{0.192500}rdləs əv/} & \pos{前置詞（複数語）}\newline …にかかわらず・…を問わず & 特定の条件が違っていても、扱いや結果が変わらないことを示す。\newline 用例: Everyone receives the same training regardless of previous experience.\newline 訳: 経験の有無にかかわらず、全員が同じ研修を受ける。 \\
\addlinespace[\vocabrowspace]
\en{for as long as}\newline{\ipa \ipastrut /fɔr æz lɔŋ æz/} & \pos{期間を示す表現}\newline …している間ずっと & ある状態が続く期間の全体にわたることを示す。\newline 用例: You may keep the loaner for as long as necessary.\newline 訳: 必要な間はずっと、貸出機を使っていてかまわない。 \\
\addlinespace[\vocabrowspace]
\en{neither ... nor ...}\newline{\ipa \ipastrut /n\ipastress{vkey}{ˈ}{i}{-0.013750}{0.392500}ðər ... nɔr .../} & \pos{並列の否定表現}\newline …も…もどちらもない & 並べた二つのものを、両方とも否定する。\newline 用例: Neither the brochure nor the website lists the opening date.\newline 訳: パンフレットにもウェブサイトにも、開業日は記載されていない。 \\
\addlinespace[\vocabrowspace]
\en{either ... or ...}\newline{\ipa \ipastrut /\ipastress{vkey}{ˈ}{i}{-0.013750}{0.392500}ðər ... ɔr .../} & \pos{並列・選択表現}\newline …か…のどちらか & 二つの候補を示し、そのうち一方を選べることを表す。\newline 用例: You can reserve either a desk or a meeting room.\newline 訳: 机か会議室のどちらかを予約できる。 \\
\addlinespace[\vocabrowspace]
\en{both ... and ...}\newline{\ipa \ipastrut /boʊθ ... ænd .../} & \pos{並列表現}\newline …と…の両方 & 並べた二つのものが、ともに対象に含まれることを示す。\newline 用例: The course covers both practical skills and basic theory.\newline 訳: その講座では実践的な技能と基礎理論の両方を扱う。 \\
\addlinespace[\vocabrowspace]
\en{at most}\newline{\ipa \ipastrut /æt moʊst/} & \pos{副詞句}\newline 多くても・最大で & 数や時間が最大でもその程度に収まることを示す。\newline 用例: The initial setup should take ten minutes at most.\newline 訳: 初期設定は、長くても10分で終わるはずだ。 \\
\addlinespace[\vocabrowspace]
\en{only after}\newline{\ipa \ipastrut /\ipastress{vkey}{ˈ}{oʊ}{0.462750}{0.192500}nli \ipastress{vkey}{ˈ}{æ}{0.341750}{0.192500}ftər/} & \pos{時間の限定表現}\newline …した後になって初めて & ある出来事が済んだ時点で初めて、次のことが起こると示す。\newline 用例: We booked the venue only after the dates were confirmed.\newline 訳: 日程が確定して初めて、会場を予約した。 \\
\addlinespace[\vocabrowspace]
\en{no longer}\newline{\ipa \ipastrut /noʊ l\ipastress{vkey}{ˈ}{ɔ}{0.125500}{0.192500}ŋɡər/} & \pos{副詞句}\newline もはや…ない・もう…しない & 以前続いていた状態や習慣が、今は終わっていることを示す。\newline 用例: The museum no longer sells tickets at the entrance.\newline 訳: その博物館では、もう入口でチケットを販売していない。 \\
\addlinespace[\vocabrowspace]
\en{still}\newline{\ipa \ipastrut /stɪl/} & \pos{副詞}\newline まだ・依然として・それでも & 状態の継続や、障害があっても変わらない結果を表す。\newline 用例: The team is still waiting for the final delivery.\newline 訳: チームはまだ最後の納品を待っている。\newline 用例: The instructions were brief, but we still completed the setup.\newline 訳: 説明は短かったが、それでも私たちは設定を終えた。 \\
\addlinespace[\vocabrowspace]
\en{otherwise}\newline{\ipa \ipastrut /\ipastress{vkey}{ˈ}{ʌ}{0.142750}{0.179500}ðərwaɪz/} & \pos{副詞}\newline ほかの点では・そうでなければ & 特定の点以外の状態や、行動しなかった場合の結果を示す。\newline 用例: The folder has a torn label but is otherwise intact.\newline 訳: フォルダーはラベルが破れているが、それ以外は傷んでいない。\newline 用例: Please bring your badge; otherwise, you cannot enter the building.\newline 訳: バッジを持参してください。そうしないと建物に入れません。 \\
\addlinespace[\vocabrowspace]
\en{once}\newline{\ipa \ipastrut /wʌns/} & \pos{接続詞}\newline いったん…したら・…すると & ある出来事が起きることを、その後の行動のきっかけとして示す。\newline 用例: Once the paint dries, we can install the shelves.\newline 訳: 塗料が乾いたら、棚を取り付けられる。 \\
\addlinespace[\vocabrowspace]
\en{until}\newline{\ipa \ipastrut /ənt\ipastress{vkey}{ˈ}{ɪ}{0.032500}{0.179500}l/} & \pos{接続詞}\newline …までずっと・…するまで & 動作や状態が続く最後の時点を示す。\newline 用例: The reception desk remains open until the last visitor leaves.\newline 訳: 最後の来場者が帰るまで、受付は開いている。 \\
\addlinespace[\vocabrowspace]
\en{in addition to}\newline{\ipa \ipastrut /ɪn əd\ipastress{vkey}{ˈ}{ɪ}{0.032500}{0.179500}ʃən tu/} & \pos{前置詞（複数語）}\newline …に加えて & すでにあるものに、別のものが加わることを表す。\newline 用例: The room includes storage space in addition to a desk.\newline 訳: その部屋には机に加えて収納スペースもある。 \\
\addlinespace[\vocabrowspace]
\en{fail to}\newline{\ipa \ipastrut /feɪl tu/} & \pos{動詞句（自動詞）}\newline …できない・…し損なう & 期待されたことや必要なことが実現しない状況を表す。\newline 用例: Older devices sometimes fail to display the updated calendar.\newline 訳: 古い端末では、更新されたカレンダーが表示されないことがある。 \\
\addlinespace[\vocabrowspace]
\en{so that}\newline{\ipa \ipastrut /soʊ ðæt/} & \pos{接続表現}\newline …するように・…するために & ある行動によって実現したい目的を示す。\newline 用例: We enlarged the labels so that everyone could read them.\newline 訳: 全員が読めるように、私たちはラベルを大きくした。 \\
\addlinespace[\vocabrowspace]
\en{one at a time}\newline{\ipa \ipastrut /wʌn æt ə taɪm/} & \pos{副詞句}\newline 一つずつ・一度に一つ & 複数のものを同時に扱わず、一つずつ順に扱うことを示す。\newline 用例: Please place the slides on the scanner one at a time.\newline 訳: スライドを一枚ずつスキャナーに置いてください。 \\
\addlinespace[\vocabrowspace]
\en{such as}\newline{\ipa \ipastrut /sʌtʃ æz/} & \pos{前置詞（複数語）}\newline 例えば…のような & 広い分類の中から、具体的な例を挙げる。\newline 用例: Bring basic supplies such as paper, pencils, and tape.\newline 訳: 紙、鉛筆、テープなどの基本的な用品を持参してください。 \\
\addlinespace[\vocabrowspace]
\en{whichever}\newline{\ipa \ipastrut /wɪtʃ\ipastress{vkey}{ˈ}{ɛ}{0.115750}{0.193000}vər/} & \pos{限定詞}\newline …する方のどれでも・どちらでも & いくつかの候補の中で、条件に合うどれを選んでもよいことを示す。\newline 用例: Choose whichever seat gives you the clearest view.\newline 訳: 一番見やすい席を、どれでも選んでください。 \\
\addlinespace[\vocabrowspace]
\en{whenever}\newline{\ipa \ipastrut /wɛn\ipastress{vkey}{ˈ}{ɛ}{0.115750}{0.193000}vər/} & \pos{接続詞}\newline …するたびに・…するときはいつでも & ある状況が生じるたびに、決まったことが起こると示す。\newline 用例: The camera saves a picture whenever the door opens.\newline 訳: ドアが開くたびに、カメラが写真を保存する。 \\
\addlinespace[\vocabrowspace]
\en{unless}\newline{\ipa \ipastrut /ənl\ipastress{vkey}{ˈ}{ɛ}{0.115750}{0.193000}s/} & \pos{接続詞}\newline …でない限り & 通常の結果に対して、例外となる条件を示す。\newline 用例: The reservation will expire unless you confirm it today.\newline 訳: 今日中に確認しない限り、その予約は失効する。 \\
\addlinespace[\vocabrowspace]
\en{as soon as}\newline{\ipa \ipastrut /æz sun æz/} & \pos{接続表現}\newline …するとすぐに & 一つの出来事に続いて、すぐ次の行動が起こることを示す。\newline 用例: I will send the slides as soon as editing finishes.\newline 訳: 編集が終わり次第、スライドを送る。 \\
\addlinespace[\vocabrowspace]
\en{only if}\newline{\ipa \ipastrut /\ipastress{vkey}{ˈ}{oʊ}{0.462750}{0.192500}nli ɪf/} & \pos{接続表現}\newline …の場合に限り & 行動や結果を成立させるために欠かせない条件を示す。\newline 用例: The printer works only if both cables are connected.\newline 訳: 両方のケーブルが接続されている場合に限り、プリンターは動く。 \\
\addlinespace[\vocabrowspace]
\en{even if}\newline{\ipa \ipastrut /\ipastress{vkey}{ˈ}{i}{-0.013750}{0.392500}vən ɪf/} & \pos{接続表現}\newline たとえ…でも & 仮に条件が変わっても、判断や結果が変わらないことを示す。\newline 用例: Keep a paper copy even if the connection seems reliable.\newline 訳: 接続が安定しているようでも、紙のコピーを保管してください。 \\
\addlinespace[\vocabrowspace]
\en{whereas}\newline{\ipa \ipastrut /wɛr\ipastress{vkey}{ˈ}{æ}{0.341750}{0.192500}z/} & \pos{接続詞}\newline 一方で・…であるのに対し & 二つの対象に見られる、異なる特徴を並べて述べる。\newline 用例: This room faces east, whereas the larger room faces west.\newline 訳: この部屋は東向きだが、大きい方の部屋は西向きだ。 \\
\addlinespace[\vocabrowspace]
\en{alongside}\newline{\ipa \ipastrut /əl\ipastress{vkey}{ˈ}{ɑ}{0.146500}{0.192500}ŋsaɪd/} & \pos{前置詞}\newline …と一緒に・…と並んで & 二つのものが一緒に置かれたり、同時に進められたりすることを示す。\newline 用例: The guide includes illustrations alongside the written instructions.\newline 訳: そのガイドには文章による説明と一緒に図も載っている。 \\
\addlinespace[\vocabrowspace]
\en{right after}\newline{\ipa \ipastrut /raɪt \ipastress{vkey}{ˈ}{æ}{0.341750}{0.192500}ftər/} & \pos{接続表現（副詞＋接続詞）}\newline …した直後に & ある出来事と次の出来事の間が、ほとんど空かないことを示す。\newline 用例: We cleaned the tables right after the guests left.\newline 訳: 客が帰った直後に、私たちはテーブルを片づけた。 \\
\addlinespace[\vocabrowspace]
\en{ever}\newline{\ipa \ipastrut /\ipastress{vkey}{ˈ}{ɛ}{0.115750}{0.193000}vər/} & \pos{副詞}\newline 一度でも・そもそも & 経験や可能性について、一度でも該当するかを強調する。\newline 用例: Have you ever used this room for a presentation?\newline 訳: この部屋をプレゼンテーションに使ったことが一度でもありますか。 \\
\addlinespace[\vocabrowspace]
\en{so full of ... that ...}\newline{\ipa \ipastrut /soʊ fʊl əv ... ðæt .../} & \pos{結果の構文}\newline …であまりにいっぱいなので、その結果… & 何かで満たされた程度が大きく、そのために結果が生じることを表す。\newline 用例: The drawer was so full of cables that it would not close.\newline 訳: 引き出しはケーブルでいっぱいで、閉まらなかった。 \\
\addlinespace[\vocabrowspace]
\en{across}\newline{\ipa \ipastrut /əkr\ipastress{vkey}{ˈ}{ɔ}{0.125500}{0.192500}s/} & \pos{前置詞}\newline …にわたって・…をまたいで & 複数の場所に広がる範囲や、ある場所を横切る動きを表す。\newline 用例: The new schedule applies across all our regional offices.\newline 訳: 新しい予定は、すべての地域事務所に適用される。\newline 用例: We carried the spare chairs across the courtyard.\newline 訳: 私たちは予備の椅子を運んで中庭を横切った。 \\
\addlinespace[\vocabrowspace]
\en{against}\newline{\ipa \ipastrut /əɡ\ipastress{vkey}{ˈ}{ɛ}{0.115750}{0.193000}nst/} & \pos{前置詞}\newline …を対象に・…に対して & 検証や操作を行うときの、相手や照合対象を示す。\newline 用例: We tested the new client against a local server.\newline 訳: 私たちはローカルサーバーを使って新しいクライアントをテストした。 \\
\addlinespace[\vocabrowspace]
\en{over}\newline{\ipa \ipastrut /\ipastress{vkey}{ˈ}{oʊ}{0.462750}{0.192500}vər/} & \pos{前置詞}\newline …を対象に・…について & 確認や検討が、どの範囲や対象に及ぶかを示す。\newline 用例: The editor made a final pass over the complete manuscript.\newline 訳: 編集者は、原稿全体を対象に最後の確認を行った。 \\
\addlinespace[\vocabrowspace]
\en{no matter what}\newline{\ipa \ipastrut /noʊ m\ipastress{vkey}{ˈ}{æ}{0.341750}{0.192500}tər wʌt/} & \pos{慣用表現}\newline 何がどうであっても・何を…しても & 状況や要求がどう変わっても、方針や結果が一定であることを示す。\newline 用例: Keep the emergency exit clear no matter what happens.\newline 訳: 何が起きても、非常口をふさがないようにしてください。 \\
\addlinespace[\vocabrowspace]
\en{whatever}\newline{\ipa \ipastrut /wɑt\ipastress{vkey}{ˈ}{ɛ}{0.115750}{0.193000}vər/} & \pos{代名詞}\newline …するものは何でも・すべて & 内容を限定せず、該当するものすべてを指す。\newline 用例: Please label whatever you leave in the shared refrigerator.\newline 訳: 共用冷蔵庫に置くものには、何でもラベルを付けてください。 \\
\addlinespace[\vocabrowspace]
\en{no number of}\newline{\ipa \ipastrut /noʊ n\ipastress{vkey}{ˈ}{ʌ}{0.142750}{0.179500}mbər əv/} & \pos{数量表現}\newline どれほど多くの…でも・いくら…しても & 回数や量をいくら増やしても、結果を変えられないことを示す。\newline 用例: No number of reminders can replace a clear agreement.\newline 訳: いくら念押しをしても、明確な合意の代わりにはならない。 \\
\addlinespace[\vocabrowspace]
\en{as if}\newline{\ipa \ipastrut /æz ɪf/} & \pos{接続表現}\newline あたかも…であるかのように・まるで…したかのように & 見え方や振る舞いを、別の状況になぞらえて表す。\newline 用例: He described the city as if he lived there.\newline 訳: 彼は、まるでそこに住んでいるかのようにその街を説明した。 \\
\end{vtab}
\Needspace{14\baselineskip}

\section{2. 動作・判断を読む}
\begin{vtab}
\en{succeed}\newline{\ipa \ipastrut /səks\ipastress{vkey}{ˈ}{i}{-0.013750}{0.392500}d/} & \pos{自動詞}\newline 成功する・うまくいく & 試みや取り組みが目指した結果に結び付くことを表す。\newline 用例: The fundraising campaign succeeded after several local businesses joined.\newline 訳: 地元の企業が数社加わり、募金活動は成功した。 \\
\addlinespace[\vocabrowspace]
\en{address}\newline{\ipa \ipastrut /ədr\ipastress{vkey}{ˈ}{ɛ}{0.115750}{0.193000}s/} & \pos{他動詞}\newline 対処する・取り上げる & 問題を取り上げ、解決に向けて扱うことを表す。\newline 用例: The revised handbook addresses common questions about remote work.\newline 訳: 改訂された手引きは、在宅勤務についてよくある質問を扱っている。 \\
\addlinespace[\vocabrowspace]
\en{require}\newline{\ipa \ipastrut /rɪkw\ipastress{vkey}{ˈ}{aɪ}{0.340750}{0.192500}ər/} & \pos{他動詞}\newline 必要とする・求める & 行動や仕組みを成り立たせるために、何かを必要とする。\newline 用例: The workshop requires participants to bring their own laptops.\newline 訳: その研修では、参加者に自分のノートパソコンの持参を求める。 \\
\addlinespace[\vocabrowspace]
\en{gather}\newline{\ipa \ipastrut /ɡ\ipastress{vkey}{ˈ}{æ}{0.341750}{0.192500}ðər/} & \pos{他動詞}\newline 集める & 必要な資料や情報を、複数の場所から集める。\newline 用例: We gathered visitor feedback during the first exhibition week.\newline 訳: 私たちは展覧会の最初の週に来場者の意見を集めた。 \\
\addlinespace[\vocabrowspace]
\en{omit}\newline{\ipa \ipastrut /oʊm\ipastress{vkey}{ˈ}{ɪ}{0.032500}{0.179500}t/} & \pos{他動詞}\newline 省く・入れ忘れる & 内容から意図的に外すことにも、うっかり抜かすことにも使う。\newline 用例: Please omit personal phone numbers from the public directory.\newline 訳: 公開する名簿からは個人の電話番号を省いてください。\newline 用例: I accidentally omitted the date from the invitation.\newline 訳: 私はうっかり招待状に日付を書き忘れた。 \\
\addlinespace[\vocabrowspace]
\en{revise}\newline{\ipa \ipastrut /rɪv\ipastress{vkey}{ˈ}{aɪ}{0.340750}{0.192500}z/} & \pos{他動詞}\newline 見直す・修正する & 既存の文章や計画を見直し、必要な変更を加える。\newline 用例: We revised the timetable after checking the train connections.\newline 訳: 私たちは電車の接続を確認してから予定表を修正した。 \\
\addlinespace[\vocabrowspace]
\en{remain}\newline{\ipa \ipastrut /rɪm\ipastress{vkey}{ˈ}{eɪ}{0.339250}{0.192500}n/} & \pos{自動詞}\newline 残る・…のままである & 何かが残っていることや、同じ状態が続くことを表す。\newline 用例: A few empty boxes remain in the storage room.\newline 訳: 倉庫には空の箱がいくつか残っている。\newline 用例: The side entrance will remain closed during the repairs.\newline 訳: 修理中は、横の入口は閉鎖されたままになる。 \\
\addlinespace[\vocabrowspace]
\en{point out}\newline{\ipa \ipastrut /pɔɪnt aʊt/} & \pos{動詞句（他動詞）}\newline 指摘する・示す & 見落とされそうな点や重要な点に、人の注意を向ける。\newline 用例: The editor pointed out a missing caption beneath the photograph.\newline 訳: 編集者は、写真の下に説明文がないと指摘した。 \\
\addlinespace[\vocabrowspace]
\en{maintain}\newline{\ipa \ipastrut /meɪnt\ipastress{vkey}{ˈ}{eɪ}{0.339250}{0.192500}n/} & \pos{他動詞}\newline 維持する・保つ & 一定の状態や水準が続くように、保ち続ける。\newline 用例: Regular cleaning helps maintain the quality of the equipment.\newline 訳: 定期的な清掃は、機材の品質を保つのに役立つ。 \\
\addlinespace[\vocabrowspace]
\en{accept}\newline{\ipa \ipastrut /əks\ipastress{vkey}{ˈ}{ɛ}{0.115750}{0.193000}pt/} & \pos{他動詞}\newline 受け入れる・受け付ける & 提出物や提案などを、受け取ってよいものとして扱う。\newline 用例: The committee accepted our proposal for a weekend workshop.\newline 訳: 委員会は、週末に研修を開くという私たちの提案を受け入れた。 \\
\addlinespace[\vocabrowspace]
\en{discover}\newline{\ipa \ipastrut /dɪsk\ipastress{vkey}{ˈ}{ʌ}{0.142750}{0.179500}vər/} & \pos{他動詞}\newline 発見する・気づく & 調べたり経験したりする中で、それまで知らなかったことを見つける。\newline 用例: We discovered a quiet workspace behind the main library.\newline 訳: 私たちは中央図書館の裏に静かな作業場所を見つけた。 \\
\addlinespace[\vocabrowspace]
\en{distinguish}\newline{\ipa \ipastrut /dɪst\ipastress{vkey}{ˈ}{ɪ}{0.032500}{0.179500}ŋɡwɪʃ/} & \pos{他動詞}\newline 区別する・見分ける & 異なるものの特徴を捉え、それぞれを見分ける。\newline 用例: Clear labels help visitors distinguish originals from printed copies.\newline 訳: 分かりやすいラベルによって、来場者は原本と印刷された複製を見分けられる。 \\
\addlinespace[\vocabrowspace]
\en{measure}\newline{\ipa \ipastrut /m\ipastress{vkey}{ˈ}{ɛ}{0.115750}{0.193000}ʒər/} & \pos{他動詞}\newline 測る・測定する & 対象の大きさや程度を、数値などで確かめる。\newline 用例: The technician measured the noise near each ventilation unit.\newline 訳: 技術者は各換気装置の近くで騒音を測定した。 \\
\addlinespace[\vocabrowspace]
\en{ensure}\newline{\ipa \ipastrut /ɪnʃ\ipastress{vkey}{ˈ}{ʊ}{0.156750}{0.179500}r/} & \pos{他動詞}\newline 確実にする・確保する & 必要な状態や結果が実現するように手配し、確かなものにする。\newline 用例: A second battery ensures uninterrupted power during the presentation.\newline 訳: 予備のバッテリーが、プレゼンテーション中の途切れない電力供給を確保する。 \\
\addlinespace[\vocabrowspace]
\en{include}\newline{\ipa \ipastrut /ɪnkl\ipastress{vkey}{ˈ}{u}{0.160750}{0.179500}d/} & \pos{他動詞}\newline 含む・含める & より大きなまとまりの一部として、何かを含める。\newline 用例: The welcome pack includes a map and a notebook.\newline 訳: 歓迎用の配布物には、地図とノートが入っている。 \\
\addlinespace[\vocabrowspace]
\en{determine}\newline{\ipa \ipastrut /dɪt\ipastress{vkey}{ˈ}{ɜ}{0.115750}{0.193000}rmɪn/} & \pos{他動詞}\newline 判断する・突き止める・決める & 調査で事実を突き止めたり、選択の内容を決めたりする。\newline 用例: We need to determine the cause of the delay.\newline 訳: 私たちは遅延の原因を突き止める必要がある。\newline 用例: Available floor space will determine the size of the display.\newline 訳: 使える床面積によって、展示の大きさが決まる。 \\
\addlinespace[\vocabrowspace]
\en{identify}\newline{\ipa \ipastrut /aɪd\ipastress{vkey}{ˈ}{ɛ}{0.115750}{0.193000}ntəf\ipastress{vhead}{ˌ}{aɪ}{0.340750}{0.701000}/} & \pos{他動詞}\newline 特定する・見つけ出す & いくつかの候補から、該当するものを特定する。\newline 用例: The librarian identified the missing volume from its description.\newline 訳: 司書は説明から、なくなった巻を特定した。 \\
\addlinespace[\vocabrowspace]
\en{reduce}\newline{\ipa \ipastrut /rɪd\ipastress{vkey}{ˈ}{u}{0.160750}{0.179500}s/} & \pos{他動詞}\newline 減らす・小さくする & 数量や程度を、今より少なくしたり小さくしたりする。\newline 用例: Shared transport can reduce the cost of attending workshops.\newline 訳: 相乗りは、研修に参加する費用を減らせる。 \\
\addlinespace[\vocabrowspace]
\en{articulate}\newline{\ipa \ipastrut /ɑrt\ipastress{vkey}{ˈ}{ɪ}{0.032500}{0.179500}kjəl\ipastress{vhead}{ˌ}{eɪ}{0.339250}{0.701000}t/} & \pos{他動詞}\newline 明確に言葉で表す & 考えや要望を、相手に伝わる明確な言葉にする。\newline 用例: The designer articulated her goals before showing the sketches.\newline 訳: デザイナーはスケッチを見せる前に、目標を明確に述べた。 \\
\addlinespace[\vocabrowspace]
\en{acknowledge}\newline{\ipa \ipastrut /əkn\ipastress{vkey}{ˈ}{ɑ}{0.146500}{0.192500}lɪdʒ/} & \pos{他動詞}\newline 認める・明らかにして示す／受け取った・承知した旨を伝える & 事実を認めることや、連絡を受け取った旨を伝えることを表す。\newline 用例: The manager acknowledged the mistake during the morning meeting.\newline 訳: 管理者は朝の会議で、その誤りを認めた。\newline 用例: Please acknowledge receipt of the documents by this afternoon.\newline 訳: 今日の午後までに、書類を受領した旨をお知らせください。 \\
\addlinespace[\vocabrowspace]
\en{collaborate}\newline{\ipa \ipastrut /kəl\ipastress{vkey}{ˈ}{æ}{0.341750}{0.192500}bər\ipastress{vhead}{ˌ}{eɪ}{0.339250}{0.701000}t/} & \pos{自動詞}\newline 協働する・一緒に取り組む & ほかの人や組織と力を合わせ、一つの仕事に取り組む。\newline 用例: The two museums collaborated on a local history exhibition.\newline 訳: 二つの博物館が協力して、郷土史の展覧会に取り組んだ。 \\
\addlinespace[\vocabrowspace]
\en{iterate}\newline{\ipa \ipastrut /\ipastress{vkey}{ˈ}{ɪ}{0.032500}{0.179500}tər\ipastress{vhead}{ˌ}{eɪ}{0.339250}{0.701000}t/} & \pos{自動詞}\newline 試して見直すことを繰り返す & 試した結果を確かめ、改善を重ねながら繰り返す。\newline 用例: The design team iterated on the prototype throughout the summer.\newline 訳: 設計チームは夏の間ずっと、試作品の改良を繰り返した。 \\
\addlinespace[\vocabrowspace]
\en{refine}\newline{\ipa \ipastrut /rɪf\ipastress{vkey}{ˈ}{aɪ}{0.340750}{0.192500}n/} & \pos{他動詞}\newline 磨く・よりよく調整する & すでにある案や仕組みを、細部までよりよく整える。\newline 用例: Visitor interviews helped us refine the layout of the gallery.\newline 訳: 来場者への聞き取りは、展示室の配置を改善する助けになった。 \\
\addlinespace[\vocabrowspace]
\en{tweak}\newline{\ipa \ipastrut /twik/} & \pos{他動詞}\newline 少し調整する・微修正する & 全体を作り直さず、細かな調整を加える。\newline 用例: She tweaked the slide colors before the final rehearsal.\newline 訳: 彼女は最後のリハーサルの前に、スライドの色を少し調整した。 \\
\addlinespace[\vocabrowspace]
\en{provide}\newline{\ipa \ipastrut /prəv\ipastress{vkey}{ˈ}{aɪ}{0.340750}{0.192500}d/} & \pos{他動詞}\newline 提供する・渡す & 相手が必要とするものを用意し、利用できるようにする。\newline 用例: The venue provides spare microphones for visiting speakers.\newline 訳: 会場は来訪する講演者に予備のマイクを提供する。 \\
\addlinespace[\vocabrowspace]
\en{handle}\newline{\ipa \ipastrut /h\ipastress{vkey}{ˈ}{æ}{0.341750}{0.192500}ndəl/} & \pos{他動詞}\newline 扱う・対処する・処理する & 仕事や問題を引き受け、必要な処理を行う。\newline 用例: Our small office handles all registrations for the annual festival.\newline 訳: 私たちの小さな事務所が、毎年の祭りの申込みをすべて処理する。 \\
\addlinespace[\vocabrowspace]
\en{support}\newline{\ipa \ipastrut /səp\ipastress{vkey}{ˈ}{ɔ}{0.125500}{0.192500}rt/} & \pos{名詞・他動詞}\newline 支える・対応する／支援・サポート & 活動を支えることや、機能への対応、提供される支援を表す。\newline 用例: The new editor supports several common image formats.\newline 訳: 新しいエディターは、一般的な画像形式のいくつかに対応している。\newline 用例: Local businesses support the festival by donating equipment.\newline 訳: 地元企業は機材を寄付して祭りを支援している。\newline 用例: The library offers technical support during evening study sessions.\newline 訳: 図書館は夜間の学習時間に技術的な支援を提供している。 \\
\addlinespace[\vocabrowspace]
\en{select}\newline{\ipa \ipastrut /səl\ipastress{vkey}{ˈ}{ɛ}{0.115750}{0.193000}kt/} & \pos{他動詞}\newline 選ぶ・選択する & 候補の中から、目的や条件に合うものを選ぶ。\newline 用例: Please select a delivery date from the available options.\newline 訳: 選べる候補から配達日を選択してください。 \\
\addlinespace[\vocabrowspace]
\en{return}\newline{\ipa \ipastrut /rɪt\ipastress{vkey}{ˈ}{ɜ}{0.115750}{0.193000}rn/} & \pos{他動詞・自動詞}\newline 返す・戻す／戻る & 物や結果を相手に返すこと、元の場所へ戻ることを表す。\newline 用例: The search function returns a list of matching titles.\newline 訳: 検索機能は、一致するタイトルの一覧を返す。\newline 用例: The technicians will return after the replacement parts arrive.\newline 訳: 交換部品が届いたら、技術者たちは戻ってくる。 \\
\addlinespace[\vocabrowspace]
\en{implement}\newline{\ipa \ipastrut /\ipastress{vkey}{ˈ}{ɪ}{0.032500}{0.179500}mpləm\ipastress{vhead}{ˌ}{ɛ}{0.115750}{0.701500}nt/} & \pos{他動詞}\newline 実装する・実施する & 決めた計画を実施したり、設計した機能を動く形にしたりする。\newline 用例: We implemented a simpler booking process before the busy season.\newline 訳: 繁忙期の前に、私たちはより簡単な予約手順を導入した。 \\
\addlinespace[\vocabrowspace]
\en{apply}\newline{\ipa \ipastrut /əpl\ipastress{vkey}{ˈ}{aɪ}{0.340750}{0.192500}/} & \pos{他動詞}\newline 適用する・使う & 規則や方法を、具体的な対象に当てはめて使う。\newline 用例: Please apply the same naming rule to every folder.\newline 訳: すべてのフォルダーに同じ命名規則を適用してください。 \\
\addlinespace[\vocabrowspace]
\en{allow}\newline{\ipa \ipastrut /əl\ipastress{vkey}{ˈ}{aʊ}{0.465750}{0.192500}/} & \pos{他動詞}\newline 許す・認める & 行動を許可するほか、条件や仕組みが何かを可能にする。\newline 用例: The studio allows visitors to photograph the finished exhibits.\newline 訳: その工房では、来場者が完成した展示物を撮影することを認めている。\newline 用例: The wider entrance allows large deliveries to pass easily.\newline 訳: 入口が広くなったことで、大きな荷物を簡単に運び込める。 \\
\addlinespace[\vocabrowspace]
\en{avoid}\newline{\ipa \ipastrut /əv\ipastress{vkey}{ˈ}{ɔɪ}{0.310000}{0.192500}d/} & \pos{他動詞}\newline 避ける & 望ましくない対象や状況に関わらずに済ませる。\newline 用例: We scheduled the delivery early to avoid the lunchtime traffic.\newline 訳: 昼どきの渋滞を避けるため、私たちは配達を早い時間に設定した。 \\
\addlinespace[\vocabrowspace]
\en{prevent}\newline{\ipa \ipastrut /prɪv\ipastress{vkey}{ˈ}{ɛ}{0.115750}{0.193000}nt/} & \pos{他動詞}\newline 防ぐ・起こらないようにする & 問題となる出来事が起こる前に、その発生を止める。\newline 用例: The protective cover prevents dust from entering the keyboard.\newline 訳: 保護カバーは、ほこりがキーボードに入るのを防ぐ。 \\
\addlinespace[\vocabrowspace]
\en{retain}\newline{\ipa \ipastrut /rɪt\ipastress{vkey}{ˈ}{eɪ}{0.339250}{0.192500}n/} & \pos{他動詞}\newline 保持する・残しておく & 持っている情報や性質を、失わずに残しておく。\newline 用例: The renovated hall retains many of its original wooden panels.\newline 訳: 改装されたホールには、元の木製パネルが数多く残されている。 \\
\addlinespace[\vocabrowspace]
\en{preserve}\newline{\ipa \ipastrut /prɪz\ipastress{vkey}{ˈ}{ɜ}{0.115750}{0.193000}rv/} & \pos{他動詞}\newline 保つ・元の状態を維持する & 大切なものの状態や内容を、損なわずに保つ。\newline 用例: Special packaging helps preserve the photographs during long journeys.\newline 訳: 専用の梱包材は、長距離の輸送中に写真を保護するのに役立つ。 \\
\addlinespace[\vocabrowspace]
\en{restore}\newline{\ipa \ipastrut /rɪst\ipastress{vkey}{ˈ}{ɔ}{0.125500}{0.192500}r/} & \pos{他動詞}\newline 元に戻す・復元する & 失われたり変わったりしたものを、以前の状態に戻す。\newline 用例: We restored the deleted folder from yesterday's backup.\newline 訳: 私たちは昨日のバックアップから、削除されたフォルダーを復元した。 \\
\addlinespace[\vocabrowspace]
\en{replace}\newline{\ipa \ipastrut /rɪpl\ipastress{vkey}{ˈ}{eɪ}{0.339250}{0.192500}s/} & \pos{他動詞}\newline 置き換える & 現在あるものを取り除き、その位置に別のものを入れる。\newline 用例: The caretaker replaced the broken handle with a new one.\newline 訳: 管理人は壊れた取っ手を新しいものに交換した。 \\
\addlinespace[\vocabrowspace]
\en{remove}\newline{\ipa \ipastrut /rɪm\ipastress{vkey}{ˈ}{u}{0.160750}{0.179500}v/} & \pos{他動詞}\newline 取り除く・削除する & 不要なものや妨げになるものを、その場所から取り除く。\newline 用例: Please remove the old notices from the entrance wall.\newline 訳: 入口の壁から古い掲示物を取り除いてください。 \\
\addlinespace[\vocabrowspace]
\en{verify}\newline{\ipa \ipastrut /v\ipastress{vkey}{ˈ}{ɛ}{0.115750}{0.193000}rəfaɪ/} & \pos{他動詞}\newline 確かめる・検証する & 証拠や基準と照らし合わせて、正しいことを確かめる。\newline 用例: The assistant verified every measurement against the original drawing.\newline 訳: 担当者はすべての寸法を元の図面と照合して確かめた。 \\
\addlinespace[\vocabrowspace]
\en{confirm}\newline{\ipa \ipastrut /kənf\ipastress{vkey}{ˈ}{ɜ}{0.115750}{0.193000}rm/} & \pos{他動詞}\newline 確認する・確かめる & 予定や情報が確かなものかどうかを確認する。\newline 用例: Please confirm the room number before sending the invitations.\newline 訳: 招待状を送る前に、部屋番号を確認してください。 \\
\addlinespace[\vocabrowspace]
\en{detect}\newline{\ipa \ipastrut /dɪt\ipastress{vkey}{ˈ}{ɛ}{0.115750}{0.193000}kt/} & \pos{他動詞}\newline 見つける・検出する & 見えにくい変化や問題の存在を、調べて見つける。\newline 用例: The sensor detected a small temperature change near the window.\newline 訳: センサーは窓の近くのわずかな温度変化を検出した。 \\
\addlinespace[\vocabrowspace]
\en{inspect}\newline{\ipa \ipastrut /ɪnsp\ipastress{vkey}{ˈ}{ɛ}{0.115750}{0.193000}kt/} & \pos{他動詞}\newline 詳しく調べる・点検する & 状態や不具合を把握するために、対象を注意深く調べる。\newline 用例: The technician inspected the cables before connecting the projector.\newline 訳: 技術者はプロジェクターを接続する前にケーブルを点検した。 \\
\addlinespace[\vocabrowspace]
\en{resolve}\newline{\ipa \ipastrut /rɪz\ipastress{vkey}{ˈ}{ɑ}{0.146500}{0.192500}lv/} & \pos{他動詞}\newline 解決する・決着をつける & 問題を解決することや、名前が指す具体的な対象を特定することを表す。\newline 用例: A brief phone call resolved the misunderstanding about the venue.\newline 訳: 短い電話で、会場についての行き違いが解消した。\newline 用例: The browser resolves the hostname before opening the connection.\newline 訳: ブラウザーは接続を開く前に、ホスト名の参照先を特定する。 \\
\addlinespace[\vocabrowspace]
\en{ignore}\newline{\ipa \ipastrut /ɪɡn\ipastress{vkey}{ˈ}{ɔ}{0.125500}{0.192500}r/} & \pos{他動詞}\newline 無視する・考慮しない & 対象に注意を向けず、判断や行動に取り入れない。\newline 用例: The printer ignores blank pages when this option is enabled.\newline 訳: この設定を有効にすると、プリンターは空白ページを無視する。 \\
\addlinespace[\vocabrowspace]
\en{discard}\newline{\ipa \ipastrut /dɪsk\ipastress{vkey}{ˈ}{ɑ}{0.146500}{0.192500}rd/} & \pos{他動詞}\newline 捨てる・採用しない & 不要と判断したものを、残さず捨てる。\newline 用例: We discarded the damaged samples after documenting their condition.\newline 訳: 私たちは状態を記録した後、傷んだ試料を廃棄した。 \\
\addlinespace[\vocabrowspace]
\en{exceed}\newline{\ipa \ipastrut /ɪks\ipastress{vkey}{ˈ}{i}{-0.013750}{0.392500}d/} & \pos{他動詞}\newline 超える・上回る & 決められた量や基準よりも、大きくなることを表す。\newline 用例: The total weight must not exceed the shelf's stated limit.\newline 訳: 総重量は、その棚に示された上限を超えてはならない。 \\
\addlinespace[\vocabrowspace]
\en{assess}\newline{\ipa \ipastrut /əs\ipastress{vkey}{ˈ}{ɛ}{0.115750}{0.193000}s/} & \pos{他動詞}\newline 評価する・見極める & 状態や価値を調べ、どの程度かを見極める。\newline 用例: The team assessed the building's condition before planning repairs.\newline 訳: チームは修理を計画する前に、建物の状態を評価した。 \\
\addlinespace[\vocabrowspace]
\en{investigate}\newline{\ipa \ipastrut /ɪnv\ipastress{vkey}{ˈ}{ɛ}{0.115750}{0.193000}stəɡeɪt/} & \pos{他動詞}\newline 調査する・詳しく調べる & 原因や事実を明らかにするため、詳しく調べる。\newline 用例: Engineers investigated the unusual vibration in the ventilation system.\newline 訳: 技術者たちは換気設備の異常な振動を調査した。 \\
\addlinespace[\vocabrowspace]
\en{affect}\newline{\ipa \ipastrut /əf\ipastress{vkey}{ˈ}{ɛ}{0.115750}{0.193000}kt/} & \pos{他動詞}\newline 影響を与える & 物事の状態や結果に変化を与えることを表す。\newline 用例: Poor lighting can affect the accuracy of color measurements.\newline 訳: 照明が悪いと、色の測定精度に影響することがある。 \\
\addlinespace[\vocabrowspace]
\en{observe}\newline{\ipa \ipastrut /əbz\ipastress{vkey}{ˈ}{ɜ}{0.115750}{0.193000}rv/} & \pos{他動詞}\newline 観察する・確認する & 実際の様子や変化に注意を向け、見て確かめる。\newline 用例: The researchers observed changes in leaf color over several weeks.\newline 訳: 研究者たちは数週間にわたって葉の色の変化を観察した。 \\
\addlinespace[\vocabrowspace]
\en{prefer}\newline{\ipa \ipastrut /prɪf\ipastress{vkey}{ˈ}{ɜ}{0.115750}{0.193000}r/} & \pos{他動詞}\newline より好む・優先して選ぶ & 複数の候補の中で、よりよいと感じるものを選ぶ。\newline 用例: Most participants prefer short practical exercises to long lectures.\newline 訳: 参加者の多くは、長い講義より短い実習を好む。 \\
\addlinespace[\vocabrowspace]
\en{restrict}\newline{\ipa \ipastrut /rɪstr\ipastress{vkey}{ˈ}{ɪ}{0.032500}{0.179500}kt/} & \pos{他動詞}\newline 制限する・範囲を絞る & 利用や行動が及ぶ範囲を、一定の内側に抑える。\newline 用例: The library restricts access to fragile books during renovations.\newline 訳: 図書館は改装中、傷みやすい本の利用を制限している。 \\
\addlinespace[\vocabrowspace]
\en{suppress}\newline{\ipa \ipastrut /səpr\ipastress{vkey}{ˈ}{ɛ}{0.115750}{0.193000}s/} & \pos{他動詞}\newline 抑える・表に出さない & 表示や活動を抑えて、表に出ないようにする。\newline 用例: The quiet setting suppresses routine notifications during presentations.\newline 訳: 静音設定は、プレゼンテーション中の通常の通知を表示しないようにする。 \\
\addlinespace[\vocabrowspace]
\en{extract}\newline{\ipa \ipastrut /ɪkstr\ipastress{vkey}{ˈ}{æ}{0.341750}{0.192500}kt/} & \pos{他動詞}\newline 抽出する・取り出す & まとまったものの中から、必要な部分を取り出す。\newline 用例: The script extracts chapter titles from the digital manuscript.\newline 訳: そのスクリプトは、電子原稿から章のタイトルを抽出する。 \\
\addlinespace[\vocabrowspace]
\en{evaluate}\newline{\ipa \ipastrut /ɪv\ipastress{vkey}{ˈ}{æ}{0.341750}{0.192500}lju\ipastress{vhead}{ˌ}{eɪ}{0.339250}{0.701000}t/} & \pos{他動詞}\newline 評価する & 目的や基準に照らし、よさや有効性を判断する。\newline 用例: We evaluated several microphones in the actual lecture room.\newline 訳: 私たちは実際の講義室で、いくつかのマイクを評価した。 \\
\addlinespace[\vocabrowspace]
\en{classify}\newline{\ipa \ipastrut /kl\ipastress{vkey}{ˈ}{æ}{0.341750}{0.192500}səf\ipastress{vhead}{ˌ}{aɪ}{0.340750}{0.701000}/} & \pos{他動詞}\newline 分類する & 共通する性質に応じて、対象を種類ごとに分ける。\newline 用例: The archive classifies photographs by location and approximate date.\newline 訳: 資料館は場所とおおよその年代で写真を分類している。 \\
\addlinespace[\vocabrowspace]
\en{flag}\newline{\ipa \ipastrut /flæɡ/} & \pos{他動詞}\newline 注意すべきものとして示す・指摘する & 見直しや注意が必要な箇所に、目印を付けて知らせる。\newline 用例: The editor flagged several unclear passages in the draft.\newline 訳: 編集者は草稿の中の分かりにくい箇所をいくつか指摘した。 \\
\addlinespace[\vocabrowspace]
\en{complain}\newline{\ipa \ipastrut /kəmpl\ipastress{vkey}{ˈ}{eɪ}{0.339250}{0.192500}n/} & \pos{自動詞}\newline 不満を述べる・苦情を言う & 好ましくない状況について、不満や苦情を伝える。\newline 用例: Several residents complained about noise from the evening construction.\newline 訳: 何人かの住民が、夜の工事の騒音について苦情を述べた。 \\
\addlinespace[\vocabrowspace]
\en{present}\newline{\ipa \ipastrut v. /prɪz\ipastress{vkey}{ˈ}{ɛ}{0.115750}{0.193000}nt/  adj. /pr\ipastress{vkey}{ˈ}{ɛ}{0.115750}{0.193000}zənt/} & \pos{他動詞・形容詞}\newline （回答・情報など）を提示する・示す／存在する・そこにある & 情報を相手に示すことや、何かがその場にある状態を表す。\newline 用例: The researcher presented her findings at a regional conference.\newline 訳: 研究者は地域の学会で調査結果を発表した。\newline 用例: Small air bubbles are present in several glass samples.\newline 訳: いくつかのガラス試料には、小さな気泡が存在する。 \\
\addlinespace[\vocabrowspace]
\en{trial}\newline{\ipa \ipastrut /tr\ipastress{vkey}{ˈ}{aɪ}{0.340750}{0.192500}əl/} & \pos{他動詞・名詞}\newline 試す・試験的に使う／試行・試験 & 本格導入の前に試すことや、そのための試験を表す。\newline 用例: The school will trial the new timetable next month.\newline 訳: 学校は来月、新しい時間割を試験的に導入する。\newline 用例: The initial trial revealed several problems with the seating layout.\newline 訳: 最初の試行で、座席配置の問題がいくつか分かった。 \\
\addlinespace[\vocabrowspace]
\en{inject}\newline{\ipa \ipastrut /ɪndʒ\ipastress{vkey}{ˈ}{ɛ}{0.115750}{0.193000}kt/} & \pos{他動詞}\newline 注入する・（情報など）を入れる・追加する & 既存のものの中へ、情報や要素を外から加える。\newline 用例: The test harness injects artificial delays into network requests.\newline 訳: テスト用の仕組みは、ネットワークリクエストに人為的な遅延を加える。 \\
\addlinespace[\vocabrowspace]
\en{capture}\newline{\ipa \ipastrut /k\ipastress{vkey}{ˈ}{æ}{0.341750}{0.192500}ptʃər/} & \pos{他動詞}\newline （情報を）記録する・保存する & 一時的な情報や出来事を、後から使える記録として残す。\newline 用例: The recorder captures every spoken comment during the interview.\newline 訳: 録音機は、面接中に話された発言をすべて記録する。 \\
\addlinespace[\vocabrowspace]
\en{surface}\newline{\ipa \ipastrut /s\ipastress{vkey}{ˈ}{ɜ}{0.115750}{0.193000}rfəs/} & \pos{他動詞}\newline （情報・論点を）明らかにする・表に出す & 見えにくかった情報や問題を、目に見える形で取り上げる。\newline 用例: The staff survey surfaced concerns about the new office layout.\newline 訳: 職員調査によって、新しい事務所の配置に対する懸念が明らかになった。 \\
\addlinespace[\vocabrowspace]
\en{reveal}\newline{\ipa \ipastrut /rɪv\ipastress{vkey}{ˈ}{i}{-0.013750}{0.392500}l/} & \pos{他動詞}\newline （未知の情報・事実などを）明らかにする・示す & 隠れていた事実や分かっていなかったことを明らかにする。\newline 用例: Closer inspection revealed a crack beneath the painted surface.\newline 訳: 詳しく調べると、塗装面の下にひびがあると分かった。 \\
\addlinespace[\vocabrowspace]
\en{offer}\newline{\ipa \ipastrut /\ipastress{vkey}{ˈ}{ɑ}{0.146500}{0.192500}fər/} & \pos{他動詞}\newline 提供を申し出る・提案する & 相手に何かを提供できると伝え、選べるようにする。\newline 用例: The host offered each speaker a choice of interview times.\newline 訳: 主催者は各講演者に、面接時間を選べるよう提案した。 \\
\addlinespace[\vocabrowspace]
\en{describe}\newline{\ipa \ipastrut /dɪskr\ipastress{vkey}{ˈ}{aɪ}{0.340750}{0.192500}b/} & \pos{他動詞}\newline 内容や様子を説明する・記述する & 対象の特徴や様子を、言葉で相手に伝える。\newline 用例: The brochure describes the walking route around the old harbor.\newline 訳: パンフレットは、古い港の周りの散策路を説明している。 \\
\addlinespace[\vocabrowspace]
\en{misread}\newline{\ipa \ipastrut /m\ipastress{vhead}{ˌ}{ɪ}{0.032500}{0.688000}sr\ipastress{vkey}{ˈ}{i}{-0.013750}{0.392500}d/} & \pos{他動詞}\newline 読み違える・誤って解釈する & 文字や内容を、正しい意味とは違う形で受け取る。\newline 用例: I misread the timetable and arrived an hour early.\newline 訳: 私は時刻表を読み違えて、一時間早く到着した。 \\
\addlinespace[\vocabrowspace]
\en{break}\newline{\ipa \ipastrut /breɪk/} & \pos{自動詞・他動詞}\newline 正常に動かなくなる・処理が破綻する／壊す & 動作が成り立たなくなることや、何かを壊してしまうことを表す。\newline 用例: The animation may break when the required font is unavailable.\newline 訳: 必要なフォントがないと、アニメーションが正常に動かなくなることがある。\newline 用例: The update broke the shortcut for opening recent files.\newline 訳: その更新によって、最近のファイルを開くショートカットが動かなくなった。 \\
\addlinespace[\vocabrowspace]
\en{exhaust}\newline{\ipa \ipastrut /ɪɡz\ipastress{vkey}{ˈ}{ɑ}{0.146500}{0.192500}st/} & \pos{他動詞}\newline 使い果たす・使い尽くす & 資源を使い切ることや、人をひどく疲れさせることを表す。\newline 用例: The unexpected repairs exhausted our maintenance budget for the year.\newline 訳: 予想外の修理で、その年の保守予算を使い切ってしまった。\newline 用例: The long journey exhausted everyone before the workshop began.\newline 訳: 長旅のせいで、研修が始まる前から全員が疲れ果てていた。 \\
\addlinespace[\vocabrowspace]
\en{list}\newline{\ipa \ipastrut /lɪst/} & \pos{他動詞}\newline 列挙する・項目として記載する & 情報を項目ごとに並べ、一覧として示す。\newline 用例: The notice lists all available rooms and their capacities.\newline 訳: 掲示には、利用可能なすべての部屋と定員が記載されている。 \\
\addlinespace[\vocabrowspace]
\en{lack}\newline{\ipa \ipastrut /læk/} & \pos{他動詞}\newline …を欠いている・…がない & 必要な性質や要素を、持っていないことを表す。\newline 用例: The older building lacks a convenient entrance for deliveries.\newline 訳: 古い方の建物には、搬入に便利な入口がない。 \\
\addlinespace[\vocabrowspace]
\en{vary}\newline{\ipa \ipastrut /v\ipastress{vkey}{ˈ}{ɛ}{0.115750}{0.193000}ri/} & \pos{自動詞・他動詞}\newline 異なる・ばらつく・変化する／変える & 違いがあること、変化すること、意図的に変えることを表す。\newline 用例: Opening hours vary throughout the year with visitor demand.\newline 訳: 開館時間は来場者の需要に合わせ、一年を通して変わる。\newline 用例: We vary the exercises to keep the workshops engaging.\newline 訳: 研修への関心を保つため、私たちは実習の内容を変えている。 \\
\addlinespace[\vocabrowspace]
\en{remediate}\newline{\ipa \ipastrut /rɪm\ipastress{vkey}{ˈ}{i}{-0.013750}{0.392500}dieɪt/} & \pos{他動詞}\newline 問題を修正する・是正する & 確認された不具合や悪い状態に、是正の処置を施す。\newline 用例: The team remediated accessibility issues before publishing the website.\newline 訳: チームはウェブサイトの公開前に、アクセシビリティの問題を修正した。 \\
\addlinespace[\vocabrowspace]
\en{reach}\newline{\ipa \ipastrut /ritʃ/} & \pos{他動詞・自動詞}\newline …に届く・…に到達する／手を伸ばす & 対象に届くことや、何かを取ろうと手を伸ばすことを表す。\newline 用例: The wireless signal barely reaches the farthest meeting room.\newline 訳: 無線の電波は、一番遠い会議室にかろうじて届く。\newline 用例: She reached for the notebook on the upper shelf.\newline 訳: 彼女は上の棚にあるノートに手を伸ばした。 \\
\addlinespace[\vocabrowspace]
\en{touch}\newline{\ipa \ipastrut /tʌtʃ/} & \pos{他動詞}\newline …に関係する・…を扱う & 変更や説明などが、ある対象にも及ぶことを表す。\newline 用例: The renovation touches every public area of the building.\newline 訳: その改装は、建物のすべての共用部分に及ぶ。 \\
\addlinespace[\vocabrowspace]
\en{anticipate}\newline{\ipa \ipastrut /ænt\ipastress{vkey}{ˈ}{ɪ}{0.032500}{0.179500}səp\ipastress{vhead}{ˌ}{eɪ}{0.339250}{0.701000}t/} & \pos{他動詞・自動詞}\newline …をあらかじめ予想する・…を見越す／話を先取りする & 先の出来事を見越すことや、後で述べる内容を先取りすることを表す。\newline 用例: We anticipated heavy traffic and left before sunrise.\newline 訳: 私たちは激しい渋滞を予想し、日の出前に出発した。\newline 用例: But I anticipate; the next chapter explains that discovery.\newline 訳: しかし、話を先取りしてしまった。その発見については次の章で説明する。 \\
\addlinespace[\vocabrowspace]
\en{examine}\newline{\ipa \ipastrut /ɪɡz\ipastress{vkey}{ˈ}{æ}{0.341750}{0.192500}mɪn/} & \pos{他動詞}\newline …を詳しく調べる・…を注意深く確認する & 細部や根拠に注意しながら、対象を詳しく調べる。\newline 用例: The curator examined the signature under a bright lamp.\newline 訳: 学芸員は明るいランプの下で署名を詳しく調べた。 \\
\addlinespace[\vocabrowspace]
\en{pinpoint}\newline{\ipa \ipastrut /p\ipastress{vkey}{ˈ}{ɪ}{0.032500}{0.179500}np\ipastress{vhead}{ˌ}{ɔɪ}{0.310000}{0.701000}nt/} & \pos{他動詞}\newline …を正確に特定する・…を的確に示す & 広い候補の中から、位置や原因を正確に突き止める。\newline 用例: The map helped us pinpoint the damaged section of pipe.\newline 訳: 地図によって、私たちは配管の損傷箇所を正確に特定できた。 \\
\addlinespace[\vocabrowspace]
\en{compact}\newline{\ipa \ipastrut /kəmp\ipastress{vkey}{ˈ}{æ}{0.341750}{0.192500}kt/} & \pos{他動詞}\newline …を圧縮する・…を短くまとめる & 情報や物を、より少ない空間に収まるよう圧縮する。\newline 用例: The utility compacts the database during the scheduled maintenance window.\newline 訳: そのユーティリティーは、予定された保守時間中にデータベースを圧縮する。 \\
\addlinespace[\vocabrowspace]
\en{save}\newline{\ipa \ipastrut /seɪv/} & \pos{他動詞}\newline （時間・費用・手間など）を節約する・省く & 必要になる時間や費用、手間を少なくする。\newline 用例: Booking the equipment together saved us several delivery fees.\newline 訳: 機材をまとめて予約したことで、配送料を数回分節約できた。 \\
\addlinespace[\vocabrowspace]
\en{record}\newline{\ipa \ipastrut /rɪk\ipastress{vkey}{ˈ}{ɔ}{0.125500}{0.192500}rd/} & \pos{他動詞}\newline …を記録する・…を記録として残す & 起きたことや得られた情報を、後に残る形で書き留める。\newline 用例: The assistant recorded each visitor's preferred appointment time.\newline 訳: 担当者は、各来訪者が希望する予約時間を記録した。 \\
\addlinespace[\vocabrowspace]
\en{append}\newline{\ipa \ipastrut /əp\ipastress{vkey}{ˈ}{ɛ}{0.115750}{0.193000}nd/} & \pos{他動詞}\newline …を末尾に追加する・…を付け足す & すでにある文章や一覧の末尾に、内容を追加する。\newline 用例: Please append the equipment list to the meeting agenda.\newline 訳: 会議の議題の末尾に機材一覧を追加してください。 \\
\addlinespace[\vocabrowspace]
\en{involve}\newline{\ipa \ipastrut /ɪnv\ipastress{vkey}{ˈ}{ɑ}{0.146500}{0.192500}lv/} & \pos{他動詞}\newline …を伴う・…を含む・…に関係する & 活動の一部として何かが含まれたり、必要になったりする。\newline 用例: The restoration will involve extensive research into traditional materials.\newline 訳: 修復には、伝統的な材料についての広範な調査が伴う。 \\
\addlinespace[\vocabrowspace]
\en{expose}\newline{\ipa \ipastrut /ɪksp\ipastress{vkey}{ˈ}{oʊ}{0.462750}{0.192500}z/} & \pos{他動詞}\newline …を公開する・…を利用できるようにする & 内部にある機能や情報を、外から利用できる形にする。\newline 用例: The dashboard exposes several controls for adjusting display brightness.\newline 訳: その操作画面には、表示の明るさを調整する操作項目がいくつか用意されている。 \\
\addlinespace[\vocabrowspace]
\en{enumerate}\newline{\ipa \ipastrut /ɪn\ipastress{vkey}{ˈ}{u}{0.160750}{0.179500}mər\ipastress{vhead}{ˌ}{eɪ}{0.339250}{0.701000}t/} & \pos{他動詞}\newline …を一つずつ挙げる・…を列挙する & 含まれるものを、一つずつ漏れなく挙げて示す。\newline 用例: The checklist enumerates every item needed for the field trip.\newline 訳: チェックリストには、見学旅行に必要な品目が一つずつ記載されている。 \\
\addlinespace[\vocabrowspace]
\en{fetch}\newline{\ipa \ipastrut /fɛtʃ/} & \pos{他動詞}\newline …を取りに行く・…を取得する & 必要な情報や物がある場所から、それを取ってくる。\newline 用例: The application fetches fresh weather readings from nearby stations.\newline 訳: そのアプリケーションは、近くの観測所から新しい気象データを取得する。 \\
\addlinespace[\vocabrowspace]
\en{revoke}\newline{\ipa \ipastrut /rɪv\ipastress{vkey}{ˈ}{oʊ}{0.462750}{0.192500}k/} & \pos{他動詞}\newline …の効力を取り消す・…を無効にする & 以前与えた許可や権利を取り消し、効力を失わせる。\newline 用例: The administrator revoked access after the temporary project ended.\newline 訳: 管理者は、期間限定の企画が終わった後にアクセス権を取り消した。 \\
\addlinespace[\vocabrowspace]
\en{attach}\newline{\ipa \ipastrut /ət\ipastress{vkey}{ˈ}{æ}{0.341750}{0.192500}tʃ/} & \pos{他動詞}\newline …を取り付ける・…を添付する & 書類に別の資料を添えたり、物を別の物に取り付けたりする。\newline 用例: Please attach the floor plan to your booking request.\newline 訳: 予約の申込みに見取り図を添付してください。\newline 用例: We attached a small lamp to the drafting table.\newline 訳: 私たちは製図台に小さなランプを取り付けた。 \\
\addlinespace[\vocabrowspace]
\en{misstate}\newline{\ipa \ipastrut /m\ipastress{vhead}{ˌ}{ɪ}{0.032500}{0.688000}sst\ipastress{vkey}{ˈ}{eɪ}{0.339250}{0.192500}t/} & \pos{他動詞}\newline …を誤って述べる・…を不正確に伝える & 数字や事実を、実際の内容とは違う形で伝える。\newline 用例: The poster misstated the date of the final performance.\newline 訳: ポスターには、最終公演の日付が誤って記載されていた。 \\
\addlinespace[\vocabrowspace]
\en{dispatch}\newline{\ipa \ipastrut /dɪsp\ipastress{vkey}{ˈ}{æ}{0.341750}{0.192500}tʃ/} & \pos{他動詞}\newline …を送り出す・…を送信する & 目的地へ人や物を送り出したり、処理に向けて情報を送ったりする。\newline 用例: The warehouse dispatched the replacement chairs early this morning.\newline 訳: 倉庫は今朝早く、交換用の椅子を発送した。\newline 用例: The scheduler dispatches a notification when the reservation expires.\newline 訳: 予約の期限が切れると、スケジューラーが通知を送信する。 \\
\addlinespace[\vocabrowspace]
\en{dedicate}\newline{\ipa \ipastrut dedicate /d\ipastress{vkey}{ˈ}{ɛ}{0.115750}{0.193000}dɪk\ipastress{vhead}{ˌ}{eɪ}{0.339250}{0.701000}t/  dedicated /d\ipastress{vkey}{ˈ}{ɛ}{0.115750}{0.193000}dɪk\ipastress{vhead}{ˌ}{eɪ}{0.339250}{0.701000}tɪd/} & \pos{他動詞・形容詞（dedicated）}\newline …を特定の目的に充てる／dedicated：専用の & 資源や時間を特定の目的に充てることや、専用の状態を表す。\newline 用例: We dedicated one afternoon to organizing the shared workspace.\newline 訳: 私たちは、ある日の午後を共用作業場の整理に充てた。\newline 用例: The new studio has a dedicated room for audio recording.\newline 訳: 新しいスタジオには、録音専用の部屋がある。 \\
\addlinespace[\vocabrowspace]
\en{ship}\newline{\ipa \ipastrut /ʃɪp/} & \pos{他動詞}\newline （製品・機能など）を提供する・同梱して提供する & 完成した製品や機能を、利用できる状態で提供する。\newline 用例: The company ships the editor with several accessible color themes.\newline 訳: その会社は、使いやすさに配慮した配色テーマをいくつか同梱してエディターを提供している。 \\
\addlinespace[\vocabrowspace]
\en{hand}\newline{\ipa \ipastrut /hænd/} & \pos{他動詞}\newline （情報・調査結果など）を渡す・引き渡す & 物や情報を、次に扱う相手へ渡すことを表す。\newline 用例: Please hand your interview notes to the project librarian.\newline 訳: 面接の記録を、プロジェクトの資料担当者に渡してください。 \\
\addlinespace[\vocabrowspace]
\en{form}\newline{\ipa \ipastrut /fɔrm/} & \pos{他動詞}\newline （考え・見解など）を形成する・持つようになる & 経験や情報から、考えや見解を作り上げる。\newline 用例: Visitors form their first impression before entering the gallery.\newline 訳: 来場者は展示室に入る前に、第一印象を持つ。 \\
\addlinespace[\vocabrowspace]
\en{anchor}\newline{\ipa \ipastrut /\ipastress{vkey}{ˈ}{æ}{0.341750}{0.192500}ŋkər/} & \pos{自動詞}\newline （判断が特定の情報に）固定される・引きずられる & 最初に触れた情報などに、判断が強く固定される。\newline 用例: Buyers sometimes anchor on the first price they hear.\newline 訳: 買い手は、ときに最初に聞いた価格に判断を引きずられる。 \\
\addlinespace[\vocabrowspace]
\en{honor}\newline{\ipa \ipastrut /\ipastress{vkey}{ˈ}{ɑ}{0.146500}{0.192500}nər/} & \pos{他動詞}\newline （要求・約束など）を尊重して従う・履行する & 約束や要望を重く受け止め、それに従って行動する。\newline 用例: The venue honored our agreement to provide additional seating.\newline 訳: 会場は、追加の座席を用意するという私たちとの約束を守った。 \\
\addlinespace[\vocabrowspace]
\en{pull}\newline{\ipa \ipastrut /pʊl/} & \pos{他動詞}\newline （変更・データを）取得して取り込む & 外部にある情報や変更を取得し、手元へ取り込む。\newline 用例: The dashboard pulls attendance data from the registration database.\newline 訳: その画面は、受付データベースから出席情報を取り込む。 \\
\addlinespace[\vocabrowspace]
\en{belong}\newline{\ipa \ipastrut /bɪl\ipastress{vkey}{ˈ}{ɔ}{0.125500}{0.192500}ŋ/} & \pos{自動詞}\newline 本来の場所にある・適切な場所に置かれる & 何かがあるべき場所や、適切に収まる先を示す。\newline 用例: These reference books belong on the shelf beside the window.\newline 訳: これらの参考書の置き場所は、窓の横の棚だ。 \\
\addlinespace[\vocabrowspace]
\en{style}\newline{\ipa \ipastrut /staɪl/} & \pos{他動詞}\newline …の見た目・書式を整える & 見た目や書式を調整して、望む形に整える。\newline 用例: We styled the headings to match the exhibition posters.\newline 訳: 私たちは展示会のポスターに合わせて、見出しの書式を整えた。 \\
\addlinespace[\vocabrowspace]
\en{contaminate}\newline{\ipa \ipastrut /kənt\ipastress{vkey}{ˈ}{æ}{0.341750}{0.192500}mən\ipastress{vhead}{ˌ}{eɪ}{0.339250}{0.701000}t/} & \pos{他動詞}\newline …を汚染する・…に悪影響を及ぼす & 望ましくないものが混ざり、品質や判断に悪影響を与える。\newline 用例: Leading questions can contaminate the results of a visitor survey.\newline 訳: 誘導的な質問は、来場者調査の結果に悪影響を与えることがある。 \\
\addlinespace[\vocabrowspace]
\en{cover}\newline{\ipa \ipastrut /k\ipastress{vkey}{ˈ}{ʌ}{0.142750}{0.179500}vər/} & \pos{他動詞}\newline （内容・範囲として）扱う・含める & 説明や仕事の対象として、ある範囲を含めて扱う。\newline 用例: The introductory class covers basic lighting and camera placement.\newline 訳: 入門講座では、基本的な照明とカメラの配置を扱う。 \\
\addlinespace[\vocabrowspace]
\en{consult}\newline{\ipa \ipastrut /kəns\ipastress{vkey}{ˈ}{ʌ}{0.142750}{0.179500}lt/} & \pos{他動詞}\newline （資料など）を参照する・調べる & 判断に必要な情報を得るため、資料を参照する。\newline 用例: Please consult the installation manual before changing these settings.\newline 訳: この設定を変える前に、設置マニュアルを参照してください。 \\
\addlinespace[\vocabrowspace]
\en{confuse}\newline{\ipa \ipastrut /kənfj\ipastress{vkey}{ˈ}{u}{0.160750}{0.179500}z/} & \pos{他動詞}\newline …を混乱させる・…を困惑させる & 説明や状況によって、相手を分かりにくい状態にする。\newline 用例: The similar room names confused several first-time visitors.\newline 訳: 部屋の名前が似ていたため、初めて来た人が何人か戸惑った。 \\
\addlinespace[\vocabrowspace]
\en{concentrate}\newline{\ipa \ipastrut /k\ipastress{vkey}{ˈ}{ɑ}{0.146500}{0.192500}nsəntr\ipastress{vhead}{ˌ}{eɪ}{0.339250}{0.701000}t/} & \pos{自動詞}\newline 集中する・一か所に集まる & 活動や問題などが、特定の場所や対象に集まる。\newline 用例: Most visitor traffic concentrates near the main entrance.\newline 訳: 来場者の往来の大部分は、正面入口の近くに集中する。 \\
\addlinespace[\vocabrowspace]
\en{pair / pairing}\newline{\ipa \ipastrut /pɛr/ /p\ipastress{vkey}{ˈ}{ɛ}{0.115750}{0.193000}rɪŋ/} & \pos{他動詞・名詞}\newline 組み合わせる／組合せ & 二つのものを対応させることや、その組み合わせを表す。\newline 用例: The organizer paired each beginner with an experienced volunteer.\newline 訳: 主催者は初心者一人ひとりを、経験豊かなボランティアと組み合わせた。\newline 用例: This pairing of photographs and letters creates a vivid portrait.\newline 訳: この写真と手紙の組み合わせが、生き生きとした人物像を描き出している。 \\
\addlinespace[\vocabrowspace]
\en{reflect}\newline{\ipa \ipastrut /rɪfl\ipastress{vkey}{ˈ}{ɛ}{0.115750}{0.193000}kt/} & \pos{他動詞}\newline 反映する・表す & 内容や結果に、考えや状況の特徴が表れることを示す。\newline 用例: The revised schedule reflects feedback from the evening staff.\newline 訳: 改訂された予定には、夜間担当の職員の意見が反映されている。 \\
\addlinespace[\vocabrowspace]
\en{fabricate}\newline{\ipa \ipastrut /f\ipastress{vkey}{ˈ}{æ}{0.341750}{0.192500}brɪk\ipastress{vhead}{ˌ}{eɪ}{0.339250}{0.701000}t/} & \pos{他動詞}\newline 捏造する・でっち上げる & 実際には存在しない事実や証拠を、あるように作り上げる。\newline 用例: The investigator discovered that someone had fabricated the inspection records.\newline 訳: 調査担当者は、誰かが点検記録を捏造したと突き止めた。 \\
\addlinespace[\vocabrowspace]
\en{invent}\newline{\ipa \ipastrut /ɪnv\ipastress{vkey}{ˈ}{ɛ}{0.115750}{0.193000}nt/} & \pos{他動詞}\newline 作り上げる・でっち上げる & 事実にない内容を、想像で作り出してしまうことを表す。\newline 用例: He invented an excuse for missing the planning meeting.\newline 訳: 彼は、企画会議を欠席した言い訳をでっち上げた。 \\
\addlinespace[\vocabrowspace]
\en{precede}\newline{\ipa \ipastrut /prɪs\ipastress{vkey}{ˈ}{i}{-0.013750}{0.392500}d/} & \pos{他動詞}\newline …に先行する・…より前に行われる & 順番や時間の上で、別のものより先に位置する。\newline 用例: A short introduction will precede the main lecture.\newline 訳: 本講義の前に、短い導入説明が行われる。 \\
\addlinespace[\vocabrowspace]
\en{conflate}\newline{\ipa \ipastrut /kənfl\ipastress{vkey}{ˈ}{eɪ}{0.339250}{0.192500}t/} & \pos{他動詞}\newline （別々のもの）を一緒にして扱う・混同する & 異なる事柄の境界をなくし、一つのものとして扱う。\newline 用例: The discussion conflated visitor satisfaction with the number of tickets sold.\newline 訳: その議論は、来場者の満足度とチケットの販売枚数を混同していた。 \\
\addlinespace[\vocabrowspace]
\en{elicit}\newline{\ipa \ipastrut /ɪl\ipastress{vkey}{ˈ}{ɪ}{0.032500}{0.179500}sɪt/} & \pos{他動詞}\newline （情報・要望・反応など）を引き出す & 問いかけや働きかけによって、情報や反応を引き出す。\newline 用例: Open questions often elicit more detailed answers from participants.\newline 訳: 自由に答えられる質問は、参加者からより詳しい回答を引き出すことが多い。 \\
\addlinespace[\vocabrowspace]
\en{warrant}\newline{\ipa \ipastrut /w\ipastress{vkey}{ˈ}{ɔ}{0.125500}{0.192500}rənt/} & \pos{他動詞}\newline …を行う根拠となる・…を必要とする & 状況や根拠が、ある対応を行うのに十分であることを示す。\newline 用例: The unexpected crack warrants a closer inspection of the wall.\newline 訳: 予想外のひびがあるため、その壁をさらに詳しく点検する必要がある。 \\
\addlinespace[\vocabrowspace]
\en{traverse}\newline{\ipa \ipastrut /trəv\ipastress{vkey}{ˈ}{ɜ}{0.115750}{0.193000}rs/} & \pos{他動詞}\newline …を通過する・…を順にたどる & ある範囲を横切ったり、構造を順番にたどったりする。\newline 用例: The program traverses the folder tree to locate unused images.\newline 訳: そのプログラムは、未使用の画像を見つけるためにフォルダーの階層を順にたどる。 \\
\addlinespace[\vocabrowspace]
\en{forfeit}\newline{\ipa \ipastrut /f\ipastress{vkey}{ˈ}{ɔ}{0.125500}{0.192500}rfɪt/} & \pos{他動詞}\newline （利点・権利など）を失う・手放す & 行動や選択の結果、持っていた権利や利益を失う。\newline 用例: Customers who cancel late forfeit the discount for advance bookings.\newline 訳: 直前にキャンセルした利用者は、早期予約の割引を受ける権利を失う。 \\
\addlinespace[\vocabrowspace]
\en{shoehorn}\newline{\ipa \ipastrut /ʃ\ipastress{vkey}{ˈ}{u}{0.160750}{0.179500}h\ipastress{vhead}{ˌ}{ɔ}{0.125500}{0.701000}rn/} & \pos{他動詞}\newline （無理のある枠に）…を押し込む & 内容や大きさに無理があっても、決まった枠に押し込む。\newline 用例: We tried to shoehorn too many activities into one afternoon.\newline 訳: 私たちは、一日の午後にあまりに多くの活動を詰め込もうとした。 \\
\addlinespace[\vocabrowspace]
\en{supersede}\newline{\ipa \ipastrut /s\ipastress{vhead}{ˌ}{u}{0.160750}{0.688000}pərs\ipastress{vkey}{ˈ}{i}{-0.013750}{0.392500}d/} & \pos{他動詞}\newline （古いもの）に取って代わる・…を置き換える & 新しいものが、以前のものに代わって使われるようになる。\newline 用例: The updated floor plan supersedes every earlier version.\newline 訳: 更新された見取り図は、それ以前のすべての版に取って代わる。 \\
\addlinespace[\vocabrowspace]
\en{exonerate}\newline{\ipa \ipastrut /ɪɡz\ipastress{vkey}{ˈ}{ɑ}{0.146500}{0.192500}nər\ipastress{vhead}{ˌ}{eɪ}{0.339250}{0.701000}t/} & \pos{他動詞}\newline …に責任がないと示す・…の疑いを晴らす & 調査や証拠によって、責任や疑いがないと明らかにする。\newline 用例: The access logs exonerated the technician blamed for the deletion.\newline 訳: アクセス記録によって、削除の責任を問われていた技術者の疑いが晴れた。 \\
\addlinespace[\vocabrowspace]
\en{conjure}\newline{\ipa \ipastrut /k\ipastress{vkey}{ˈ}{ʌ}{0.142750}{0.179500}ndʒər/} & \pos{他動詞}\newline 魔法のように生み出す・出現させる & ないものを、まるで魔法のように生み出すことを表す。\newline 用例: A longer meeting cannot conjure resources that nobody has.\newline 訳: 会議を長くしても、誰も持っていない資源を生み出すことはできない。 \\
\addlinespace[\vocabrowspace]
\en{obliterate}\newline{\ipa \ipastrut /əbl\ipastress{vkey}{ˈ}{ɪ}{0.032500}{0.179500}tər\ipastress{vhead}{ˌ}{eɪ}{0.339250}{0.701000}t/} & \pos{他動詞}\newline 完全に消し去る・跡形もなく消す & 痕跡や内容が残らないほど、完全に消してしまう。\newline 用例: The thick black paint obliterated the lettering on the sign.\newline 訳: 厚く塗られた黒い塗料が、看板の文字を完全に消してしまった。 \\
\addlinespace[\vocabrowspace]
\en{outweigh}\newline{\ipa \ipastrut /\ipastress{vhead}{ˌ}{aʊ}{0.465750}{0.701000}tw\ipastress{vkey}{ˈ}{eɪ}{0.339250}{0.192500}/} & \pos{他動詞}\newline …を重要性・価値で上回る & 比較したときに、価値や影響の大きさで上回る。\newline 用例: The lower maintenance costs outweigh the higher purchase price.\newline 訳: 維持費の低さによる利点は、購入価格の高さを上回る。 \\
\addlinespace[\vocabrowspace]
\en{arbitrate}\newline{\ipa \ipastrut /\ipastress{vkey}{ˈ}{ɑ}{0.146500}{0.192500}rbɪtr\ipastress{vhead}{ˌ}{eɪ}{0.339250}{0.701000}t/} & \pos{他動詞}\newline 調整して決める・裁定する & 対立する要望や主張を調整し、決着を付ける。\newline 用例: An independent panel will arbitrate the dispute over exhibition space.\newline 訳: 独立した委員会が、展示スペースをめぐる争いを裁定する。 \\
\addlinespace[\vocabrowspace]
\en{reinforce}\newline{\ipa \ipastrut /r\ipastress{vhead}{ˌ}{i}{-0.013750}{0.901000}ɪnf\ipastress{vkey}{ˈ}{ɔ}{0.125500}{0.192500}rs/} & \pos{他動詞}\newline 強める・補強する & すでにある傾向や理解、強度をさらに強める。\newline 用例: Practical exercises reinforce the ideas introduced during the lecture.\newline 訳: 実習は、講義で紹介された考え方への理解を深める。 \\
\addlinespace[\vocabrowspace]
\en{unwind}\newline{\ipa \ipastrut /ʌnw\ipastress{vkey}{ˈ}{aɪ}{0.340750}{0.192500}nd/} & \pos{他動詞}\newline 元に戻す・解消する & 組み上げた変更や取り決めを、順に解消して元に戻す。\newline 用例: We carefully unwound the changes made during the unsuccessful upgrade.\newline 訳: 私たちは、失敗した更新で加えられた変更を慎重に元へ戻した。 \\
\addlinespace[\vocabrowspace]
\en{entrench}\newline{\ipa \ipastrut /ɪntr\ipastress{vkey}{ˈ}{ɛ}{0.115750}{0.193000}ntʃ/} & \pos{他動詞}\newline 定着させる・変えにくくする & 習慣や考えを深く定着させ、変わりにくくする。\newline 用例: Repeated exceptions can entrench inefficient habits across the organization.\newline 訳: 例外を繰り返すと、非効率な習慣が組織全体に定着することがある。 \\
\addlinespace[\vocabrowspace]
\en{reason}\newline{\ipa \ipastrut /r\ipastress{vkey}{ˈ}{i}{-0.013750}{0.392500}zən/} & \pos{自動詞}\newline 筋道を立てて考える・推論する & 根拠や関係をたどりながら、筋道を立てて考える。\newline 用例: We reasoned about the likely cause using the maintenance history.\newline 訳: 私たちは保守の履歴を使って、考えられる原因を推論した。 \\
\addlinespace[\vocabrowspace]
\en{trim}\newline{\ipa \ipastrut /trɪm/} & \pos{他動詞}\newline 不要な部分を削る・必要な項目だけに絞る & 全体から余分な部分を削り、必要な長さや量に整える。\newline 用例: The editor trimmed several repetitive paragraphs from the article.\newline 訳: 編集者は、記事から繰り返しの多い段落をいくつか削った。 \\
\addlinespace[\vocabrowspace]
\en{condense}\newline{\ipa \ipastrut /kənd\ipastress{vkey}{ˈ}{ɛ}{0.115750}{0.193000}ns/} & \pos{他動詞}\newline 短くまとめる・要約する & 長い内容の要点をまとめ、短い形にする。\newline 用例: She condensed the lecture into a concise printed handout.\newline 訳: 彼女は講義の内容を、簡潔な配布資料にまとめた。 \\
\addlinespace[\vocabrowspace]
\en{dump}\newline{\ipa \ipastrut /dʌmp/} & \pos{他動詞}\newline 大量の情報をまとめて入れる・流し込む & 大量の情報などを、一度にまとめて出したり入れたりする。\newline 用例: The diagnostic command dumps detailed statistics into a text file.\newline 訳: 診断用のコマンドは、詳細な統計をテキストファイルにまとめて出力する。 \\
\addlinespace[\vocabrowspace]
\en{pollute}\newline{\ipa \ipastrut /pəl\ipastress{vkey}{ˈ}{u}{0.160750}{0.179500}t/} & \pos{他動詞}\newline 汚す・不要な情報で乱す & 不要なものが混ざり、環境や情報の質を悪くする。\newline 用例: Automated advertising messages pollute the otherwise useful discussion forum.\newline 訳: 自動配信の広告メッセージが、せっかく有用な掲示板を荒らしている。 \\
\addlinespace[\vocabrowspace]
\en{merge}\newline{\ipa \ipastrut /mɜrdʒ/} & \pos{他動詞}\newline 統合する・取り込む & 別々になっているものを、一つのまとまりに統合する。\newline 用例: We merged the two contact lists before sending invitations.\newline 訳: 私たちは招待状を送る前に、二つの連絡先一覧を統合した。 \\
\addlinespace[\vocabrowspace]
\en{orchestrate}\newline{\ipa \ipastrut /\ipastress{vkey}{ˈ}{ɔ}{0.125500}{0.192500}rkəstr\ipastress{vhead}{ˌ}{eɪ}{0.339250}{0.701000}t/} & \pos{他動詞}\newline 複数の作業を調整してまとめる・連携させる & 複数の人や処理の動きを調整し、全体を連携させる。\newline 用例: The stage manager orchestrated lighting changes throughout the performance.\newline 訳: 舞台監督は、公演を通して照明の切り替えを調整した。 \\
\addlinespace[\vocabrowspace]
\en{achieve}\newline{\ipa \ipastrut /ətʃ\ipastress{vkey}{ˈ}{i}{-0.013750}{0.392500}v/} & \pos{他動詞}\newline 達成する・実現する & 取り組みの結果として、目標や望む状態を実現する。\newline 用例: The renovation achieved a better balance between light and privacy.\newline 訳: 改装によって、明るさとプライバシーのよりよい両立が実現した。 \\
\addlinespace[\vocabrowspace]
\en{pad}\newline{\ipa \ipastrut /pæd/} & \pos{他動詞}\newline （文章などを）水増しする・引き延ばす & 中身よりも量を増やすために、余分な内容を付け足す。\newline 用例: The writer padded the article with unrelated anecdotes.\newline 訳: その書き手は、無関係な逸話で記事を水増しした。 \\
\addlinespace[\vocabrowspace]
\en{probe}\newline{\ipa \ipastrut /proʊb/} & \pos{自動詞}\newline 探る・詳しく調べる & 表面だけでは分からないことを、さらに深く探って調べる。\newline 用例: The interviewer probed into the reasons behind the resignation.\newline 訳: 面接者は、辞職の背後にある理由を詳しく探った。 \\
\addlinespace[\vocabrowspace]
\en{burn}\newline{\ipa \ipastrut /bɜrn/} & \pos{他動詞}\newline （時間・資源などを）消費する・費やす & 限りある時間や資源を、多く費やしてしまうことを表す。\newline 用例: Repeated layout changes burned most of our preparation time.\newline 訳: 配置の変更を繰り返したことで、準備時間の大半を費やした。 \\
\addlinespace[\vocabrowspace]
\en{compress}\newline{\ipa \ipastrut /kəmpr\ipastress{vkey}{ˈ}{ɛ}{0.115750}{0.193000}s/} & \pos{他動詞}\newline 圧縮する・短くまとめる & 大きなものや長い内容を、小さな容量や短い形にまとめる。\newline 用例: The software compresses the recording before uploading it.\newline 訳: そのソフトウェアは、録音をアップロードする前に圧縮する。\newline 用例: We compressed the training material into a short illustrated guide.\newline 訳: 私たちは研修資料を、図入りの短いガイドにまとめた。 \\
\addlinespace[\vocabrowspace]
\en{pick}\newline{\ipa \ipastrut /pɪk/} & \pos{他動詞}\newline 選ぶ・選択する & いくつかある候補から、目的に合うものを選ぶ。\newline 用例: Please pick a quiet room for the recorded interview.\newline 訳: 録音する面接のために、静かな部屋を選んでください。 \\
\addlinespace[\vocabrowspace]
\en{impose}\newline{\ipa \ipastrut /ɪmp\ipastress{vkey}{ˈ}{oʊ}{0.462750}{0.192500}z/} & \pos{他動詞}\newline 押し付ける・課す & 相手に義務や制約を課したり、負担を押し付けたりする。\newline 用例: The venue imposes strict limits on amplified sound.\newline 訳: その会場は、拡声した音に厳しい制限を設けている。 \\
\addlinespace[\vocabrowspace]
\en{shape}\newline{\ipa \ipastrut /ʃeɪp/} & \pos{他動詞}\newline 形作る・方向付ける & 考え方や判断、将来の姿に影響を与えて方向付ける。\newline 用例: Early visitor feedback shaped the final layout of the exhibition.\newline 訳: 初期の来場者の意見が、展覧会の最終的な配置を方向付けた。 \\
\addlinespace[\vocabrowspace]
\en{emerge}\newline{\ipa \ipastrut /ɪm\ipastress{vkey}{ˈ}{ɜ}{0.115750}{0.193000}rdʒ/} & \pos{自動詞}\newline 現れる・明らかになる & 隠れていた事実や新しい動きが、次第に現れる。\newline 用例: A clear pattern emerged after several weeks of observation.\newline 訳: 数週間の観察を経て、明確な傾向が現れた。 \\
\addlinespace[\vocabrowspace]
\en{subsume}\newline{\ipa \ipastrut /səbs\ipastress{vkey}{ˈ}{u}{0.160750}{0.179500}m/} & \pos{他動詞}\newline 包含する・より大きな枠組みの中に含む & 小さな区分や個別のものを、より大きな枠組みに含める。\newline 用例: The new department subsumes several previously separate research groups.\newline 訳: 新しい部門は、以前は別々だった研究グループをいくつか包含している。 \\
\addlinespace[\vocabrowspace]
\en{overfit}\newline{\ipa \ipastrut /\ipastress{vkey}{ˈ}{oʊ}{0.462750}{0.192500}vərf\ipastress{vhead}{ˌ}{ɪ}{0.032500}{0.688000}t/} & \pos{自動詞}\newline 特定の細部に適合しすぎる・過剰適合する & 限られたデータに合いすぎて、ほかの場合への対応力を失う。\newline 用例: The small model overfits to unusual features in the training images.\newline 訳: その小型モデルは、学習画像の特殊な特徴に過剰適合する。 \\
\addlinespace[\vocabrowspace]
\en{deduplicate}\newline{\ipa \ipastrut UK /didj\ipastress{vkey}{ˈ}{u}{0.160750}{0.179500}plɪk\ipastress{vhead}{ˌ}{eɪ}{0.339250}{0.701000}t/} & \pos{他動詞}\newline 重複を取り除く & 同じ内容が複数あるとき、重複分を取り除く。\newline 用例: We deduplicated the mailing list before printing the address labels.\newline 訳: 私たちは宛名ラベルを印刷する前に、送付先一覧の重複を取り除いた。 \\
\addlinespace[\vocabrowspace]
\en{counteract}\newline{\ipa \ipastrut /k\ipastress{vhead}{ˌ}{aʊ}{0.465750}{0.701000}ntər\ipastress{vkey}{ˈ}{æ}{0.341750}{0.192500}kt/} & \pos{他動詞}\newline 悪影響を抑える・打ち消す & 何かの作用に対して、反対の働きで影響を弱める。\newline 用例: The window shades counteract the glare from the afternoon sun.\newline 訳: 窓の日よけが、午後の日差しのまぶしさを抑える。 \\
\addlinespace[\vocabrowspace]
\en{swap}\newline{\ipa \ipastrut /swɑp/} & \pos{他動詞}\newline 置き換える・交換する & あるものを別のものに置き換えたり、互いに交換したりする。\newline 用例: We swapped the heavy display boards for lightweight fabric banners.\newline 訳: 私たちは重い展示パネルを、軽い布製の垂れ幕に置き換えた。 \\
\addlinespace[\vocabrowspace]
\en{dismiss}\newline{\ipa \ipastrut /dɪsm\ipastress{vkey}{ˈ}{ɪ}{0.032500}{0.179500}s/} & \pos{他動詞}\newline 取り合わない・退ける & 意見や問題を、取り上げる価値がないものとして退ける。\newline 用例: The committee dismissed the complaint without examining the photographs.\newline 訳: 委員会は写真を調べずに、その苦情を退けた。 \\
\end{vtab}
\Needspace{14\baselineskip}

\section{3. 状態・範囲・要件を読む}
\begin{vtab}
\en{concern}\newline{\ipa \ipastrut /kəns\ipastress{vkey}{ˈ}{ɜ}{0.115750}{0.193000}rn/} & \pos{名詞}\newline 懸念・気にかかる事柄 & 気にかかる問題や、対応を考える必要のある事柄を指す。\newline 用例: Noise from the new equipment remains our main concern.\newline 訳: 新しい機材から出る騒音が、依然として私たちの一番の懸念だ。 \\
\addlinespace[\vocabrowspace]
\en{assumption}\newline{\ipa \ipastrut /əs\ipastress{vkey}{ˈ}{ʌ}{0.142750}{0.179500}mpʃən/} & \pos{名詞}\newline 仮定・裏づけなしに置く前提 & 十分に確かめる前に、正しいものとして置く前提を指す。\newline 用例: Our budget rests on the assumption that attendance will increase.\newline 訳: 私たちの予算は、来場者数が増えるという仮定に基づいている。 \\
\addlinespace[\vocabrowspace]
\en{likely}\newline{\ipa \ipastrut /l\ipastress{vkey}{ˈ}{aɪ}{0.340750}{0.192500}kli/} & \pos{副詞}\newline 可能性が高い・…しそうな & 予想される出来事が起こる可能性の高さを表す。\newline 用例: The meeting will likely finish before the afternoon workshop.\newline 訳: 会議はおそらく午後の研修が始まる前に終わるだろう。 \\
\addlinespace[\vocabrowspace]
\en{genuine}\newline{\ipa \ipastrut genuine /dʒ\ipastress{vkey}{ˈ}{ɛ}{0.115750}{0.193000}njuɪn/  genuinely /dʒ\ipastress{vkey}{ˈ}{ɛ}{0.115750}{0.193000}njuɪnli/} & \pos{形容詞（副詞形：genuinely）}\newline 本当の・実際の & 見せかけではなく実際にそうである性質や、その確かさを表す。\newline 用例: The noisy fan points to a genuine problem with ventilation.\newline 訳: うるさいファンは、換気に実際の問題があることを示している。\newline 用例: The revised layout is genuinely easier for visitors to navigate.\newline 訳: 改訂された配置は、来場者にとって本当に移動しやすい。 \\
\addlinespace[\vocabrowspace]
\en{current}\newline{\ipa \ipastrut /k\ipastress{vkey}{ˈ}{ɜ}{0.115750}{0.193000}rənt/} & \pos{形容詞}\newline 現在の・現行の & 今の時点で使われていたり、行われていたりするものを示す。\newline 用例: The current timetable is displayed beside the main entrance.\newline 訳: 現在の時刻表は、正面入口の横に掲示されている。 \\
\addlinespace[\vocabrowspace]
\en{prior}\newline{\ipa \ipastrut /pr\ipastress{vkey}{ˈ}{aɪ}{0.340750}{0.192500}ər/} & \pos{形容詞}\newline 前の・それ以前の & 現在の出来事より前に行われたことを指す。\newline 用例: Prior experience with spreadsheets helped her organize the survey data.\newline 訳: 以前の表計算ソフトの使用経験が、調査データの整理に役立った。 \\
\addlinespace[\vocabrowspace]
\en{incomplete}\newline{\ipa \ipastrut /\ipastress{vhead}{ˌ}{ɪ}{0.032500}{0.688000}nkəmpl\ipastress{vkey}{ˈ}{i}{-0.013750}{0.392500}t/} & \pos{形容詞}\newline 不完全な・必要な部分が欠けた & 仕上げるために必要な部分がまだ欠けていること。\newline 用例: The application remains incomplete because the address field is empty.\newline 訳: 住所欄が空なので、申請書はまだ未完成だ。 \\
\addlinespace[\vocabrowspace]
\en{approach}\newline{\ipa \ipastrut /əpr\ipastress{vkey}{ˈ}{oʊ}{0.462750}{0.192500}tʃ/} & \pos{名詞・他動詞}\newline 方法・取り組み方／取り組む・近づく & 仕事への取り組み方や、対象に近づく動作を表す。\newline 用例: Our approach combines short interviews with direct observation of daily work.\newline 訳: 私たちの方法は、短い面談と日常業務の直接観察を組み合わせている。\newline 用例: We approach unfamiliar problems by asking simple questions first.\newline 訳: 未知の問題には、まず簡単な問いを立てることから取り組む。\newline 用例: The delivery driver approached the building along the narrow path.\newline 訳: 配達員は細い小道を通って建物に近づいた。 \\
\addlinespace[\vocabrowspace]
\en{additional}\newline{\ipa \ipastrut /əd\ipastress{vkey}{ˈ}{ɪ}{0.032500}{0.179500}ʃənəl/} & \pos{形容詞}\newline 追加の・さらに加わる & すでにあるものに、さらに加わることを表す。\newline 用例: The workshop needs additional chairs for guests arriving this afternoon.\newline 訳: 今日の午後に来る参加者のために、講習会には追加の椅子が必要だ。 \\
\addlinespace[\vocabrowspace]
\en{specific}\newline{\ipa \ipastrut /spəs\ipastress{vkey}{ˈ}{ɪ}{0.032500}{0.179500}fɪk/} & \pos{形容詞}\newline 特定の・具体的な & 対象がはっきり定まっていることや、内容が具体的なこと。\newline 用例: Please give specific examples of problems you encountered during installation.\newline 訳: インストール中に起きた問題の具体例を挙げてください。\newline 用例: This key opens one specific cabinet in the laboratory.\newline 訳: この鍵で開くのは、実験室にある特定の一つの戸棚だ。 \\
\addlinespace[\vocabrowspace]
\en{recent}\newline{\ipa \ipastrut /r\ipastress{vkey}{ˈ}{i}{-0.013750}{0.392500}sənt/} & \pos{形容詞}\newline 最近の・直近の & 現在からそう遠くない過去に起きたことを指す。\newline 用例: Recent measurements show that the room becomes warmer after lunch.\newline 訳: 最近の測定では、昼食後に室温が上がることが分かった。 \\
\addlinespace[\vocabrowspace]
\en{independent}\newline{\ipa \ipastrut /\ipastress{vhead}{ˌ}{ɪ}{0.032500}{0.688000}ndɪp\ipastress{vkey}{ˈ}{ɛ}{0.115750}{0.193000}ndənt/} & \pos{形容詞}\newline 独立した・互いに依存しない & ほかのものに頼らず、それぞれが成り立っていること。\newline 用例: The two pumps have independent power supplies for emergency use.\newline 訳: 二つのポンプには、非常時用の独立した電源がある。 \\
\addlinespace[\vocabrowspace]
\en{available}\newline{\ipa \ipastrut /əv\ipastress{vkey}{ˈ}{eɪ}{0.339250}{0.192500}ləbəl/} & \pos{形容詞}\newline 利用できる・使える & 人や物を必要なときに使える状態を表す。\newline 用例: A quiet meeting room is available on the third floor.\newline 訳: 静かな会議室が三階で利用できる。 \\
\addlinespace[\vocabrowspace]
\en{relevant}\newline{\ipa \ipastrut /r\ipastress{vkey}{ˈ}{ɛ}{0.115750}{0.193000}ləvənt/} & \pos{形容詞}\newline 関係のある・目的に合った & 取り上げている話題や、達成したい目的に関係すること。\newline 用例: She attached the photographs relevant to the damaged roof.\newline 訳: 彼女は、傷んだ屋根に関係する写真を添付した。 \\
\addlinespace[\vocabrowspace]
\en{evidence}\newline{\ipa \ipastrut /\ipastress{vkey}{ˈ}{ɛ}{0.115750}{0.193000}vɪdəns/} & \pos{名詞}\newline 根拠・証拠 & 説明や判断を裏づける、観察結果や記録などの情報。\newline 用例: The soil samples provided evidence of earlier flooding nearby.\newline 訳: 土壌の試料は、付近で以前に洪水があった証拠となった。 \\
\addlinespace[\vocabrowspace]
\en{accurate}\newline{\ipa \ipastrut /\ipastress{vkey}{ˈ}{æ}{0.341750}{0.192500}kjərət/} & \pos{形容詞}\newline 正確な & 記録や説明の内容が、実際の状態と合っていること。\newline 用例: The inventory must remain accurate when equipment moves between offices.\newline 訳: 備品を事務所間で移動するときも、在庫記録は正確でなければならない。 \\
\addlinespace[\vocabrowspace]
\en{unreliable}\newline{\ipa \ipastrut /\ipastress{vhead}{ˌ}{ʌ}{0.142750}{0.688000}nrɪl\ipastress{vkey}{ˈ}{aɪ}{0.340750}{0.192500}əbəl/} & \pos{形容詞}\newline 信頼できない・当てにできない & 結果や動作が安定せず、安心して頼れないこと。\newline 用例: The old sensor became unreliable during periods of heavy rain.\newline 訳: 古いセンサーは、大雨が続くと当てにならなくなった。 \\
\addlinespace[\vocabrowspace]
\en{precise}\newline{\ipa \ipastrut /prɪs\ipastress{vkey}{ˈ}{aɪ}{0.340750}{0.192500}s/} & \pos{形容詞}\newline 正確な・細かく特定された & 位置や数量などが、細かいところまで明確に定まっていること。\newline 用例: The diagram shows the precise location of each mounting hole.\newline 訳: 図には、取り付け穴それぞれの正確な位置が示されている。 \\
\addlinespace[\vocabrowspace]
\en{intermediate}\newline{\ipa \ipastrut /\ipastress{vhead}{ˌ}{ɪ}{0.032500}{0.688000}ntərm\ipastress{vkey}{ˈ}{i}{-0.013750}{0.392500}diət/} & \pos{形容詞}\newline 途中の・中間の & 始まりと完成の間にある段階を表す。\newline 用例: Save intermediate results so colleagues can inspect the calculation process.\newline 訳: 同僚が計算の過程を確認できるよう、中間結果を保存してください。 \\
\addlinespace[\vocabrowspace]
\en{final}\newline{\ipa \ipastrut /f\ipastress{vkey}{ˈ}{aɪ}{0.340750}{0.192500}nəl/} & \pos{形容詞}\newline 最後の・最終の & 一連の作業や順序の、いちばん終わりに当たること。\newline 用例: The final draft includes a new illustration of the building.\newline 訳: 最終稿には、その建物の新しいイラストが入っている。 \\
\addlinespace[\vocabrowspace]
\en{fluency}\newline{\ipa \ipastrut /fl\ipastress{vkey}{ˈ}{u}{0.160750}{0.179500}ənsi/} & \pos{名詞}\newline 熟達・使いこなす力 & 知識や道具を、無理なく使いこなせる熟達の度合い。\newline 用例: Regular practice improved her fluency with the new design software.\newline 訳: 定期的に練習することで、新しい設計ソフトを使いこなせるようになった。 \\
\addlinespace[\vocabrowspace]
\en{competency}\newline{\ipa \ipastrut /k\ipastress{vkey}{ˈ}{ɑ}{0.146500}{0.192500}mpətənsi/} & \pos{名詞}\newline 実践に必要な能力 & ある仕事を実際に進めるために求められる能力。\newline 用例: Clear communication is a core competency for workshop instructors.\newline 訳: 分かりやすく伝える力は、講習会の講師に必要な中心的能力だ。 \\
\addlinespace[\vocabrowspace]
\en{Description}\newline{\ipa \ipastrut /dɪskr\ipastress{vkey}{ˈ}{ɪ}{0.032500}{0.179500}pʃən/} & \pos{名詞（能力名）}\newline 記述・意図を明確に伝えること & 望む成果や進め方を、相手に分かるよう言葉にする力。\newline 用例: Description helps collaborators understand the intended result and working process.\newline 訳: 意図を記述する力は、共同作業者が望まれる成果と進め方を理解する助けになる。 \\
\addlinespace[\vocabrowspace]
\en{Discernment}\newline{\ipa \ipastrut /dɪs\ipastress{vkey}{ˈ}{ɜ}{0.115750}{0.193000}rnmənt/} & \pos{名詞（能力名）}\newline 見極め・吟味する力 & 受け取った情報や成果を注意深く確かめ、見極める力。\newline 用例: Discernment helps us recognize weak evidence beneath a persuasive presentation.\newline 訳: 見極める力があれば、説得力のある発表でも根拠の弱さに気づける。 \\
\addlinespace[\vocabrowspace]
\en{Diligence}\newline{\ipa \ipastrut /d\ipastress{vkey}{ˈ}{ɪ}{0.032500}{0.179500}lədʒəns/} & \pos{名詞（能力名）}\newline 注意深く責任を果たすこと & 必要な確認を怠らず、責任を持って仕事に取り組むこと。\newline 用例: Diligence means checking permissions before sharing research materials with visitors.\newline 訳: 注意深く責任を果たすには、訪問者に研究資料を共有する前に許可を確認する。 \\
\addlinespace[\vocabrowspace]
\en{intentions}\newline{\ipa \ipastrut /ɪnt\ipastress{vkey}{ˈ}{ɛ}{0.115750}{0.193000}nʃənz/} & \pos{名詞（複数形）}\newline 意図・実現したいこと & 何を行い、どんな結果を実現したいかという考え。\newline 用例: The artist explained her intentions before the exhibition opened.\newline 訳: その画家は、展覧会の開始前に自分の制作意図を説明した。 \\
\addlinespace[\vocabrowspace]
\en{audience}\newline{\ipa \ipastrut /\ipastress{vkey}{ˈ}{ɔ}{0.125500}{0.192500}diəns/} & \pos{名詞}\newline 受け手・想定読者 & 文章や発表などを受け取る、想定された相手。\newline 用例: The guide uses simple language for an audience of beginners.\newline 訳: その手引きは、初心者を読者として想定し、平易な言葉を使っている。 \\
\addlinespace[\vocabrowspace]
\en{concise}\newline{\ipa \ipastrut /kəns\ipastress{vkey}{ˈ}{aɪ}{0.340750}{0.192500}s/} & \pos{形容詞}\newline 簡潔な & 必要な内容を保ちながら、短くまとまっていること。\newline 用例: A concise update helped everyone understand the schedule change quickly.\newline 訳: 簡潔な連絡のおかげで、全員が日程変更をすぐに理解できた。 \\
\addlinespace[\vocabrowspace]
\en{coherence}\newline{\ipa \ipastrut /koʊh\ipastress{vkey}{ˈ}{ɪ}{0.032500}{0.179500}rəns/} & \pos{名詞}\newline 筋が通り、全体がつながっていること & 各部分がつながり、全体として筋が通っている性質。\newline 用例: Moving the final paragraph improved the coherence of her argument.\newline 訳: 最後の段落を移したことで、彼女の論述のつながりがよくなった。 \\
\addlinespace[\vocabrowspace]
\en{judgment}\newline{\ipa \ipastrut /dʒ\ipastress{vkey}{ˈ}{ʌ}{0.142750}{0.179500}dʒmənt/} & \pos{名詞}\newline 判断・判断力 & 状況や根拠を踏まえて、適切な選択をすること。\newline 用例: The team relied on her judgment when selecting laboratory equipment.\newline 訳: 研究室の機器を選ぶ際、チームは彼女の判断を頼りにした。 \\
\addlinespace[\vocabrowspace]
\en{oversight}\newline{\ipa \ipastrut /\ipastress{vkey}{ˈ}{oʊ}{0.462750}{0.192500}vərs\ipastress{vhead}{ˌ}{aɪ}{0.340750}{0.701000}t/} & \pos{名詞}\newline 監督・見守って確認すること & 仕事の進め方を監督すること。確認時の見落としも指す。\newline 用例: The project needs independent oversight throughout the construction period.\newline 訳: 工事期間を通して、その事業には独立した立場からの監督が必要だ。\newline 用例: An oversight left the attachment out of yesterday's email.\newline 訳: 見落としにより、昨日のメールには添付ファイルが入っていなかった。 \\
\addlinespace[\vocabrowspace]
\en{transparency}\newline{\ipa \ipastrut /trænsp\ipastress{vkey}{ˈ}{ɛ}{0.115750}{0.193000}rənsi/} & \pos{名詞}\newline 透明性・関与や経緯を明らかにすること & 判断の過程や関与した人などが、明らかになっていること。\newline 用例: Publishing the selection criteria improved transparency for all applicants.\newline 訳: 選考基準を公開したことで、すべての応募者にとって透明性が高まった。 \\
\addlinespace[\vocabrowspace]
\en{disclosure}\newline{\ipa \ipastrut /dɪskl\ipastress{vkey}{ˈ}{oʊ}{0.462750}{0.192500}ʒər/} & \pos{名詞}\newline 開示・明らかにすること & それまで知られていなかった情報を、相手に明らかにすること。\newline 用例: The report includes a disclosure of the researchers' funding sources.\newline 訳: 報告書では、研究者たちの資金提供元が開示されている。 \\
\addlinespace[\vocabrowspace]
\en{ethical}\newline{\ipa \ipastrut /\ipastress{vkey}{ˈ}{ɛ}{0.115750}{0.193000}θɪkəl/} & \pos{形容詞}\newline 倫理にかなった・倫理上の & 人への影響や、公正さなどの倫理に関係すること。\newline 用例: The committee discussed ethical concerns about collecting children's location data.\newline 訳: 委員会は、子どもの位置情報を集めることの倫理的な懸念を話し合った。 \\
\addlinespace[\vocabrowspace]
\en{complementary}\newline{\ipa \ipastrut /k\ipastress{vhead}{ˌ}{ɑ}{0.146500}{0.701000}mpləm\ipastress{vkey}{ˈ}{ɛ}{0.115750}{0.193000}ntəri/} & \pos{形容詞}\newline 互いに補い合う & 異なる強みを組み合わせて、互いの不足を補うこと。\newline 用例: The designer and engineer brought complementary skills to the project.\newline 訳: 設計者と技術者は、互いを補う技能をプロジェクトに持ち寄った。 \\
\addlinespace[\vocabrowspace]
\en{expectations}\newline{\ipa \ipastrut /\ipastress{vhead}{ˌ}{ɛ}{0.115750}{0.701500}kspɛkt\ipastress{vkey}{ˈ}{eɪ}{0.339250}{0.192500}ʃənz/} & \pos{名詞（複数形）}\newline 期待・求められていること & 相手や成果に対して、実現してほしいと考えていること。\newline 用例: We discussed our expectations before inviting volunteers to the event.\newline 訳: ボランティアを催しに招く前に、何を期待するかを話し合った。 \\
\addlinespace[\vocabrowspace]
\en{accountability}\newline{\ipa \ipastrut /ək\ipastress{vhead}{ˌ}{aʊ}{0.465750}{0.701000}ntəb\ipastress{vkey}{ˈ}{ɪ}{0.032500}{0.179500}ləti/} & \pos{名詞}\newline 判断や結果について責任を負い、説明すること & 自分の判断や結果に責任を持ち、説明できること。\newline 用例: Assigning an owner to each decision strengthens accountability within the team.\newline 訳: 判断ごとに責任者を定めると、チーム内の説明責任が明確になる。 \\
\addlinespace[\vocabrowspace]
\en{critically}\newline{\ipa \ipastrut /kr\ipastress{vkey}{ˈ}{ɪ}{0.032500}{0.179500}tɪkli/} & \pos{副詞}\newline 根拠や妥当性を吟味して & 情報をそのまま受け取らず、根拠や妥当性を吟味して。\newline 用例: Students examined the advertisement critically before discussing its environmental claims.\newline 訳: 学生たちは広告を吟味してから、そこにある環境に関する主張を話し合った。 \\
\addlinespace[\vocabrowspace]
\en{steerable}\newline{\ipa \ipastrut /st\ipastress{vkey}{ˈ}{ɪ}{0.032500}{0.179500}rəbəl/} & \pos{形容詞}\newline 指示によって方向づけられる & 指示を与えることで、進み方や振る舞いを調整できること。\newline 用例: The writing assistant is steerable through clear instructions about tone.\newline 訳: 文章作成アシスタントは、口調について明確に指示すると調整できる。 \\
\addlinespace[\vocabrowspace]
\en{feedback}\newline{\ipa \ipastrut /f\ipastress{vkey}{ˈ}{i}{-0.013750}{0.392500}db\ipastress{vhead}{ˌ}{æ}{0.341750}{0.701000}k/} & \pos{名詞}\newline 反応・改善のために返す意見 & 相手の成果や進め方に対して返す、改善に役立つ意見。\newline 用例: Her feedback helped the presenter simplify a confusing diagram.\newline 訳: 彼女の意見は、発表者が分かりにくい図を簡潔にする助けになった。 \\
\addlinespace[\vocabrowspace]
\en{constraints}\newline{\ipa \ipastrut /kənstr\ipastress{vkey}{ˈ}{eɪ}{0.339250}{0.192500}nts/} & \pos{名詞（複数形）}\newline 制約・守るべき条件 & 計画や作業を進める際に、満たさなければならない条件。\newline 用例: Budget constraints limited the number of samples we could collect.\newline 訳: 予算上の制約で、集められる試料の数が限られた。 \\
\addlinespace[\vocabrowspace]
\en{plausible}\newline{\ipa \ipastrut /pl\ipastress{vkey}{ˈ}{ɔ}{0.125500}{0.192500}zəbəl/} & \pos{形容詞}\newline もっともらしい & 説明などが、筋が通っていて本当らしく見えること。\newline 用例: A loose cable seemed a plausible explanation for the flickering light.\newline 訳: 照明のちらつきは、ケーブルの緩みで説明できそうだった。 \\
\addlinespace[\vocabrowspace]
\en{expertise}\newline{\ipa \ipastrut /\ipastress{vhead}{ˌ}{ɛ}{0.115750}{0.701500}kspərt\ipastress{vkey}{ˈ}{i}{-0.013750}{0.392500}z/} & \pos{名詞}\newline 専門知識・経験に基づく技能 & 経験や学習によって身についた、分野に特有の知識や技能。\newline 用例: We sought her expertise when restoring the damaged wooden staircase.\newline 訳: 傷んだ木製階段を修復する際、私たちは彼女の専門知識を求めた。 \\
\addlinespace[\vocabrowspace]
\en{generic}\newline{\ipa \ipastrut /dʒən\ipastress{vkey}{ˈ}{ɛ}{0.115750}{0.193000}rɪk/} & \pos{形容詞}\newline 一般的すぎる・個別事情を反映していない & 個々の事情に合わせた内容が少なく、広く一般的であること。\newline 用例: The generic advice offered little help with our unusual floor plan.\newline 訳: その一般的な助言は、私たちの変わった間取りにはあまり役立たなかった。 \\
\addlinespace[\vocabrowspace]
\en{recurring}\newline{\ipa \ipastrut /rɪk\ipastress{vkey}{ˈ}{ɜ}{0.115750}{0.193000}rɪŋ/} & \pos{形容詞}\newline 繰り返し生じる & 同じような出来事や作業が、繰り返し発生すること。\newline 用例: The team scheduled recurring meetings to discuss progress on the exhibition.\newline 訳: チームは展覧会の準備状況を話し合う定例会議を設定した。 \\
\addlinespace[\vocabrowspace]
\en{objective}\newline{\ipa \ipastrut /əbdʒ\ipastress{vkey}{ˈ}{ɛ}{0.115750}{0.193000}ktɪv/} & \pos{形容詞・名詞}\newline 客観的な／目標・目的 & 事実に基づいていること。達成を目指す目標も表す。\newline 用例: An objective assessment should consider measurements from several different sources.\newline 訳: 客観的な評価では、複数の異なる情報源による測定値を考慮するべきだ。\newline 用例: Our main objective is to reduce waste in the kitchen.\newline 訳: 私たちの主な目標は、厨房で出る廃棄物を減らすことだ。 \\
\addlinespace[\vocabrowspace]
\en{criteria}\newline{\ipa \ipastrut /kraɪt\ipastress{vkey}{ˈ}{ɪ}{0.032500}{0.179500}riə/} & \pos{名詞（複数形）}\newline 判断・評価の基準 & 良し悪しや適否を判断するために用いる基準。\newline 用例: The committee agreed on clear criteria for selecting the venue.\newline 訳: 委員会は、会場を選ぶための明確な基準に合意した。 \\
\addlinespace[\vocabrowspace]
\en{thoroughly}\newline{\ipa \ipastrut /θ\ipastress{vkey}{ˈ}{ɜ}{0.115750}{0.193000}roʊli/} & \pos{副詞}\newline 十分に・漏れなく詳しく & 必要な範囲を、細部まで十分に調べたり扱ったりして。\newline 用例: The technician inspected the equipment thoroughly before the training session.\newline 訳: 技術者は、研修の前に機器を入念に点検した。 \\
\addlinespace[\vocabrowspace]
\en{ambiguity}\newline{\ipa \ipastrut /\ipastress{vhead}{ˌ}{æ}{0.341750}{0.701000}mbɪɡj\ipastress{vkey}{ˈ}{u}{0.160750}{0.179500}əti/} & \pos{名詞}\newline 曖昧さ・複数の解釈ができること & 言葉や状況が、複数の意味に受け取れること。\newline 用例: Adding a date removed the ambiguity from the delivery instructions.\newline 訳: 日付を加えたことで、配送指示の曖昧さがなくなった。 \\
\addlinespace[\vocabrowspace]
\en{attempt}\newline{\ipa \ipastrut /ət\ipastress{vkey}{ˈ}{ɛ}{0.115750}{0.193000}mpt/} & \pos{名詞}\newline 試み・実行しようとすること & 何かを実現しようと試してみる、その一回の取り組み。\newline 用例: Our first attempt produced a working model of the bridge.\newline 訳: 最初の試みで、機能する橋の模型ができた。 \\
\addlinespace[\vocabrowspace]
\en{claim}\newline{\ipa \ipastrut /kleɪm/} & \pos{名詞}\newline 主張・事実として述べられた内容 & 真実であるとして述べられた内容や主張。\newline 用例: The manufacturer supported its durability claim with independent test results.\newline 訳: 製造元は、独立した試験の結果によって耐久性の主張を裏づけた。 \\
\addlinespace[\vocabrowspace]
\en{issue}\newline{\ipa \ipastrut /\ipastress{vkey}{ˈ}{ɪ}{0.032500}{0.179500}ʃu/} & \pos{名詞・他動詞}\newline 問題・検討する事柄／発行する & 対処や検討が必要な事柄。文書などを発行することも表す。\newline 用例: The team discussed an issue with lighting in the entrance hall.\newline 訳: チームは、玄関ホールの照明に関する問題を話し合った。\newline 用例: The office will issue new access cards to all employees.\newline 訳: 事務局は、全従業員に新しい入館カードを発行する。 \\
\addlinespace[\vocabrowspace]
\en{findings}\newline{\ipa \ipastrut /f\ipastress{vkey}{ˈ}{aɪ}{0.340750}{0.192500}ndɪŋz/} & \pos{名詞（複数形）}\newline 調査で分かったこと・検出された指摘 & 調べたり確かめたりすることで明らかになった結果。\newline 用例: The researchers shared their findings at a local science festival.\newline 訳: 研究者たちは、地域の科学祭で調査結果を発表した。 \\
\addlinespace[\vocabrowspace]
\en{scope}\newline{\ipa \ipastrut /skoʊp/} & \pos{名詞}\newline 範囲・担当する範囲 & 仕事や調査で、どこまでを対象として扱うかという範囲。\newline 用例: The scope of this survey includes every public playground nearby.\newline 訳: この調査は、近隣のすべての公園内の遊び場を対象とする。 \\
\addlinespace[\vocabrowspace]
\en{outcome}\newline{\ipa \ipastrut /\ipastress{vkey}{ˈ}{aʊ}{0.465750}{0.192500}tkʌm/} & \pos{名詞}\newline 結果・実行後に生じた状態 & 取り組みや出来事の後に得られた結果や状態。\newline 用例: The final outcome depended on careful preparation by the whole team.\newline 訳: 最終的な結果は、チーム全体による入念な準備にかかっていた。 \\
\addlinespace[\vocabrowspace]
\en{tradeoff}\newline{\ipa \ipastrut /tr\ipastress{vkey}{ˈ}{eɪ}{0.339250}{0.192500}d\ipastress{vhead}{ˌ}{ɔ}{0.125500}{0.701000}f/} & \pos{名詞}\newline 一方を得ると他方を失う関係・兼ね合い & 一方を良くすると、もう一方が犠牲になる兼ね合い。\newline 用例: Choosing lighter materials involves a tradeoff between weight and durability.\newline 訳: 軽い素材を選ぶ際には、重量と耐久性の兼ね合いが生じる。 \\
\addlinespace[\vocabrowspace]
\en{consistent}\newline{\ipa \ipastrut /kəns\ipastress{vkey}{ˈ}{ɪ}{0.032500}{0.179500}stənt/} & \pos{形容詞}\newline 一貫した・整合している & 基準や内容がそろい、互いに食い違っていないこと。\newline 用例: The labels use consistent colors throughout the museum's different galleries.\newline 訳: ラベルは、美術館の各展示室を通して一貫した配色を使っている。 \\
\addlinespace[\vocabrowspace]
\en{sufficient}\newline{\ipa \ipastrut /səf\ipastress{vkey}{ˈ}{ɪ}{0.032500}{0.179500}ʃənt/} & \pos{形容詞}\newline 十分な・必要なだけの & ある目的に対して、必要な量や程度を満たしていること。\newline 用例: The battery holds sufficient power for an afternoon of fieldwork.\newline 訳: その電池には、午後の野外調査に十分な電力がある。 \\
\addlinespace[\vocabrowspace]
\en{explicit}\newline{\ipa \ipastrut /ɪkspl\ipastress{vkey}{ˈ}{ɪ}{0.032500}{0.179500}sɪt/} & \pos{形容詞}\newline 明示的な・はっきり示された & 意図や条件が、言葉などではっきり示されていること。\newline 用例: The invitation gives explicit instructions about entering the building after hours.\newline 訳: 招待状には、時間外に建物へ入る方法が明記されている。 \\
\addlinespace[\vocabrowspace]
\en{entire}\newline{\ipa \ipastrut /ɪnt\ipastress{vkey}{ˈ}{aɪ}{0.340750}{0.192500}ər/} & \pos{形容詞}\newline 全体の・丸ごとの & 対象の一部分ではなく、初めから終わりまでの全体。\newline 用例: The entire collection fits inside a single temperature-controlled storage room.\newline 訳: コレクション全体が、一つの温度管理された保管室に収まる。 \\
\addlinespace[\vocabrowspace]
\en{identical}\newline{\ipa \ipastrut /aɪd\ipastress{vkey}{ˈ}{ɛ}{0.115750}{0.193000}ntəkəl/} & \pos{形容詞}\newline まったく同じ・同一の & 比べたものの形や内容に、まったく違いがないこと。\newline 用例: Both packages contain identical parts for assembling the shelving unit.\newline 訳: どちらの箱にも、棚を組み立てるための同一の部品が入っている。 \\
\addlinespace[\vocabrowspace]
\en{minimal}\newline{\ipa \ipastrut /m\ipastress{vkey}{ˈ}{ɪ}{0.032500}{0.179500}nɪməl/} & \pos{形容詞}\newline 最小限の・ごく少ない & 量や程度が、ごく小さい範囲にとどまっていること。\newline 用例: The new layout required minimal changes to the existing furniture.\newline 訳: 新しい配置にするために必要な、既存の家具の変更は最小限だった。 \\
\addlinespace[\vocabrowspace]
\en{subsequent}\newline{\ipa \ipastrut /s\ipastress{vkey}{ˈ}{ʌ}{0.142750}{0.179500}bsɪkwənt/} & \pos{形容詞}\newline その後の・後に続く & 基準となる出来事より後に続いて起こること。\newline 用例: Subsequent interviews revealed why visitors avoided the upstairs gallery.\newline 訳: その後の面談で、来館者が二階の展示室を避ける理由が分かった。 \\
\addlinespace[\vocabrowspace]
\en{sensitive}\newline{\ipa \ipastrut /s\ipastress{vkey}{ˈ}{ɛ}{0.115750}{0.193000}nsətɪv/} & \pos{形容詞}\newline 慎重な扱いを要する・影響を受けやすい & 慎重な扱いが必要なことや、小さな変化にも反応すること。\newline 用例: Store sensitive employee information in a folder with restricted access.\newline 訳: 慎重な扱いが必要な従業員情報は、アクセスを制限したフォルダーに保存してください。\newline 用例: These delicate instruments are sensitive to sudden changes in temperature.\newline 訳: これらの精密機器は、急激な温度変化の影響を受けやすい。 \\
\addlinespace[\vocabrowspace]
\en{requirement}\newline{\ipa \ipastrut /rɪkw\ipastress{vkey}{ˈ}{aɪ}{0.340750}{0.192500}rmənt/} & \pos{名詞}\newline 要件・満たすべき条件 & 利用や完成のために、満たす必要がある条件。\newline 用例: A stable internet connection is a requirement for the online workshop.\newline 訳: 安定したインターネット接続が、オンライン講習会への参加条件だ。 \\
\addlinespace[\vocabrowspace]
\en{uncertainty}\newline{\ipa \ipastrut /ʌns\ipastress{vkey}{ˈ}{ɜ}{0.115750}{0.193000}rtənti/} & \pos{名詞}\newline 不確実性・状況や進め方がまだはっきりしないこと & 結果や今後の進め方について、まだ確かなことが分からない状態。\newline 用例: Uncertainty about delivery dates made planning the exhibition more difficult.\newline 訳: 納品日の不確かさで、展覧会の計画が立てにくくなった。 \\
\addlinespace[\vocabrowspace]
\en{invoice}\newline{\ipa \ipastrut /\ipastress{vkey}{ˈ}{ɪ}{0.032500}{0.179500}nvɔɪs/} & \pos{名詞}\newline 商品・サービスの代金を請求する書類・請求書 & 提供した商品やサービスの代金を請求する書類。\newline 用例: The printer sent an invoice after delivering the exhibition brochures.\newline 訳: 印刷会社は、展覧会のパンフレットを納品してから請求書を送った。 \\
\addlinespace[\vocabrowspace]
\en{refund}\newline{\ipa \ipastrut /r\ipastress{vkey}{ˈ}{i}{-0.013750}{0.392500}fʌnd/} & \pos{名詞}\newline 支払ったお金を返すこと・返金 & 支払い済みのお金を返すことや、その返されたお金。\newline 用例: The venue offered a refund after canceling the evening concert.\newline 訳: 会場は、夜のコンサートを中止した後、払い戻しを申し出た。 \\
\addlinespace[\vocabrowspace]
\en{eligibility}\newline{\ipa \ipastrut /\ipastress{vhead}{ˌ}{ɛ}{0.115750}{0.701500}lədʒəb\ipastress{vkey}{ˈ}{ɪ}{0.032500}{0.179500}ləti/} & \pos{名詞}\newline 対象となるための条件を満たしていること & 制度の対象になるための条件を満たしていること。\newline 用例: Applicants can check their eligibility for the training grant online.\newline 訳: 申請者は、研修助成金の対象資格があるかをオンラインで確認できる。 \\
\addlinespace[\vocabrowspace]
\en{frequently}\newline{\ipa \ipastrut /fr\ipastress{vkey}{ˈ}{i}{-0.013750}{0.392500}kwəntli/} & \pos{副詞}\newline 頻繁に・たびたび & ある出来事が、高い頻度で繰り返し起こること。\newline 用例: The team frequently visits local workshops to learn new techniques.\newline 訳: チームは新しい技術を学ぶため、地元の工房を頻繁に訪れる。 \\
\addlinespace[\vocabrowspace]
\en{contradiction}\newline{\ipa \ipastrut /k\ipastress{vhead}{ˌ}{ɑ}{0.146500}{0.701000}ntrəd\ipastress{vkey}{ˈ}{ɪ}{0.032500}{0.179500}kʃən/} & \pos{名詞}\newline 内容が両立しないこと・矛盾 & 二つの内容を、同時に正しいものとして認められないこと。\newline 用例: The witness noticed a contradiction between the two written accounts.\newline 訳: 目撃者は、二つの記述の間に矛盾があることに気づいた。 \\
\addlinespace[\vocabrowspace]
\en{figure}\newline{\ipa \ipastrut /f\ipastress{vkey}{ˈ}{ɪ}{0.032500}{0.179500}ɡjər/} & \pos{名詞}\newline 数値・数字 & 測定や集計などで示される、数値そのもの。\newline 用例: The latest attendance figure includes visitors who arrived after lunch.\newline 訳: 最新の来場者数には、昼食後に来た人も含まれている。 \\
\addlinespace[\vocabrowspace]
\en{prose}\newline{\ipa \ipastrut /proʊz/} & \pos{名詞（修飾用法）}\newline 通常の文章・文章形式の記述 & 普通の文章として書かれた、連続した記述。\newline 用例: Each photograph is accompanied by a short prose description.\newline 訳: それぞれの写真には、短い文章形式の説明が添えられている。 \\
\addlinespace[\vocabrowspace]
\en{branch}\newline{\ipa \ipastrut /bræntʃ/} & \pos{名詞}\newline 枝・分岐・分岐先 & 一本の流れから分かれた枝や、作業の分岐先。\newline 用例: We tested the new animation on a separate development branch.\newline 訳: 新しいアニメーションは、別の開発ブランチで試した。 \\
\addlinespace[\vocabrowspace]
\en{run}\newline{\ipa \ipastrut /rʌn/} & \pos{名詞}\newline （処理・システムの）実行・試行 & 処理や実験などを、一回実行すること。\newline 用例: The first run processed all the archived images without errors.\newline 訳: 最初の実行では、保存済みの画像をすべてエラーなく処理できた。 \\
\addlinespace[\vocabrowspace]
\en{dispute}\newline{\ipa \ipastrut /dɪspj\ipastress{vkey}{ˈ}{u}{0.160750}{0.179500}t/} & \pos{名詞}\newline 争い・異議・意見の食い違い & 当事者の意見や要求が食い違い、争いになっていること。\newline 用例: The neighbors resolved their dispute over the shared garden gate.\newline 訳: 隣人たちは、共用の庭の門をめぐる争いを解決した。 \\
\addlinespace[\vocabrowspace]
\en{resolution}\newline{\ipa \ipastrut /r\ipastress{vhead}{ˌ}{ɛ}{0.115750}{0.701500}zəl\ipastress{vkey}{ˈ}{u}{0.160750}{0.179500}ʃən/} & \pos{名詞}\newline 問題の解決・決着 & 問題や争いに決着がつき、解決に至ること。\newline 用例: Both sides welcomed the resolution of the long-running scheduling dispute.\newline 訳: 双方とも、長引いていた日程の争いが決着したことを歓迎した。 \\
\addlinespace[\vocabrowspace]
\en{replacement}\newline{\ipa \ipastrut /rɪpl\ipastress{vkey}{ˈ}{eɪ}{0.339250}{0.192500}smənt/} & \pos{名詞}\newline 代わりの品・交換品 & 元の物に代わる品や、別のものに取り替える行為。\newline 用例: The shop ordered a replacement for the cracked display panel.\newline 訳: 店は、ひび割れた表示パネルの交換品を注文した。\newline 用例: Replacement of the damaged pipes will begin early next week.\newline 訳: 破損した配管の交換は、来週の初めに始まる。 \\
\addlinespace[\vocabrowspace]
\en{cap}\newline{\ipa \ipastrut /kæp/} & \pos{名詞・他動詞}\newline 上限・限度／上限を設ける・…までに制限する & 量や回数などの上限、または上限を定めること。\newline 用例: The workshop has a cap on the number of participants.\newline 訳: その講習会には、参加人数の上限がある。\newline 用例: The organizer capped attendance at twelve people for the pottery class.\newline 訳: 主催者は、陶芸教室の参加人数を十二人までに制限した。 \\
\addlinespace[\vocabrowspace]
\en{mid-investigation}\newline{\ipa \ipastrut mid- /mɪd/  investigation /ɪnv\ipastress{vhead}{ˌ}{ɛ}{0.115750}{0.701500}stəɡ\ipastress{vkey}{ˈ}{eɪ}{0.339250}{0.192500}ʃən/} & \pos{副詞的表現（接頭辞＋名詞）}\newline 調査の途中で & 調査を始めてから完了するまでの途中を表す。\newline 用例: The inspectors discovered another leak mid-investigation and extended their visit.\newline 訳: 検査員たちは調査の途中で別の水漏れを発見し、滞在を延長した。 \\
\addlinespace[\vocabrowspace]
\en{pending}\newline{\ipa \ipastrut /p\ipastress{vkey}{ˈ}{ɛ}{0.115750}{0.193000}ndɪŋ/} & \pos{形容詞}\newline 処理待ちの・完了待ちの・保留中の & 対応や判断がまだ終わらず、結果を待っている状態。\newline 用例: The application remains pending until the committee meets next Tuesday.\newline 訳: 申請は、来週火曜日の委員会まで審査待ちのままだ。 \\
\addlinespace[\vocabrowspace]
\en{occasionally}\newline{\ipa \ipastrut /ək\ipastress{vkey}{ˈ}{eɪ}{0.339250}{0.192500}ʒənəli/} & \pos{副詞}\newline 時々・たまに & 常にではなく、ときどき出来事が起こること。\newline 用例: We occasionally invite local artists to join our lunchtime discussions.\newline 訳: 私たちはときどき、地元の芸術家を昼食時の話し合いに招く。 \\
\addlinespace[\vocabrowspace]
\en{wording}\newline{\ipa \ipastrut /w\ipastress{vkey}{ˈ}{ɜ}{0.115750}{0.193000}rdɪŋ/} & \pos{名詞}\newline 言い回し・使う言葉や表現 & 伝える内容に対して、どのような言葉や表現を使うか。\newline 用例: The editor changed the wording to make the announcement friendlier.\newline 訳: 編集者は、お知らせが親しみやすくなるよう言い回しを変えた。 \\
\addlinespace[\vocabrowspace]
\en{ordering}\newline{\ipa \ipastrut /\ipastress{vkey}{ˈ}{ɔ}{0.125500}{0.192500}rdərɪŋ/} & \pos{名詞（動名詞）}\newline 並べる順序・並び方 & 項目を並べる順序や、商品などを注文する行為。\newline 用例: The ordering of the photographs follows the route through town.\newline 訳: 写真の並び順は、町を通る道順に沿っている。\newline 用例: Online ordering makes it easier to buy supplies for remote offices.\newline 訳: オンライン注文を使うと、遠方の事務所の備品を買いやすくなる。 \\
\addlinespace[\vocabrowspace]
\en{dozens of}\newline{\ipa \ipastrut /d\ipastress{vkey}{ˈ}{ʌ}{0.142750}{0.179500}zənz əv/} & \pos{数量表現}\newline 何十もの・数十の & 正確な個数を定めず、数十もの多数を表す。\newline 用例: Dozens of volunteers helped clean the riverbank after the festival.\newline 訳: 祭りの後、何十人ものボランティアが川岸の清掃を手伝った。 \\
\addlinespace[\vocabrowspace]
\en{30+}\newline{\ipa \ipastrut /θ\ipastress{vkey}{ˈ}{ɜ}{0.115750}{0.193000}rti plʌs/} & \pos{数量表現}\newline 30以上・30個以上の & 数値に添えられた記号によって、三十以上であることを表す。\newline 用例: The library offers 30+ evening workshops for adults each year.\newline 訳: 図書館は毎年、成人向けの夜間講習会を三十回以上開いている。 \\
\addlinespace[\vocabrowspace]
\en{raw}\newline{\ipa \ipastrut /rɑ/} & \pos{形容詞}\newline 生の・未加工の & 整理や加工を加える前の、そのままの状態。\newline 用例: The researchers saved the raw recordings before editing the interviews.\newline 訳: 研究者たちは、面談を編集する前に未加工の録音を保存した。 \\
\addlinespace[\vocabrowspace]
\en{hallmark}\newline{\ipa \ipastrut /h\ipastress{vkey}{ˈ}{ɔ}{0.125500}{0.192500}lm\ipastress{vhead}{ˌ}{ɑ}{0.146500}{0.701000}rk/} & \pos{名詞}\newline 典型的な特徴・特徴となるしるし & あるものを特徴づける、代表的なしるしや性質。\newline 用例: Careful attention to lighting is a hallmark of her photography.\newline 訳: 照明への細やかな配慮が、彼女の写真の特徴だ。 \\
\addlinespace[\vocabrowspace]
\en{unusual}\newline{\ipa \ipastrut /ʌnj\ipastress{vkey}{ˈ}{u}{0.160750}{0.179500}ʒuəl/} & \pos{形容詞}\newline 通常と異なる・珍しい & 普段あまり見られず、通常とは異なっていること。\newline 用例: The unusual shape of the roof attracted many visitors.\newline 訳: 屋根の珍しい形が、多くの来場者を引きつけた。 \\
\addlinespace[\vocabrowspace]
\en{uniform}\newline{\ipa \ipastrut /j\ipastress{vkey}{ˈ}{u}{0.160750}{0.179500}nəf\ipastress{vhead}{ˌ}{ɔ}{0.125500}{0.701000}rm/} & \pos{形容詞・名詞}\newline 統一された・一様な／制服 & 全体の状態がそろっていること。決められた制服も指す。\newline 用例: The tiles have a uniform color across the entire wall.\newline 訳: 壁全体のタイルの色が均一になっている。\newline 用例: Each volunteer received a uniform before the museum opened.\newline 訳: 美術館が開く前に、各ボランティアは制服を受け取った。 \\
\addlinespace[\vocabrowspace]
\en{routine}\newline{\ipa \ipastrut /rut\ipastress{vkey}{ˈ}{i}{-0.013750}{0.392500}n/} & \pos{形容詞・名詞}\newline 日常的な・定型的な・通常の／決まった手順 & いつもの定型的なことや、繰り返して行う決まった手順。\newline 用例: Routine maintenance keeps the laboratory's ventilation system working reliably.\newline 訳: 定期的な保守によって、実験室の換気設備が安定して動く。\newline 用例: Our morning routine includes checking the plants and opening the windows.\newline 訳: 毎朝の決まった手順には、植物の確認と窓を開けることが含まれる。 \\
\addlinespace[\vocabrowspace]
\en{flaw}\newline{\ipa \ipastrut flaw /flɑ/  flawed /flɑd/} & \pos{名詞（形容詞形：flawed）}\newline 欠陥・欠点・不備 & 仕組みや製品の欠陥や不備。flawed は欠陥がある状態を表す。\newline 用例: The prototype revealed a flaw in the folding mechanism.\newline 訳: 試作品によって、折り畳み機構の欠陥が明らかになった。\newline 用例: The proposal relied on a flawed estimate of demand.\newline 訳: その提案は、問題のある需要予測に基づいていた。 \\
\addlinespace[\vocabrowspace]
\en{angle}\newline{\ipa \ipastrut /\ipastress{vkey}{ˈ}{æ}{0.341750}{0.192500}ŋɡəl/} & \pos{名詞}\newline 物事を見る観点・調べる切り口 & 物事を考える切り口や、二つの線などが作る角度。\newline 用例: The interview gave us a fresh angle on urban gardening.\newline 訳: その面談によって、都市の園芸について新しい観点が得られた。\newline 用例: Adjust the angle of the screen to reduce reflections.\newline 訳: 反射を減らすため、画面の角度を調整してください。 \\
\addlinespace[\vocabrowspace]
\en{regulation}\newline{\ipa \ipastrut /r\ipastress{vhead}{ˌ}{ɛ}{0.115750}{0.701500}ɡjəl\ipastress{vkey}{ˈ}{eɪ}{0.339250}{0.192500}ʃən/} & \pos{名詞}\newline 公的な規則・規制 & 公的に定められた規則や、活動を規則によって管理すること。\newline 用例: A new regulation requires clearer labels on recycled packaging.\newline 訳: 新しい規則は、再生包装材により分かりやすい表示を求めている。\newline 用例: The article examines regulation of commercial fishing in coastal waters.\newline 訳: その論文は、沿岸水域における商業漁業の規制を検討している。 \\
\addlinespace[\vocabrowspace]
\en{overhead}\newline{\ipa \ipastrut /\ipastress{vkey}{ˈ}{oʊ}{0.462750}{0.192500}vərh\ipastress{vhead}{ˌ}{ɛ}{0.115750}{0.701500}d/} & \pos{名詞}\newline 本来の作業に付随する時間・処理の負担 & 主な仕事を進めるために付随して生じる、時間や費用の負担。\newline 用例: Frequent status meetings added unnecessary overhead to the small project.\newline 訳: 頻繁な進捗会議が、小さなプロジェクトに余分な負担を加えた。\newline 用例: Sharing office space reduced our monthly overhead considerably last year.\newline 訳: 事務所を共用したことで、昨年は毎月の間接経費が大幅に減った。 \\
\addlinespace[\vocabrowspace]
\en{credit}\newline{\ipa \ipastrut /kr\ipastress{vkey}{ˈ}{ɛ}{0.115750}{0.193000}dɪt/} & \pos{名詞}\newline （この文脈では）店内で使える買い物残高 & 店などで、次の買い物の支払いに使える残高。\newline 用例: The store added credit to my account after the return.\newline 訳: 返品後、店は私のアカウントに買い物用の残高を追加した。 \\
\addlinespace[\vocabrowspace]
\en{mid-session}\newline{\ipa \ipastrut /mɪd s\ipastress{vkey}{ˈ}{ɛ}{0.115750}{0.193000}ʃən/} & \pos{形容詞}\newline セッション途中の・一連のやり取りの途中の & 会議や一続きの作業の、途中に当たること。\newline 用例: A mid-session break gave everyone time to stretch and relax.\newline 訳: 会議の途中の休憩で、全員が体を伸ばしてくつろげた。 \\
\addlinespace[\vocabrowspace]
\en{misinterpretation}\newline{\ipa \ipastrut /m\ipastress{vhead}{ˌ}{ɪ}{0.032500}{0.688000}sɪnt\ipastress{vhead}{ˌ}{ɜ}{0.115750}{0.701500}rprət\ipastress{vkey}{ˈ}{eɪ}{0.339250}{0.192500}ʃən/} & \pos{名詞}\newline 誤った解釈・意味の取り違え & 言葉や情報を、本来とは違う意味に受け取ること。\newline 用例: The ambiguous symbol led to a misinterpretation of the map.\newline 訳: 曖昧な記号のせいで、地図が誤って解釈された。 \\
\addlinespace[\vocabrowspace]
\en{refinement}\newline{\ipa \ipastrut /rɪf\ipastress{vkey}{ˈ}{aɪ}{0.340750}{0.192500}nmənt/} & \pos{名詞}\newline 改善・改良・より良くするための調整 & 細かな調整を加えて、内容や仕上がりを良くすること。\newline 用例: A small refinement made the handle more comfortable to hold.\newline 訳: 小さな改良によって、取っ手が握りやすくなった。 \\
\addlinespace[\vocabrowspace]
\en{longer}\newline{\ipa \ipastrut /l\ipastress{vkey}{ˈ}{ɔ}{0.125500}{0.192500}ŋɡər/} & \pos{副詞（比較級）}\newline より長く・さらに長い時間 & ある基準と比べて、さらに長い時間にわたること。\newline 用例: The team stayed longer to finish preparing the exhibition room.\newline 訳: チームは展示室の準備を終えるため、さらに長く残った。 \\
\addlinespace[\vocabrowspace]
\en{case}\newline{\ipa \ipastrut /keɪs/} & \pos{名詞}\newline （対応・調査の対象となる）案件 & 対応や調査の対象として、一つにまとめて扱う案件。\newline 用例: The inspector opened a new case after receiving the complaint.\newline 訳: 検査官は、苦情を受けて新たな案件の調査を始めた。 \\
\addlinespace[\vocabrowspace]
\en{merely}\newline{\ipa \ipastrut /m\ipastress{vkey}{ˈ}{ɪ}{0.032500}{0.179500}rli/} & \pos{副詞}\newline 単に・ただ…だけ & 述べていることが、それだけにとどまると示す言葉。\newline 用例: The sketch merely shows the general shape of the garden.\newline 訳: そのスケッチは、庭のおおまかな形を示しているだけだ。 \\
\addlinespace[\vocabrowspace]
\en{lookup}\newline{\ipa \ipastrut /l\ipastress{vkey}{ˈ}{ʊ}{0.156750}{0.179500}kʌp/} & \pos{名詞}\newline 情報の検索・照会 & 必要な情報を、表やデータベースなどから探し出すこと。\newline 用例: A quick database lookup confirmed the equipment's last inspection date.\newline 訳: データベースをすばやく照会し、機器の最終点検日を確認できた。 \\
\addlinespace[\vocabrowspace]
\en{handful}\newline{\ipa \ipastrut /h\ipastress{vkey}{ˈ}{æ}{0.341750}{0.192500}ndfʊl/} & \pos{名詞}\newline 一握り・ごく少数 & 数が多くなく、ごく少数に限られていること。\newline 用例: Only a handful of participants had used the software before.\newline 訳: 以前にそのソフトを使った参加者は、ごく少数だった。 \\
\addlinespace[\vocabrowspace]
\en{routinely}\newline{\ipa \ipastrut /rut\ipastress{vkey}{ˈ}{i}{-0.013750}{0.392500}nli/} & \pos{副詞}\newline 日常的に・いつものこととして & 特別なこととしてではなく、普段から繰り返して。\newline 用例: The staff routinely checks emergency exits before visitors arrive.\newline 訳: 職員は普段から、来場者が来る前に非常口を確認している。 \\
\addlinespace[\vocabrowspace]
\en{unanticipated}\newline{\ipa \ipastrut /\ipastress{vhead}{ˌ}{ʌ}{0.142750}{0.688000}nænt\ipastress{vkey}{ˈ}{ɪ}{0.032500}{0.179500}səpeɪtɪd/} & \pos{形容詞}\newline 予想していなかった・想定外の & 計画を立てた時点では、予想していなかったこと。\newline 用例: An unanticipated delay gave us time to inspect the surrounding area.\newline 訳: 想定外の遅延によって、周辺を調べる時間ができた。 \\
\addlinespace[\vocabrowspace]
\en{gracefully}\newline{\ipa \ipastrut /ɡr\ipastress{vkey}{ˈ}{eɪ}{0.339250}{0.192500}sfəli/} & \pos{副詞}\newline 円滑に・問題を大きくせず適切に & 問題が起きても、混乱を広げずに円滑に対処して。\newline 用例: The application handles missing images gracefully by displaying descriptive placeholders.\newline 訳: 画像が欠けている場合、アプリは説明付きの代替表示で円滑に対処する。 \\
\addlinespace[\vocabrowspace]
\en{leverage}\newline{\ipa \ipastrut /l\ipastress{vkey}{ˈ}{ɛ}{0.115750}{0.193000}vərɪdʒ/} & \pos{名詞}\newline 効果を大きくする力・影響力 & 手持ちの条件や資源によって、結果に働きかける力。\newline 用例: A larger membership gave the association greater leverage in negotiations.\newline 訳: 会員数の増加によって、協会の交渉力が高まった。 \\
\addlinespace[\vocabrowspace]
\en{dilution}\newline{\ipa \ipastrut /daɪl\ipastress{vkey}{ˈ}{u}{0.160750}{0.179500}ʃən/} & \pos{名詞}\newline 薄まること・弱まること & 混ぜたり広げたりすることで、濃さや効果が弱まること。\newline 用例: Adding unrelated topics caused a dilution of the workshop's central message.\newline 訳: 無関係な話題を加えたため、講習会の中心的な主張が薄まった。 \\
\addlinespace[\vocabrowspace]
\en{variant}\newline{\ipa \ipastrut /v\ipastress{vkey}{ˈ}{ɛ}{0.115750}{0.193000}riənt/} & \pos{名詞}\newline 別版・派生版・一部が異なるもの & 基本となるものから、一部を変えて作られた別の形。\newline 用例: We designed a compact variant for use in smaller classrooms.\newline 訳: 私たちは、狭い教室向けに小型の別版を設計した。 \\
\addlinespace[\vocabrowspace]
\en{invocable}\newline{\ipa \ipastrut /ɪnv\ipastress{vkey}{ˈ}{oʊ}{0.462750}{0.192500}kəbəl/} & \pos{形容詞}\newline 呼び出せる・呼び出して使える & プログラムの機能などを呼び出して、実行できる状態。\newline 用例: The export function is invocable from the application's main menu.\newline 訳: 書き出し機能は、アプリのメインメニューから呼び出せる。 \\
\addlinespace[\vocabrowspace]
\en{loyalty}\newline{\ipa \ipastrut /l\ipastress{vkey}{ˈ}{ɔɪ}{0.310000}{0.192500}əlti/} & \pos{名詞}\newline 愛着・ひいき・継続的な支持 & 人や店などに愛着を持ち、継続して支持すること。\newline 用例: Consistent service helped the neighborhood bakery earn lasting customer loyalty.\newline 訳: 安定した接客により、近所のパン屋は長く続く顧客の支持を得た。 \\
\addlinespace[\vocabrowspace]
\en{amount}\newline{\ipa \ipastrut /əm\ipastress{vkey}{ˈ}{aʊ}{0.465750}{0.192500}nt/} & \pos{名詞}\newline 量・金額 & 物の量や、支払うお金などの金額。\newline 用例: The form asks for the amount spent on travel.\newline 訳: その用紙には、旅費として使った金額を記入する。\newline 用例: Measure the amount of water before mixing the clay.\newline 訳: 粘土を混ぜる前に、水の量を測ってください。 \\
\addlinespace[\vocabrowspace]
\en{midway}\newline{\ipa \ipastrut /m\ipastress{vhead}{ˌ}{ɪ}{0.032500}{0.688000}dw\ipastress{vkey}{ˈ}{eɪ}{0.339250}{0.192500}/} & \pos{副詞}\newline 途中で・中ほどで & ある活動の開始から終了までの、中ほどの時点で。\newline 用例: The presenter paused midway to invite questions from the audience.\newline 訳: 発表者は途中で一度止まり、聴衆からの質問を募った。 \\
\addlinespace[\vocabrowspace]
\en{distinct}\newline{\ipa \ipastrut /dɪst\ipastress{vkey}{ˈ}{ɪ}{0.032500}{0.179500}ŋkt/} & \pos{形容詞}\newline 明確に別の・個別の & ほかのものと、はっきり分けて扱えること。\newline 用例: The exhibition explores three distinct themes through photographs and sound.\newline 訳: その展覧会は、写真と音を通じて三つの異なるテーマを探っている。 \\
\addlinespace[\vocabrowspace]
\en{upset}\newline{\ipa \ipastrut /ʌps\ipastress{vkey}{ˈ}{ɛ}{0.115750}{0.193000}t/} & \pos{形容詞}\newline 腹を立てた・動揺した & 出来事によって、腹を立てたり心が乱れたりした状態。\newline 用例: She felt upset when the recording of her talk disappeared.\newline 訳: 自分の講演の録画が消えてしまい、彼女は動揺した。 \\
\addlinespace[\vocabrowspace]
\en{sentiment}\newline{\ipa \ipastrut /s\ipastress{vkey}{ˈ}{ɛ}{0.115750}{0.193000}ntəmənt/} & \pos{名詞}\newline 感情・気持ち・意見 & 人や集団が物事に対して持つ、気持ちや意見。\newline 用例: The survey measured public sentiment about extending the library's opening hours.\newline 訳: 調査は、図書館の開館時間延長に対する市民の意向を測った。 \\
\addlinespace[\vocabrowspace]
\en{frustration}\newline{\ipa \ipastrut /frʌstr\ipastress{vkey}{ˈ}{eɪ}{0.339250}{0.192500}ʃən/} & \pos{名詞}\newline 思いどおりにならないことへの不満・いら立ち & 期待どおりに進まないことで生じる、不満やいら立ち。\newline 用例: Repeated delays caused frustration among volunteers preparing for the event.\newline 訳: 度重なる遅延で、催しを準備するボランティアにいら立ちが広がった。 \\
\addlinespace[\vocabrowspace]
\en{transactional}\newline{\ipa \ipastrut /trænz\ipastress{vkey}{ˈ}{æ}{0.341750}{0.192500}kʃənəl/} & \pos{形容詞}\newline 取引に関する & 売買や支払いなど、具体的な取引に関係すること。\newline 用例: The analyst grouped transactional records by payment date and currency.\newline 訳: 分析担当者は、取引記録を支払日と通貨で分類した。 \\
\addlinespace[\vocabrowspace]
\en{outside}\newline{\ipa \ipastrut /\ipastress{vhead}{ˌ}{aʊ}{0.465750}{0.701000}ts\ipastress{vkey}{ˈ}{aɪ}{0.340750}{0.192500}d/} & \pos{前置詞}\newline …の外に・…の範囲外に & 場所や、決められた範囲の外側にあること。\newline 用例: Questions outside the survey's scope were saved for later discussion.\newline 訳: 調査の範囲外の質問は、後の話し合いに回された。 \\
\addlinespace[\vocabrowspace]
\en{room}\newline{\ipa \ipastrut /rum/} & \pos{名詞（不可算）}\newline 余地・余裕 & 何かを加えたり行ったりするために残っている余裕。\newline 用例: The schedule leaves room for an additional visit to the workshop.\newline 訳: 日程には、工房をもう一度訪ねる余裕がある。 \\
\addlinespace[\vocabrowspace]
\en{thorough}\newline{\ipa \ipastrut /θ\ipastress{vkey}{ˈ}{ɜ}{0.115750}{0.193000}roʊ/} & \pos{形容詞}\newline 徹底的な・細部まで行き届いた & 対象を漏れなく、細かいところまで確かめること。\newline 用例: A thorough inspection revealed damage hidden beneath the wooden floor.\newline 訳: 徹底した点検で、木の床の下に隠れていた損傷が見つかった。 \\
\addlinespace[\vocabrowspace]
\en{degradation}\newline{\ipa \ipastrut /d\ipastress{vhead}{ˌ}{ɛ}{0.115750}{0.701500}ɡrəd\ipastress{vkey}{ˈ}{eɪ}{0.339250}{0.192500}ʃən/} & \pos{名詞}\newline 劣化・質の低下 & 物の状態や性能などが、以前より悪くなること。\newline 用例: Repeated washing caused noticeable degradation of the fabric's protective coating.\newline 訳: 洗濯を繰り返すと、生地の保護加工にはっきりした劣化が生じた。 \\
\addlinespace[\vocabrowspace]
\en{runtime}\newline{\ipa \ipastrut /r\ipastress{vkey}{ˈ}{ʌ}{0.142750}{0.179500}ntaɪm/} & \pos{名詞}\newline 実行時・動作時 & プログラムを作る段階ではなく、実際に動かしている時点。\newline 用例: At runtime, the application reads its display settings from a file.\newline 訳: 実行時に、アプリは表示設定をファイルから読み込む。 \\
\addlinespace[\vocabrowspace]
\en{loose}\newline{\ipa \ipastrut /lus/} & \pos{形容詞}\newline 緩い・厳密でない & 条件や結びつきが緩く、厳密には定まっていないこと。\newline 用例: The proposal offered only a loose definition of project success.\newline 訳: 提案書で示されたプロジェクトの成功の定義は、大まかなものにすぎなかった。 \\
\addlinespace[\vocabrowspace]
\en{off-script}\newline{\ipa \ipastrut /\ipastress{vhead}{ˌ}{ɑ}{0.146500}{0.701000}f skr\ipastress{vkey}{ˈ}{ɪ}{0.032500}{0.179500}pt/} & \pos{形容詞}\newline 予定・決められた手順から外れた & あらかじめ用意した筋書きや、進行予定から外れたこと。\newline 用例: An off-script question led to a lively discussion about local history.\newline 訳: 予定外の質問から、郷土史について活発な話し合いが始まった。 \\
\addlinespace[\vocabrowspace]
\en{contingency}\newline{\ipa \ipastrut /kənt\ipastress{vkey}{ˈ}{ɪ}{0.032500}{0.179500}ndʒənsi/} & \pos{名詞}\newline 起こり得る不測の事態 & 今後起こる可能性があり、備えが必要な不測の事態。\newline 用例: The outdoor event plan includes a budget for every major contingency.\newline 訳: 屋外イベントの計画には、主な不測の事態に備える予算がある。 \\
\addlinespace[\vocabrowspace]
\en{count}\newline{\ipa \ipastrut /kaʊnt/} & \pos{名詞}\newline 数・件数 & 数えて得られた個数や、対象となるものの件数。\newline 用例: The final count showed enough chairs for all registered guests.\newline 訳: 最終的に数えた結果、登録済みの参加者全員に椅子が足りると分かった。 \\
\addlinespace[\vocabrowspace]
\en{conspicuously}\newline{\ipa \ipastrut /kənsp\ipastress{vkey}{ˈ}{ɪ}{0.032500}{0.179500}kjuəsli/} & \pos{副詞}\newline 目立って・はっきり分かるほど & 周囲と比べて、目に留まるほどはっきりと。\newline 用例: The new sign was conspicuously larger than the surrounding notices.\newline 訳: 新しい看板は、周囲の掲示より目立って大きかった。 \\
\addlinespace[\vocabrowspace]
\en{gentle}\newline{\ipa \ipastrut /dʒ\ipastress{vkey}{ˈ}{ɛ}{0.115750}{0.193000}ntəl/} & \pos{形容詞}\newline 穏やかな・厳しくない & 人への接し方や表現が穏やかで、厳しくないこと。\newline 用例: A gentle reminder encouraged participants to return their borrowed materials.\newline 訳: 穏やかな呼びかけによって、参加者は借りた資料の返却を促された。 \\
\addlinespace[\vocabrowspace]
\en{harsh}\newline{\ipa \ipastrut /hɑrʃ/} & \pos{形容詞}\newline 厳しい・手厳しい & 言葉や評価などの当たりが強く、手厳しいこと。\newline 用例: The editor's harsh comments discouraged the writer from sharing another draft.\newline 訳: 編集者の手厳しい意見で、書き手は次の原稿を見せる気を失った。 \\
\addlinespace[\vocabrowspace]
\en{triviality}\newline{\ipa \ipastrut /tr\ipastress{vhead}{ˌ}{ɪ}{0.032500}{0.688000}vi\ipastress{vkey}{ˈ}{æ}{0.341750}{0.192500}ləti/} & \pos{名詞}\newline 些細なこと・重要でないこと & 大きな目的から見て、重要性が低い細かな事柄。\newline 用例: We spent too much meeting time discussing a triviality.\newline 訳: 私たちは会議の時間を、些細なことの議論に使いすぎた。 \\
\addlinespace[\vocabrowspace]
\en{fence}\newline{\ipa \ipastrut /fɛns/} & \pos{名詞}\newline 柵・囲い／コードを囲む区切り記号 & 物理的な囲いや、文章中で範囲を示す区切り記号。\newline 用例: A wooden fence separates the vegetable beds from the walkway.\newline 訳: 木の柵が、野菜畑と通路を隔てている。\newline 用例: The closing code fence is missing from this Markdown example.\newline 訳: このMarkdownの例には、コードを囲む終わりの区切り記号がない。 \\
\addlinespace[\vocabrowspace]
\en{commentary}\newline{\ipa \ipastrut /k\ipastress{vkey}{ˈ}{ɑ}{0.146500}{0.192500}mənt\ipastress{vhead}{ˌ}{ɛ}{0.115750}{0.701500}ri/} & \pos{名詞}\newline 説明・解説・論評 & 出来事や内容について添える、説明や解説。\newline 用例: The documentary includes commentary from the engineers who restored the bridge.\newline 訳: その記録映画には、橋を修復した技術者たちの解説が入っている。 \\
\addlinespace[\vocabrowspace]
\en{occasional}\newline{\ipa \ipastrut /ək\ipastress{vkey}{ˈ}{eɪ}{0.339250}{0.192500}ʒənəl/} & \pos{形容詞}\newline 時折の・たまに起こる & 毎回ではなく、たまに起きたり行ったりすること。\newline 用例: An occasional visit helps us maintain contact with local suppliers.\newline 訳: たまに訪問することで、地元の仕入れ先とのつながりを保てる。 \\
\addlinespace[\vocabrowspace]
\en{malformed}\newline{\ipa \ipastrut /m\ipastress{vhead}{ˌ}{æ}{0.341750}{0.701000}lf\ipastress{vkey}{ˈ}{ɔ}{0.125500}{0.192500}rmd/} & \pos{形容詞}\newline 正しい形になっていない・形式が不正な & 定められた形式に合わず、正しい形になっていないこと。\newline 用例: The importer rejected a malformed address in the contact list.\newline 訳: 取り込み処理は、連絡先一覧の不正な形式のアドレスを拒否した。 \\
\addlinespace[\vocabrowspace]
\en{absent}\newline{\ipa \ipastrut /\ipastress{vkey}{ˈ}{æ}{0.341750}{0.192500}bsənt/} & \pos{形容詞}\newline 存在しない・欠けている & 必要な場所に、あるはずの物や情報が存在しないこと。\newline 用例: The inspection date was absent from the equipment record.\newline 訳: 機器の記録には、点検日が記載されていなかった。 \\
\addlinespace[\vocabrowspace]
\en{legitimately}\newline{\ipa \ipastrut /lədʒ\ipastress{vkey}{ˈ}{ɪ}{0.032500}{0.179500}təmətli/} & \pos{副詞}\newline 正当に・認められた形で & 規則や正当な理由に基づいた、認められる形で。\newline 用例: Registered participants can legitimately access the recordings after the workshop.\newline 訳: 登録済みの参加者は、講習会後に正当な権限で録画を閲覧できる。 \\
\addlinespace[\vocabrowspace]
\en{themselves}\newline{\ipa \ipastrut /ðəms\ipastress{vkey}{ˈ}{ɛ}{0.115750}{0.193000}lvz/} & \pos{代名詞（強調用法）}\newline 彼ら自身・その人たち本人 & 動作や判断をする、その人たち本人を強調する言葉。\newline 用例: The volunteers themselves chose the topics for the training sessions.\newline 訳: 研修のテーマは、ボランティア本人たちが選んだ。 \\
\addlinespace[\vocabrowspace]
\en{uncertain}\newline{\ipa \ipastrut /ʌns\ipastress{vkey}{ˈ}{ɜ}{0.115750}{0.193000}rtən/} & \pos{形容詞}\newline 不確かな・確信がない & 事実や結果について、確かな判断がまだできないこと。\newline 用例: We remain uncertain about the cause of the unusual noise.\newline 訳: 珍しい物音の原因について、私たちはまだ確信が持てない。 \\
\addlinespace[\vocabrowspace]
\en{underlying}\newline{\ipa \ipastrut /\ipastress{vhead}{ˌ}{ʌ}{0.142750}{0.688000}ndərl\ipastress{vkey}{ˈ}{aɪ}{0.340750}{0.192500}ɪŋ/} & \pos{形容詞}\newline 背後にある・根本的な & 表面に見える現象の背後にある、根本の事柄。\newline 用例: The inspection revealed an underlying problem with the building's drainage.\newline 訳: 点検で、建物の排水に根本的な問題があると分かった。 \\
\addlinespace[\vocabrowspace]
\en{brittle}\newline{\ipa \ipastrut /br\ipastress{vkey}{ˈ}{ɪ}{0.032500}{0.179500}təl/} & \pos{形容詞}\newline もろい・想定外の変化に弱い & 少し条件が変わるだけで壊れやすい、もろい性質。\newline 用例: The brittle script failed whenever a folder name contained spaces.\newline 訳: その壊れやすいスクリプトは、フォルダー名に空白があると失敗した。 \\
\addlinespace[\vocabrowspace]
\en{orthogonal}\newline{\ipa \ipastrut /ɔrθ\ipastress{vkey}{ˈ}{ɑ}{0.146500}{0.192500}ɡənəl/} & \pos{形容詞}\newline 別の問題に属する・直接関係しない & 二つの課題や性質が、互いに独立した別の問題であること。\newline 用例: The question of visual style is orthogonal to data accuracy.\newline 訳: 見た目のスタイルの問題は、データの正確さとは別の問題だ。 \\
\addlinespace[\vocabrowspace]
\en{futile}\newline{\ipa \ipastrut /fj\ipastress{vkey}{ˈ}{u}{0.160750}{0.179500}təl/} & \pos{形容詞}\newline 効果がなく徒労に終わる・無益な & 努力しても、求める効果や結果につながらないこと。\newline 用例: Searching the empty storeroom again seemed futile after the careful inspection.\newline 訳: 入念に点検した後で空の倉庫を再び探しても、無駄に思えた。 \\
\addlinespace[\vocabrowspace]
\en{caveat}\newline{\ipa \ipastrut /k\ipastress{vkey}{ˈ}{æ}{0.341750}{0.192500}vi\ipastress{vhead}{ˌ}{æ}{0.341750}{0.701000}t/} & \pos{名詞}\newline 事前に留意すべき注意・制限事項 & 説明や提案を受け取る際に、あらかじめ注意するべき点。\newline 用例: The estimate includes a caveat about unpredictable changes in material prices.\newline 訳: 見積もりには、材料価格が予測できない形で変動するという注意書きがある。 \\
\addlinespace[\vocabrowspace]
\en{viable}\newline{\ipa \ipastrut /v\ipastress{vkey}{ˈ}{aɪ}{0.340750}{0.192500}əbəl/} & \pos{形容詞}\newline 実際に機能しうる・実行可能な & 考えた案が、現実に実行できて機能する見込みのあること。\newline 用例: Shared transport became a viable option after the timetable changed.\newline 訳: 時刻表の変更後、共同送迎が実行可能な選択肢になった。 \\
\addlinespace[\vocabrowspace]
\en{calibrated}\newline{\ipa \ipastrut /k\ipastress{vkey}{ˈ}{æ}{0.341750}{0.192500}ləbr\ipastress{vhead}{ˌ}{eɪ}{0.339250}{0.701000}tɪd/} & \pos{過去分詞（他動詞由来）}\newline 適切な基準に合わせて調整された & 信頼できる基準に合わせて、測定や判断を調整した状態。\newline 用例: The laboratory uses calibrated instruments to compare small temperature changes.\newline 訳: 研究室は、わずかな温度変化を比較するために校正済みの機器を使う。 \\
\addlinespace[\vocabrowspace]
\en{distractor}\newline{\ipa \ipastrut /dɪstr\ipastress{vkey}{ˈ}{æ}{0.341750}{0.192500}ktər/} & \pos{名詞}\newline 紛らわしい誤答の選択肢 & 正答に似せるなどして、迷わせるように設けた誤答選択肢。\newline 用例: The quiz writer revised a distractor that nobody ever selected.\newline 訳: クイズの作成者は、誰も選ばなかった誤答選択肢を修正した。 \\
\addlinespace[\vocabrowspace]
\en{deliberate}\newline{\ipa \ipastrut /dɪl\ipastress{vkey}{ˈ}{ɪ}{0.032500}{0.179500}bərət/} & \pos{形容詞}\newline 意図して行われた・計画的な & 偶然ではなく、考えたうえで意図して行われたこと。\newline 用例: Leaving the central space empty was a deliberate design decision.\newline 訳: 中央の空間を空けたのは、意図的な設計上の判断だった。 \\
\addlinespace[\vocabrowspace]
\en{opaque}\newline{\ipa \ipastrut /oʊp\ipastress{vkey}{ˈ}{eɪ}{0.339250}{0.192500}k/} & \pos{形容詞}\newline （表現などが）意味をつかみにくい・不明瞭な & 仕組みや表現が分かりにくく、意味をつかみにくいこと。\newline 用例: The opaque instructions left volunteers unsure about where to begin.\newline 訳: 分かりにくい説明のせいで、ボランティアは何から始めればよいか迷った。 \\
\addlinespace[\vocabrowspace]
\en{erratic}\newline{\ipa \ipastrut /ɪr\ipastress{vkey}{ˈ}{æ}{0.341750}{0.192500}tɪk/} & \pos{形容詞}\newline 動きが不規則な・一貫しない & 動作や変化に規則性がなく、安定していないこと。\newline 用例: An aging battery caused erratic behavior in the portable recorder.\newline 訳: 古くなった電池が、携帯録音機の不安定な動作を引き起こした。 \\
\addlinespace[\vocabrowspace]
\en{wholesale}\newline{\ipa \ipastrut /h\ipastress{vkey}{ˈ}{oʊ}{0.462750}{0.192500}ls\ipastress{vhead}{ˌ}{eɪ}{0.339250}{0.701000}l/} & \pos{副詞}\newline 区別せずひとまとめに・一括して & 個別に分けたり選んだりせず、まとめて一度に。\newline 用例: The editor rejected the proposed changes wholesale without discussing individual points.\newline 訳: 編集者は、個々の点を議論せずに変更案を丸ごと却下した。 \\
\addlinespace[\vocabrowspace]
\en{skew}\newline{\ipa \ipastrut /skju/} & \pos{名詞}\newline 一方向への偏り・ゆがみ & 分布や結果が、ある方向に偏っていること。\newline 用例: The survey showed a skew toward participants living near the station.\newline 訳: その調査は、駅の近くに住む参加者に偏っていた。 \\
\addlinespace[\vocabrowspace]
\en{inherently}\newline{\ipa \ipastrut /ɪnh\ipastress{vkey}{ˈ}{ɪ}{0.032500}{0.179500}rəntli/} & \pos{副詞}\newline 本質的に・性質上 & 外から加わる条件ではなく、そのものの性質として。\newline 用例: Outdoor measurements are inherently variable because environmental conditions keep changing.\newline 訳: 環境条件が変わり続けるため、屋外での測定は性質上ばらつく。 \\
\addlinespace[\vocabrowspace]
\en{counterproductive}\newline{\ipa \ipastrut /k\ipastress{vhead}{ˌ}{aʊ}{0.465750}{0.701000}ntərprəd\ipastress{vkey}{ˈ}{ʌ}{0.142750}{0.179500}ktɪv/} & \pos{形容詞}\newline かえって目的の達成を妨げる・逆効果の & 目的の達成を助けるつもりが、かえって妨げになること。\newline 用例: Adding more meetings became counterproductive when preparation time kept shrinking.\newline 訳: 準備時間が減り続ける中で、会議を増やすことは逆効果になった。 \\
\addlinespace[\vocabrowspace]
\en{meticulously}\newline{\ipa \ipastrut /mət\ipastress{vkey}{ˈ}{ɪ}{0.032500}{0.179500}kjələsli/} & \pos{副詞}\newline 細部まで注意を払って・入念に & 細かな点にも気を配り、入念に作業して。\newline 用例: She meticulously labeled every specimen before placing it in storage.\newline 訳: 彼女は標本を保管する前に、一つずつ入念にラベルを付けた。 \\
\addlinespace[\vocabrowspace]
\en{chronic}\newline{\ipa \ipastrut /kr\ipastress{vkey}{ˈ}{ɑ}{0.146500}{0.192500}nɪk/} & \pos{形容詞}\newline 長く続く・常態化した & 好ましくない状態が、長く続いて常態化していること。\newline 用例: Chronic understaffing made it difficult to keep the gallery open.\newline 訳: 慢性的な人手不足で、展示室を開け続けるのが難しかった。 \\
\addlinespace[\vocabrowspace]
\en{predisposed}\newline{\ipa \ipastrut /pr\ipastress{vhead}{ˌ}{i}{-0.013750}{0.901000}dɪsp\ipastress{vkey}{ˈ}{oʊ}{0.462750}{0.192500}zd/} & \pos{形容詞}\newline …しがちな傾向がある・…する方へ気持ちが傾いた & ある考えや行動に、もともと傾きやすい状態。\newline 用例: People are often predisposed to favor ideas they helped develop.\newline 訳: 人はしばしば、自分が考えるのを手伝った案を好みがちだ。 \\
\addlinespace[\vocabrowspace]
\en{outright}\newline{\ipa \ipastrut /\ipastress{vhead}{ˌ}{aʊ}{0.465750}{0.701000}tr\ipastress{vkey}{ˈ}{aɪ}{0.340750}{0.192500}t/} & \pos{副詞}\newline 完全に・全面的に & 部分的にではなく、完全にそうすること。\newline 用例: The committee rejected the proposal outright at its first meeting.\newline 訳: 委員会は最初の会議で、その提案を全面的に拒否した。 \\
\addlinespace[\vocabrowspace]
\en{durable}\newline{\ipa \ipastrut /d\ipastress{vkey}{ˈ}{ʊ}{0.156750}{0.179500}rəbəl/} & \pos{形容詞}\newline 長く保たれる・持続する & 時間がたっても、状態や効果が長く保たれること。\newline 用例: The archive needs durable labels that remain readable for decades.\newline 訳: 記録保管庫には、何十年も読める状態を保つ丈夫なラベルが必要だ。 \\
\addlinespace[\vocabrowspace]
\en{salient}\newline{\ipa \ipastrut /s\ipastress{vkey}{ˈ}{eɪ}{0.339250}{0.192500}liənt/} & \pos{形容詞}\newline とくに目につく・際立った & 多くの要素の中で、とくに目を引く際立ったこと。\newline 用例: The summary highlighted the most salient features of the new design.\newline 訳: 要約は、新しい設計のとくに際立った特徴を取り上げた。 \\
\addlinespace[\vocabrowspace]
\en{fidelity}\newline{\ipa \ipastrut /fɪd\ipastress{vkey}{ˈ}{ɛ}{0.115750}{0.193000}ləti/} & \pos{名詞}\newline 元の内容への忠実さ・再現の正確さ & 元の情報や音などを、どれだけ忠実に再現しているか。\newline 用例: The new recorder captures quiet voices with impressive fidelity.\newline 訳: 新しい録音機は、小さな声を驚くほど忠実に録音する。 \\
\addlinespace[\vocabrowspace]
\en{granularity}\newline{\ipa \ipastrut /ɡr\ipastress{vhead}{ˌ}{æ}{0.341750}{0.701000}njəl\ipastress{vkey}{ˈ}{æ}{0.341750}{0.192500}rəti/} & \pos{名詞}\newline 扱う単位の細かさ・粒度 & 情報や作業を、どのくらい細かい単位で扱うか。\newline 用例: The dashboard lets users change the granularity of time-based charts.\newline 訳: ダッシュボードでは、時間別のグラフの粒度を変更できる。 \\
\addlinespace[\vocabrowspace]
\en{empirical}\newline{\ipa \ipastrut /ɪmp\ipastress{vkey}{ˈ}{ɪ}{0.032500}{0.179500}rɪkəl/} & \pos{形容詞}\newline 観察や実測による・実証に基づく & 理屈だけでなく、観察や実験で得た事実に基づくこと。\newline 用例: The researchers gathered empirical evidence through repeated tests in different rooms.\newline 訳: 研究者たちは、異なる部屋で試験を繰り返し、実証的な根拠を集めた。 \\
\addlinespace[\vocabrowspace]
\en{propensity}\newline{\ipa \ipastrut /prəp\ipastress{vkey}{ˈ}{ɛ}{0.115750}{0.193000}nsəti/} & \pos{名詞}\newline ある行動をとりやすい傾向・性向 & 特定の行動を取ったり、反応を示したりしやすい傾向。\newline 用例: The study examined people's propensity to postpone unfamiliar tasks.\newline 訳: その研究は、慣れない作業を先延ばしにする人の傾向を調べた。 \\
\addlinespace[\vocabrowspace]
\en{nitpick}\newline{\ipa \ipastrut /n\ipastress{vkey}{ˈ}{ɪ}{0.032500}{0.179500}tp\ipastress{vhead}{ˌ}{ɪ}{0.032500}{0.688000}k/} & \pos{名詞}\newline 些細な点への細かい指摘・あら探し & 本筋にはあまり影響しない、些細な点への指摘。\newline 用例: His only nitpick concerned the spacing beneath the picture captions.\newline 訳: 彼の細かな指摘は、写真の説明文の下の間隔だけだった。 \\
\addlinespace[\vocabrowspace]
\en{redaction}\newline{\ipa \ipastrut /rɪd\ipastress{vkey}{ˈ}{æ}{0.341750}{0.192500}kʃən/} & \pos{名詞}\newline 公開できない箇所の削除・墨消し & 公開できない情報を、文書から削除したり伏せたりすること。\newline 用例: The document requires redaction of personal addresses before public release.\newline 訳: その文書は公開する前に、個人の住所を伏せる必要がある。 \\
\addlinespace[\vocabrowspace]
\en{off-the-shelf}\newline{\ipa \ipastrut /\ipastress{vhead}{ˌ}{ɑ}{0.146500}{0.701000}f ðə ʃ\ipastress{vkey}{ˈ}{ɛ}{0.115750}{0.193000}lf/} & \pos{形容詞}\newline 既製の・特注せずに利用できる & 特別に作ってもらわず、そのまま買って使える既製のこと。\newline 用例: We used off-the-shelf components to build the first weather station.\newline 訳: 最初の気象観測装置は、既製の部品を使って作った。 \\
\addlinespace[\vocabrowspace]
\en{opt-in}\newline{\ipa \ipastrut /\ipastress{vkey}{ˈ}{ɑ}{0.146500}{0.192500}ptɪn/} & \pos{名詞（修飾用法）}\newline 本人が選んで利用する方式の & 本人が利用すると選んだ場合に、初めて有効になる方式。\newline 用例: The newsletter uses an opt-in system for all new subscribers.\newline 訳: そのニュースレターは、新規読者すべてに希望登録方式を使っている。 \\
\addlinespace[\vocabrowspace]
\en{substantive}\newline{\ipa \ipastrut /s\ipastress{vkey}{ˈ}{ʌ}{0.142750}{0.179500}bstəntɪv/} & \pos{形容詞}\newline 内容の実質に関わる・本質的な & 見た目や表現だけでなく、内容の実質に関わること。\newline 用例: The revised proposal makes substantive changes to the research plan.\newline 訳: 改訂案では、研究計画の中身に実質的な変更を加えている。 \\
\addlinespace[\vocabrowspace]
\en{ephemeral}\newline{\ipa \ipastrut /ɪf\ipastress{vkey}{ˈ}{ɛ}{0.115750}{0.193000}mərəl/} & \pos{形容詞}\newline 一時的な・短期間しか存続しない & 一時的に存在し、短い期間で消える性質。\newline 用例: The installation uses ephemeral patterns created by moving light and water.\newline 訳: その作品は、動く光と水が作る一時的な模様を利用している。 \\
\addlinespace[\vocabrowspace]
\en{surgical}\newline{\ipa \ipastrut /s\ipastress{vkey}{ˈ}{ɜ}{0.115750}{0.193000}rdʒɪkəl/} & \pos{形容詞}\newline 必要な箇所に絞った精密な & 必要な箇所を正確に狙い、限定した範囲で行うこと。\newline 用例: A surgical edit removed the repeated paragraph without disturbing the argument.\newline 訳: 狙いを絞った修正で、論述を崩さずに重複した段落を削除した。 \\
\addlinespace[\vocabrowspace]
\en{fruitless}\newline{\ipa \ipastrut /fr\ipastress{vkey}{ˈ}{u}{0.160750}{0.179500}tləs/} & \pos{形容詞}\newline 成果が得られない・実を結ばない & 時間や労力を使っても、成果が得られないこと。\newline 用例: The search for spare parts proved fruitless after several visits.\newline 訳: 何軒か訪ねたが、予備部品の探索は実を結ばなかった。 \\
\addlinespace[\vocabrowspace]
\en{proxy}\newline{\ipa \ipastrut /pr\ipastress{vkey}{ˈ}{ɑ}{0.146500}{0.192500}ksi/} & \pos{名詞}\newline 代わりとして用いるもの・代用指標 & 直接測れないものなどを推し量るために使う、代わりの指標。\newline 用例: Attendance served as a proxy for interest in the new workshops.\newline 訳: 参加者数を、新しい講習会への関心を推し量る指標として使った。 \\
\addlinespace[\vocabrowspace]
\en{blanket}\newline{\ipa \ipastrut /bl\ipastress{vkey}{ˈ}{æ}{0.341750}{0.192500}ŋkɪt/} & \pos{形容詞}\newline 例外なく一律に適用される・全面的な & 個々の違いを問わず、全体に一律に適用されること。\newline 用例: The venue introduced a blanket ban on disposable plastic cups.\newline 訳: 会場は、使い捨てのプラスチック製コップを一律に禁止した。 \\
\addlinespace[\vocabrowspace]
\en{distilled}\newline{\ipa \ipastrut /dɪst\ipastress{vkey}{ˈ}{ɪ}{0.032500}{0.179500}ld/} & \pos{形容詞（過去分詞）}\newline 要点を抽出した・重要な部分に絞った & 多くの情報から、重要な点を選び取ってまとめた状態。\newline 用例: The handbook offers distilled advice from years of practical experience.\newline 訳: その手引きには、長年の実務経験から要点をまとめた助言が載っている。 \\
\addlinespace[\vocabrowspace]
\en{compounding}\newline{\ipa \ipastrut /kəmp\ipastress{vkey}{ˈ}{aʊ}{0.465750}{0.192500}ndɪŋ/} & \pos{形容詞（現在分詞）}\newline 積み重なることで増していく & 効果や影響が積み重なり、次第に大きくなっていくこと。\newline 用例: Small improvements produced compounding benefits throughout the following year.\newline 訳: 小さな改善が積み重なり、翌年を通して効果が増していった。 \\
\addlinespace[\vocabrowspace]
\en{prematurely}\newline{\ipa \ipastrut /pr\ipastress{vhead}{ˌ}{i}{-0.013750}{0.901000}mətʃ\ipastress{vkey}{ˈ}{ʊ}{0.156750}{0.179500}rli/} & \pos{副詞}\newline 早すぎる段階で・時期尚早に & 準備や条件が整う前の、早すぎる段階で。\newline 用例: The meeting ended prematurely when the building lost power.\newline 訳: 停電により、会議は予定より早く終わってしまった。 \\
\addlinespace[\vocabrowspace]
\en{unbiased}\newline{\ipa \ipastrut /ʌnb\ipastress{vkey}{ˈ}{aɪ}{0.340750}{0.192500}əst/} & \pos{形容詞}\newline 偏りのない・公平な & 特定の立場や選び方に偏らず、公平であること。\newline 用例: We invited an unbiased observer to assess the selection process.\newline 訳: 選考過程を評価するため、公平な立場の観察者を招いた。 \\
\addlinespace[\vocabrowspace]
\en{residual}\newline{\ipa \ipastrut /rɪz\ipastress{vkey}{ˈ}{ɪ}{0.032500}{0.179500}dʒuəl/} & \pos{形容詞}\newline 残った・残存する & 処理や変化が起きた後にも、なお残っていること。\newline 用例: The technician measured residual moisture after the drying cycle finished.\newline 訳: 技術者は、乾燥工程が終わった後の残留水分を測った。 \\
\addlinespace[\vocabrowspace]
\en{tolerable}\newline{\ipa \ipastrut /t\ipastress{vkey}{ˈ}{ɑ}{0.146500}{0.192500}lərəbəl/} & \pos{形容詞}\newline 許容できる・受け入れられる & 困った点があっても、受け入れられる範囲に収まること。\newline 用例: The background noise remained tolerable during the afternoon recording session.\newline 訳: 午後の録音中、背景の雑音は許容できる程度に収まっていた。 \\
\addlinespace[\vocabrowspace]
\en{verbatim}\newline{\ipa \ipastrut /vərb\ipastress{vkey}{ˈ}{eɪ}{0.339250}{0.192500}təm/} & \pos{副詞}\newline 一字一句そのままに・原文どおりに & 言葉を変えず、元の表現のまま正確に。\newline 用例: The secretary recorded the motion verbatim in the meeting minutes.\newline 訳: 書記は、議案を一字一句そのまま議事録に記録した。 \\
\addlinespace[\vocabrowspace]
\en{single source of truth}\newline{\ipa \ipastrut /s\ipastress{vkey}{ˈ}{ɪ}{0.032500}{0.179500}ŋɡəl sɔrs əv truθ/} & \pos{名詞句}\newline 皆が基準にする、一元化された情報の正本 & 皆が同じ基準として参照する、一元化された情報。\newline 用例: The shared calendar is our single source of truth for bookings.\newline 訳: 予約については、共有カレンダーを全員共通の正本としている。 \\
\addlinespace[\vocabrowspace]
\en{blind spot}\newline{\ipa \ipastrut /bl\ipastress{vkey}{ˈ}{aɪ}{0.340750}{0.192500}nd sp\ipastress{vhead}{ˌ}{ɑ}{0.146500}{0.701000}t/} & \pos{名詞句}\newline 監視・把握が届かない部分・盲点 & 十分に目が届かず、把握や検討から漏れている部分。\newline 用例: Nighttime transport became a blind spot in our initial visitor survey.\newline 訳: 最初の来場者調査では、夜間の交通が盲点になっていた。 \\
\addlinespace[\vocabrowspace]
\en{rule of thumb}\newline{\ipa \ipastrut /r\ipastress{vkey}{ˈ}{u}{0.160750}{0.179500}l əv θ\ipastress{vkey}{ˈ}{ʌ}{0.142750}{0.179500}m/} & \pos{名詞句}\newline 実務上の大まかな目安・経験則 & 厳密な規則ではなく、経験から使うおおまかな判断の目安。\newline 用例: Our rule of thumb is to leave extra time for setup.\newline 訳: 私たちは、設営時間に余裕を持たせることを経験則としている。 \\
\addlinespace[\vocabrowspace]
\en{stratified random sample}\newline{\ipa \ipastrut /str\ipastress{vkey}{ˈ}{æ}{0.341750}{0.192500}təf\ipastress{vhead}{ˌ}{aɪ}{0.340750}{0.701000}d r\ipastress{vkey}{ˈ}{æ}{0.341750}{0.192500}ndəm s\ipastress{vkey}{ˈ}{æ}{0.341750}{0.192500}mpəl/} & \pos{名詞句}\newline 種類ごとに集団を分け、各集団から無作為に選んだ標本 & 集団を種類別に分け、それぞれから無作為に選んだ標本。\newline 用例: Researchers surveyed a stratified random sample of households across several districts.\newline 訳: 研究者は複数地区の世帯を層化して無作為抽出し、調査した。 \\
\addlinespace[\vocabrowspace]
\en{outage}\newline{\ipa \ipastrut /\ipastress{vkey}{ˈ}{aʊ}{0.465750}{0.192500}tɪdʒ/} & \pos{名詞}\newline （サービスなどの）停止・不通 & 設備やサービスが止まり、一時的に使えなくなること。\newline 用例: A brief power outage interrupted the rehearsal at the community theater.\newline 訳: 短い停電で、市民劇場のリハーサルが中断した。 \\
\addlinespace[\vocabrowspace]
\en{opposite}\newline{\ipa \ipastrut /\ipastress{vkey}{ˈ}{ɑ}{0.146500}{0.192500}pəzɪt/} & \pos{名詞}\newline 正反対・逆 & 性質や方向などが、あるものと正反対であるもの。\newline 用例: The actual result was the opposite of what we expected.\newline 訳: 実際の結果は、予想したものとは正反対だった。 \\
\addlinespace[\vocabrowspace]
\en{guidance}\newline{\ipa \ipastrut /ɡ\ipastress{vkey}{ˈ}{aɪ}{0.340750}{0.192500}dəns/} & \pos{名詞}\newline 指針・手引き & どう進めたり判断したりすればよいかを示す手引き。\newline 用例: New volunteers receive guidance on welcoming visitors with different accessibility needs.\newline 訳: 新しいボランティアは、多様な配慮を必要とする来館者の迎え方を教わる。 \\
\addlinespace[\vocabrowspace]
\en{case-facts}\newline{\ipa \ipastrut /keɪs fækts/} & \pos{名詞句}\newline 案件に関する事実情報 & 一つの案件について、確認できた事実をまとめた情報。\newline 用例: The case-facts section lists inspection dates and equipment serial numbers.\newline 訳: 案件の事実情報の欄には、点検日と機器の製造番号が記されている。 \\
\addlinespace[\vocabrowspace]
\en{incrementally}\newline{\ipa \ipastrut /\ipastress{vhead}{ˌ}{ɪ}{0.032500}{0.688000}ŋkrəm\ipastress{vkey}{ˈ}{ɛ}{0.115750}{0.193000}ntəli/} & \pos{副詞}\newline 少しずつ・段階的に & 一度にすべてを変えず、小さな段階を重ねて。\newline 用例: The team improved the prototype incrementally after each round of testing.\newline 訳: チームは、試験を一巡するごとに試作品を少しずつ改良した。 \\
\addlinespace[\vocabrowspace]
\en{rigid}\newline{\ipa \ipastrut /r\ipastress{vkey}{ˈ}{ɪ}{0.032500}{0.179500}dʒɪd/} & \pos{形容詞}\newline 融通の利かない・厳格な & 決まりや考えが固く、状況に応じて変えにくいこと。\newline 用例: The rigid timetable left little room for unexpected questions.\newline 訳: 融通の利かない時間割では、想定外の質問を受ける余裕がほとんどなかった。 \\
\addlinespace[\vocabrowspace]
\en{step-by-step}\newline{\ipa \ipastrut /st\ipastress{vkey}{ˈ}{ɛ}{0.115750}{0.193000}p baɪ st\ipastress{vkey}{ˈ}{ɛ}{0.115750}{0.193000}p/} & \pos{形容詞}\newline 一つずつ順を追った・段階的な & 作業や説明を、一つずつ順番に進めること。\newline 用例: The booklet provides step-by-step instructions for assembling the display stand.\newline 訳: 小冊子には、展示台を組み立てる手順が順を追って書かれている。 \\
\addlinespace[\vocabrowspace]
\en{coherent}\newline{\ipa \ipastrut /koʊh\ipastress{vkey}{ˈ}{ɪ}{0.032500}{0.179500}rənt/} & \pos{形容詞}\newline 筋が通った・まとまりのある & 内容のつながりが自然で、全体として筋が通っていること。\newline 用例: Her notes developed into a coherent explanation of the design choices.\newline 訳: 彼女のメモは、設計上の選択を説明する筋の通った文章になった。 \\
\addlinespace[\vocabrowspace]
\en{shallow}\newline{\ipa \ipastrut /ʃ\ipastress{vkey}{ˈ}{æ}{0.341750}{0.192500}loʊ/} & \pos{形容詞}\newline 浅い・掘り下げが足りない & 説明や検討が表面にとどまり、十分に掘り下げられていないこと。\newline 用例: The article offered a shallow discussion of the neighborhood's transport problems.\newline 訳: その記事の地域交通問題に関する議論は、掘り下げが足りなかった。 \\
\addlinespace[\vocabrowspace]
\en{thin}\newline{\ipa \ipastrut /θɪn/} & \pos{形容詞}\newline 内容の薄い・中身の乏しい & 内容を支える説明や材料が少なく、乏しいこと。\newline 用例: The proposal was thin on details about maintaining the equipment.\newline 訳: その提案は、機器をどう保守するかの詳細が乏しかった。 \\
\addlinespace[\vocabrowspace]
\en{away}\newline{\ipa \ipastrut /əw\ipastress{vkey}{ˈ}{eɪ}{0.339250}{0.192500}/} & \pos{副詞}\newline 取り除かれて・なくなって & 動作によって、対象が取り除かれたり失われたりすること。\newline 用例: The editor cut away details that distracted from the main argument.\newline 訳: 編集者は、主張の中心から注意をそらす細部を削り取った。 \\
\addlinespace[\vocabrowspace]
\en{credible}\newline{\ipa \ipastrut /kr\ipastress{vkey}{ˈ}{ɛ}{0.115750}{0.193000}dəbəl/} & \pos{形容詞}\newline 信頼できる・信用に足る & 情報や説明に、信頼するに足る根拠があること。\newline 用例: The engineer offered a credible explanation for the unexpected vibration.\newline 訳: 技術者は、想定外の振動について信頼できる説明を示した。 \\
\addlinespace[\vocabrowspace]
\en{experimental}\newline{\ipa \ipastrut /ɪksp\ipastress{vhead}{ˌ}{ɛ}{0.115750}{0.701500}rəm\ipastress{vkey}{ˈ}{ɛ}{0.115750}{0.193000}ntəl/} & \pos{形容詞}\newline 実験的な・試験的な & 本格的な採用前に、可能性を試すためのこと。\newline 用例: An experimental display lets visitors explore the painting through sound.\newline 訳: 実験的な展示では、来館者が音を通して絵画を探ることができる。 \\
\addlinespace[\vocabrowspace]
\en{confusion}\newline{\ipa \ipastrut /kənfj\ipastress{vkey}{ˈ}{u}{0.160750}{0.179500}ʒən/} & \pos{名詞}\newline 混乱・戸惑い & 情報や状況を整理できず、理解が定まらない状態。\newline 用例: Similar room names caused confusion among visitors attending the conference.\newline 訳: 似た部屋の名前が、会議に来た参加者の混乱を招いた。 \\
\addlinespace[\vocabrowspace]
\en{bare}\newline{\ipa \ipastrut /bɛr/} & \pos{形容詞}\newline 何も付けていない・必要な指定のない & 付け加えられる情報や装飾などが、何もない状態。\newline 用例: The documentation shows a bare command without any optional arguments.\newline 訳: 文書には、省略可能な引数を付けていないコマンドだけが示されている。 \\
\addlinespace[\vocabrowspace]
\en{universal}\newline{\ipa \ipastrut /j\ipastress{vhead}{ˌ}{u}{0.160750}{0.688000}nəv\ipastress{vkey}{ˈ}{ɜ}{0.115750}{0.193000}rsəl/} & \pos{形容詞}\newline 全体に共通する・すべてに適用される & 特定の一部に限らず、全体に共通して当てはまること。\newline 用例: Clear labeling is a universal requirement across all our exhibitions.\newline 訳: 分かりやすい表示は、私たちのすべての展覧会に共通する要件だ。 \\
\addlinespace[\vocabrowspace]
\en{placement}\newline{\ipa \ipastrut /pl\ipastress{vkey}{ˈ}{eɪ}{0.339250}{0.192500}smənt/} & \pos{名詞}\newline 配置・置き場所 & 物や情報を、どこに置くかという配置。\newline 用例: The placement of the lamps changed how visitors viewed the sculpture.\newline 訳: 照明の配置によって、来館者から見た彫刻の印象が変わった。 \\
\addlinespace[\vocabrowspace]
\en{surprise}\newline{\ipa \ipastrut /sərpr\ipastress{vkey}{ˈ}{aɪ}{0.340750}{0.192500}z/} & \pos{名詞}\newline 予想外の出来事・想定外の事実 & 事前には考えていなかった、予想外の出来事や発見。\newline 用例: Finding an original sketch inside the cabinet was a pleasant surprise.\newline 訳: 戸棚の中で原画を見つけたのは、うれしい驚きだった。 \\
\addlinespace[\vocabrowspace]
\en{patience}\newline{\ipa \ipastrut /p\ipastress{vkey}{ˈ}{eɪ}{0.339250}{0.192500}ʃəns/} & \pos{名詞}\newline 忍耐・辛抱 & 急がず、難しい状況や待ち時間を辛抱すること。\newline 用例: Learning to restore delicate paper requires practice and considerable patience.\newline 訳: 傷みやすい紙の修復を学ぶには、練習と相当の忍耐が必要だ。 \\
\addlinespace[\vocabrowspace]
\en{conservative}\newline{\ipa \ipastrut /kəns\ipastress{vkey}{ˈ}{ɜ}{0.115750}{0.193000}rvətɪv/} & \pos{形容詞}\newline 慎重な・控えめな & 判断や見積もりで、無理をせず控えめにすること。\newline 用例: Our conservative estimate leaves extra time for unexpected delivery delays.\newline 訳: 控えめな見積もりにして、想定外の配送遅延に備える時間を残している。 \\
\addlinespace[\vocabrowspace]
\en{banner note}\newline{\ipa \ipastrut /b\ipastress{vkey}{ˈ}{æ}{0.341750}{0.192500}nər noʊt/} & \pos{名詞句}\newline 目立つ位置に表示する注意書き・お知らせ & 画面などの目立つ場所に掲げる、短い注意書きやお知らせ。\newline 用例: A banner note announces the temporary entrance on the museum website.\newline 訳: 美術館のウェブサイトでは、目立つお知らせで臨時の入り口を案内している。 \\
\addlinespace[\vocabrowspace]
\en{fraction}\newline{\ipa \ipastrut /fr\ipastress{vkey}{ˈ}{æ}{0.341750}{0.192500}kʃən/} & \pos{名詞}\newline ごく一部・わずかな割合 & 全体と比べて、ごく小さな部分やわずかな割合。\newline 用例: The prototype uses a fraction of the previous storage space.\newline 訳: 試作品が使う保存容量は、以前のごく一部にすぎない。 \\
\end{vtab}
\Needspace{14\baselineskip}

\section{4. 熟語をまとまりで読む}
\begin{vtab}
\en{on your behalf}\newline{\ipa \ipastrut /ɑn jʊr bɪh\ipastress{vkey}{ˈ}{æ}{0.341750}{0.192500}f/} & \pos{前置詞句}\newline あなたに代わって・あなたのために & あなたの代わりに、またはあなたの利益のために。\newline 用例: I can collect the exhibition tickets on your behalf tomorrow.\newline 訳: 明日、あなたに代わって展覧会のチケットを受け取れます。 \\
\addlinespace[\vocabrowspace]
\en{in terms of}\newline{\ipa \ipastrut /ɪn tɜrmz əv/} & \pos{前置詞（複数語）}\newline …の観点で・…について & 物事を説明したり比較したりするときの、着目する観点。\newline 用例: The new layout works better in terms of visitor flow.\newline 訳: 来館者の流れという点では、新しい配置のほうがうまく機能する。 \\
\addlinespace[\vocabrowspace]
\en{come down to}\newline{\ipa \ipastrut /kʌm daʊn tu/} & \pos{動詞句（自動詞＋前置詞）}\newline 結局は…にかかっている & 多くの条件の中で、最終的な決め手になること。\newline 用例: The decision will come down to cost and long-term maintenance.\newline 訳: その判断は結局、費用と長期的な保守にかかっている。 \\
\addlinespace[\vocabrowspace]
\en{up front}\newline{\ipa \ipastrut /ʌp frʌnt/} & \pos{副詞句（upfront は形容詞）}\newline 最初に・あらかじめ & 作業が進んでからではなく、あらかじめ最初の段階で。\newline 用例: We discussed the expected costs up front before ordering materials.\newline 訳: 資材を注文する前に、予想される費用をあらかじめ話し合った。\newline 用例: Upfront planning helped the volunteers prepare for the busy weekend.\newline 訳: 事前の計画が、ボランティアの忙しい週末への準備に役立った。 \\
\addlinespace[\vocabrowspace]
\en{vouch for}\newline{\ipa \ipastrut /vaʊtʃ fɔr/} & \pos{動詞句（自動詞＋前置詞）}\newline 確かだとして責任を持つ・保証する & 確かだと確認した相手や内容について、自分が責任を持って保証する。\newline 用例: I can vouch for the accuracy of these measurements.\newline 訳: この測定値の正確さは私が保証できます。 \\
\addlinespace[\vocabrowspace]
\en{not only ... but also ...}\newline{\ipa \ipastrut /nɑt \ipastress{vkey}{ˈ}{oʊ}{0.462750}{0.192500}nli ... bʌt \ipastress{vkey}{ˈ}{ɔ}{0.125500}{0.192500}lsoʊ .../} & \pos{並列・強調表現}\newline …だけでなく…も & 二つの長所や結果などを並べ、後の内容も加えて強調する。\newline 用例: The update not only saves space but also improves battery life.\newline 訳: この更新は容量を節約するだけでなく、電池の持ちも改善します。 \\
\addlinespace[\vocabrowspace]
\en{work through}\newline{\ipa \ipastrut /wɜrk θru/} & \pos{動詞句（他動詞）}\newline 順を追って取り組む・解決へ進める & 複雑な問題を一つずつ検討し、解決に向けて進める。\newline 用例: We worked through the scheduling problems during our lunch break.\newline 訳: 私たちは昼休みに予定調整の問題を一つずつ検討しました。 \\
\addlinespace[\vocabrowspace]
\en{build on}\newline{\ipa \ipastrut /bɪld ɑn/} & \pos{動詞句（自動詞＋前置詞）}\newline …を土台に発展させる & すでにある成果や知識を出発点にして、さらに発展させる。\newline 用例: The new course builds on students' knowledge of basic algebra.\newline 訳: この新しい講座は、学生の基礎的な代数学の知識を土台にしています。 \\
\addlinespace[\vocabrowspace]
\en{turn out to be}\newline{\ipa \ipastrut /tɜrn aʊt tə bi/} & \pos{動詞句（自動詞）}\newline 結果として…だと分かる & 調査や経験を経て、実際の性質や必要なことが明らかになる。\newline 用例: The damaged cable turned out to be the only problem.\newline 訳: 損傷したケーブルだけが問題だったと分かりました。\newline 用例: The old printer turned out to need a replacement motor.\newline 訳: その古いプリンターにはモーターの交換が必要だと分かりました。 \\
\addlinespace[\vocabrowspace]
\en{go off track}\newline{\ipa \ipastrut /ɡoʊ ɔf træk/} & \pos{動詞句（自動詞）}\newline 本来の方向からそれる & 話し合いや作業が、当初の目的や進める方向から外れる。\newline 用例: Our discussion went off track after someone mentioned office furniture.\newline 訳: 誰かがオフィスの家具に触れてから、議論が脱線しました。 \\
\addlinespace[\vocabrowspace]
\en{be meant to}\newline{\ipa \ipastrut /bi mɛnt tu/} & \pos{動詞句（他動詞の受け身）}\newline …することが想定されている & 物や仕組みに、あらかじめ想定されている用途や役割を示す。\newline 用例: This button is meant to pause the current recording.\newline 訳: このボタンは、現在の録音を一時停止するためのものです。 \\
\addlinespace[\vocabrowspace]
\en{make sense of}\newline{\ipa \ipastrut /meɪk sɛns əv/} & \pos{動詞句（makeは他動詞）}\newline …を理解する・意味をつかむ & 断片的な情報を結び付けて、全体として何を示すか理解する。\newline 用例: A simple chart helped us make sense of the survey results.\newline 訳: 簡単な図のおかげで、私たちは調査結果を理解できました。 \\
\addlinespace[\vocabrowspace]
\en{so far}\newline{\ipa \ipastrut /soʊ fɑr/} & \pos{副詞句}\newline これまでのところ & ある開始時点から、話している今までの経過や結果を示す。\newline 用例: We have received several useful suggestions from visitors so far.\newline 訳: これまでに、来場者から有益な提案をいくつかいただいています。 \\
\addlinespace[\vocabrowspace]
\en{due to}\newline{\ipa \ipastrut /du tu/} & \pos{前置詞（複数語）}\newline …のために・…が理由で & 出来事が起きた理由や、その状態を生んだ原因を示す。\newline 用例: The outdoor exhibition closed early due to strong winds.\newline 訳: 屋外展示は強風のため、予定より早く終了しました。 \\
\addlinespace[\vocabrowspace]
\en{leave ... up to ...}\newline{\ipa \ipastrut /liv ... ʌp tu .../} & \pos{動詞句（他動詞）}\newline …を…に任せる & 判断や作業を、自分以外の人や仕組みに委ねる。\newline 用例: We left the final color choice up to the illustrator.\newline 訳: 最終的な色の選択は、イラストレーターに任せました。 \\
\addlinespace[\vocabrowspace]
\en{depend on}\newline{\ipa \ipastrut /dɪp\ipastress{vkey}{ˈ}{ɛ}{0.115750}{0.193000}nd ɑn/} & \pos{動詞句（自動詞＋前置詞）}\newline …に依存する・…に左右される & 結果や行動が、別の条件や人の働きに左右される。\newline 用例: The delivery date depends on the availability of replacement parts.\newline 訳: 納品日は、交換部品を確保できるかどうかに左右されます。 \\
\addlinespace[\vocabrowspace]
\en{carry forward}\newline{\ipa \ipastrut /k\ipastress{vkey}{ˈ}{ɛ}{0.115750}{0.193000}ri f\ipastress{vkey}{ˈ}{ɔ}{0.125500}{0.192500}rwərd/} & \pos{動詞句（他動詞）}\newline …を引き継ぐ・…を後へ持ち越す & 今回得たものや残ったものを、次の機会にも引き継ぐ。\newline 用例: We will carry forward these lessons into our next exhibition.\newline 訳: 今回の教訓を、次の展示会にも引き継いでいきます。 \\
\addlinespace[\vocabrowspace]
\en{in-house}\newline{\ipa \ipastrut /\ipastress{vhead}{ˌ}{ɪ}{0.032500}{0.688000}nh\ipastress{vkey}{ˈ}{aʊ}{0.465750}{0.192500}s/} & \pos{副詞}\newline 社内で・自分たちの組織内で & 外部に依頼せず、自分たちの会社や組織の中で行う。\newline 用例: The museum designs and prints all its signs in-house.\newline 訳: その博物館は、すべての案内表示を組織内で制作・印刷しています。 \\
\addlinespace[\vocabrowspace]
\en{pull up}\newline{\ipa \ipastrut /pʊl ʌp/} & \pos{動詞句（他動詞）}\newline 情報を呼び出す・画面に表示する & 保存された資料や情報を探し出し、画面で見られるようにする。\newline 用例: Could you pull up the floor plan on your tablet?\newline 訳: タブレットで間取り図を表示してもらえますか。 \\
\addlinespace[\vocabrowspace]
\en{cut off}\newline{\ipa \ipastrut /kʌt ɔf/} & \pos{動詞句（他動詞の受け身）}\newline 途中で打ち切る・中断する & 通話や処理などが、終わりまで続かず途中で止められる。\newline 用例: Our video call was cut off during the storm.\newline 訳: 嵐の最中に、ビデオ通話が途中で切れてしまいました。 \\
\addlinespace[\vocabrowspace]
\en{lose track of}\newline{\ipa \ipastrut /luz træk əv/} & \pos{動詞句（loseは他動詞）}\newline …を見失う・…を把握できなくなる & 追い続けていた対象の状態や位置が、分からなくなる。\newline 用例: I lost track of the changes after several documents arrived.\newline 訳: 書類がいくつも届いた後、変更内容を把握できなくなりました。 \\
\addlinespace[\vocabrowspace]
\en{stem from}\newline{\ipa \ipastrut /stɛm frəm/} & \pos{動詞句（自動詞＋前置詞）}\newline …に起因する・…から生じる & ある問題や結果が、特定の原因から生じていることを示す。\newline 用例: Most of the confusion stems from an outdated address list.\newline 訳: 混乱の大半は、古い住所録に起因しています。 \\
\addlinespace[\vocabrowspace]
\en{route around}\newline{\ipa \ipastrut /rut ər\ipastress{vkey}{ˈ}{aʊ}{0.465750}{0.192500}nd/} & \pos{動詞句（自動詞的用法＋前置詞）}\newline 別の経路を使って…を迂回する & 通れない箇所や障害を避けて、別の経路を通る。\newline 用例: Network traffic can route around a failed connection automatically.\newline 訳: ネットワーク通信は、切れた接続を自動的に迂回できます。 \\
\addlinespace[\vocabrowspace]
\en{sit in}\newline{\ipa \ipastrut /sɪt ɪn/} & \pos{動詞句（自動詞＋前置詞）}\newline …の中・…の位置にある & 物や情報がある場所に存在すること。また、会合を傍聴すること。\newline 用例: The backup drive sits in a locked cabinet downstairs.\newline 訳: バックアップ用ドライブは、階下の鍵付き棚に置いてあります。\newline 用例: A trainee will sit in on tomorrow's planning meeting.\newline 訳: 研修生が明日の計画会議を傍聴します。 \\
\addlinespace[\vocabrowspace]
\en{drawn from}\newline{\ipa \ipastrut /drɑn frəm/} & \pos{過去分詞句（他動詞由来）}\newline …から取られた・…から得られた & 情報や着想が、特定の資料や経験をもとに得られている。\newline 用例: The illustrations were drawn from the artist's childhood memories.\newline 訳: その挿絵は、画家の子ども時代の記憶をもとにしています。 \\
\addlinespace[\vocabrowspace]
\en{carry out}\newline{\ipa \ipastrut /k\ipastress{vkey}{ˈ}{ɛ}{0.115750}{0.193000}ri aʊt/} & \pos{動詞句（他動詞）}\newline …を実行する・…を遂行する & 計画した作業や受けた指示を、実際の行動に移す。\newline 用例: Technicians carried out the safety checks before opening the laboratory.\newline 訳: 技術者たちは、研究室を開ける前に安全点検を実施しました。 \\
\addlinespace[\vocabrowspace]
\en{go through}\newline{\ipa \ipastrut /ɡoʊ θru/} & \pos{動詞句（自動詞）}\newline （手続き・取引などが）成立する・実行される & 申込みや取引などが必要な手続きを経て、実際に成立する。\newline 用例: The bank confirmed that our transfer had gone through.\newline 訳: 銀行は、私たちの送金が完了したことを確認しました。 \\
\addlinespace[\vocabrowspace]
\en{point at}\newline{\ipa \ipastrut /pɔɪnt æt/} & \pos{動詞句（自動詞＋前置詞）}\newline …を指す・…の位置を示す & 指や矢印などで、対象がある場所を示す。\newline 用例: The red arrow points at the emergency exit upstairs.\newline 訳: 赤い矢印は、上階の非常口を指しています。 \\
\addlinespace[\vocabrowspace]
\en{leave off}\newline{\ipa \ipastrut /liv ɑf/} & \pos{動詞句（自動詞）}\newline 中断する・いったん終わる & 続けていた作業を途中で止め、その地点を区切りとする。\newline 用例: The instructor resumed the lesson where we left off.\newline 訳: 講師は、前回中断したところから授業を再開しました。 \\
\addlinespace[\vocabrowspace]
\en{get + 過去分詞}\newline{\ipa \ipastrut /ɡɛt/} & \pos{動詞表現（受け身）}\newline …される & 人や物が、何らかの行為を受けることを表す。\newline 用例: Wet labels sometimes get damaged during the washing process.\newline 訳: ぬれたラベルは、洗浄中に傷むことがあります。 \\
\addlinespace[\vocabrowspace]
\en{tug-of-war}\newline{\ipa \ipastrut /t\ipastress{vhead}{ˌ}{ʌ}{0.142750}{0.688000}ɡ əv w\ipastress{vkey}{ˈ}{ɔ}{0.125500}{0.192500}r/} & \pos{名詞句}\newline 綱引き・互いに引っ張り合う状態 & 複数の立場や要求が、互いに異なる方向へ引っ張り合う状態。\newline 用例: The budget meeting became a tug-of-war between the two departments.\newline 訳: 予算会議は、二つの部門の綱引きになりました。 \\
\addlinespace[\vocabrowspace]
\en{take precedence over}\newline{\ipa \ipastrut /teɪk pr\ipastress{vkey}{ˈ}{ɛ}{0.115750}{0.193000}sədɛns \ipastress{vkey}{ˈ}{oʊ}{0.462750}{0.192500}vər/} & \pos{動詞句（takeは他動詞）}\newline …より優先される・…に優先する & 複数の条件が関係するとき、一方を他方より先に扱う。\newline 用例: Passenger safety takes precedence over maintaining the original schedule.\newline 訳: 乗客の安全は、当初の運行予定を守ることより優先されます。 \\
\addlinespace[\vocabrowspace]
\en{in principle}\newline{\ipa \ipastrut /ɪn pr\ipastress{vkey}{ˈ}{ɪ}{0.032500}{0.179500}nsəpəl/} & \pos{副詞句}\newline 原理的には・理屈の上では & 理論や基本的な考え方に照らせば成り立つ、という捉え方。\newline 用例: In principle, the battery could power the sensor for years.\newline 訳: 理屈の上では、その電池でセンサーを何年も動かせるはずです。 \\
\addlinespace[\vocabrowspace]
\en{in favor of}\newline{\ipa \ipastrut /ɪn f\ipastress{vkey}{ˈ}{eɪ}{0.339250}{0.192500}vər əv/} & \pos{前置詞句}\newline …のほうを選んで・…を採用するために & 候補を比較して、ある側を選んだり支持したりする。\newline 用例: The team rejected plastic packaging in favor of reusable containers.\newline 訳: チームはプラスチック包装をやめ、再利用できる容器を選びました。 \\
\addlinespace[\vocabrowspace]
\en{fall short}\newline{\ipa \ipastrut /fɑl ʃɔrt/} & \pos{動詞句（自動詞）}\newline 要求に届かない・不十分である & 求められる水準や目標に、実際の成果が届かない。\newline 用例: The final design fell short of our accessibility goals.\newline 訳: 最終案は、私たちのアクセシビリティ目標に届きませんでした。 \\
\addlinespace[\vocabrowspace]
\en{baked into}\newline{\ipa \ipastrut /beɪkt \ipastress{vkey}{ˈ}{ɪ}{0.032500}{0.179500}ntu/} & \pos{過去分詞句（他動詞由来）}\newline …に組み込まれた・…に織り込まれた & ある性質や条件が、仕組みの初めから組み込まれている。\newline 用例: Automatic backups are baked into the new editing software.\newline 訳: 新しい編集ソフトには、自動バックアップが初めから組み込まれています。 \\
\addlinespace[\vocabrowspace]
\en{spell out}\newline{\ipa \ipastrut /spɛl aʊt/} & \pos{動詞句（他動詞）}\newline 明確に詳しく説明する・明記する & 相手が分かるように、必要な条件や手順を具体的に示す。\newline 用例: The invitation spells out the dress code for guests.\newline 訳: 招待状には、来客が守る服装の決まりが明記されています。 \\
\addlinespace[\vocabrowspace]
\en{at the cost of}\newline{\ipa \ipastrut /æt ðə kɑst əv/} & \pos{前置詞句}\newline …という代償を払って・…を引き換えに & 何かを得るために、別の価値や利点を犠牲にする。\newline 用例: We increased storage capacity at the cost of portability.\newline 訳: 持ち運びやすさを犠牲にして、保存容量を増やしました。 \\
\addlinespace[\vocabrowspace]
\en{in the first place}\newline{\ipa \ipastrut /ɪn ðə fɜrst pleɪs/} & \pos{副詞句}\newline そもそも・初めから & 話の前提や、物事が始まった時点に立ち返って述べる。\newline 用例: The application never requested location access in the first place.\newline 訳: そもそも、そのアプリは位置情報へのアクセスを求めていませんでした。 \\
\addlinespace[\vocabrowspace]
\en{by design}\newline{\ipa \ipastrut /baɪ dɪz\ipastress{vkey}{ˈ}{aɪ}{0.340750}{0.192500}n/} & \pos{副詞句}\newline 設計上の意図によって・意図的に & 偶然の結果ではなく、初めから意図された仕組みによること。\newline 用例: The device shuts down before overheating, entirely by design.\newline 訳: その機器は、設計上の意図どおり過熱する前に停止します。 \\
\addlinespace[\vocabrowspace]
\en{on an ongoing basis}\newline{\ipa \ipastrut /ɑn ən \ipastress{vkey}{ˈ}{ɑ}{0.146500}{0.192500}nɡ\ipastress{vhead}{ˌ}{oʊ}{0.462750}{0.701000}ɪŋ b\ipastress{vkey}{ˈ}{eɪ}{0.339250}{0.192500}sɪs/} & \pos{副詞句}\newline 継続して・継続的に & 一度で終わらせず、一定の活動を継続して行う。\newline 用例: We update the training materials on an ongoing basis.\newline 訳: 研修資料は継続的に更新しています。 \\
\addlinespace[\vocabrowspace]
\en{slip through}\newline{\ipa \ipastrut /slɪp θru/} & \pos{動詞句（自動詞）}\newline 検査・防止策をすり抜ける & 確認や防止の仕組みに捕まらず、問題などが通ってしまう。\newline 用例: Occasional spelling mistakes still slip through despite careful editing.\newline 訳: 丁寧に編集しても、つづりの誤りが時々見逃されます。 \\
\addlinespace[\vocabrowspace]
\en{narrow down}\newline{\ipa \ipastrut /n\ipastress{vkey}{ˈ}{ɛ}{0.115750}{0.193000}roʊ daʊn/} & \pos{動詞句（他動詞）}\newline （候補・探索範囲を）絞り込む & 多くの候補から条件に合うものを選び、範囲を小さくする。\newline 用例: These interviews will help us narrow down the candidate list.\newline 訳: この面接で、候補者の一覧を絞り込めます。 \\
\addlinespace[\vocabrowspace]
\en{pay off}\newline{\ipa \ipastrut /peɪ ɑf/} & \pos{動詞句（自動詞）}\newline 成果につながる・かけた手間が報われる & 費やした努力や時間が、後に良い結果につながる。\newline 用例: Months of careful practice finally paid off at the concert.\newline 訳: 何か月にもわたる丁寧な練習が、演奏会でついに実を結びました。 \\
\addlinespace[\vocabrowspace]
\en{out of input order}\newline{\ipa \ipastrut /aʊt əv \ipastress{vkey}{ˈ}{ɪ}{0.032500}{0.179500}npʊt \ipastress{vkey}{ˈ}{ɔ}{0.125500}{0.192500}rdər/} & \pos{副詞句}\newline 入力された順序と異なる順番で & 処理へ渡した並びとは異なる順序で、結果などが現れる。\newline 用例: The worker threads may return results out of input order.\newline 訳: 作業スレッドは、入力順とは異なる順序で結果を返す場合があります。 \\
\addlinespace[\vocabrowspace]
\en{in turn}\newline{\ipa \ipastrut /ɪn tɜrn/} & \pos{副詞句}\newline さらにその次に・今度は & ある変化や行為を受けて、その次の段階へつなげる。\newline 用例: Better lighting improves visibility, which in turn reduces inspection time.\newline 訳: 照明の改善で見やすくなり、それが今度は点検時間の短縮につながります。 \\
\addlinespace[\vocabrowspace]
\en{from scratch}\newline{\ipa \ipastrut /frəm skrætʃ/} & \pos{副詞句}\newline 初めから・ゼロから & 既存の完成物を使わず、最初の段階から作り上げる。\newline 用例: She built the garden shed from scratch last summer.\newline 訳: 彼女は昨夏、庭の物置を一から作りました。 \\
\addlinespace[\vocabrowspace]
\en{on their own terms}\newline{\ipa \ipastrut /ɑn ðɛr oʊn tɜrmz/} & \pos{副詞句}\newline それ自体の基準・論理に照らして & 他の尺度に置き換えず、それぞれの狙いや前提に照らす。\newline 用例: We should assess these proposals on their own terms.\newline 訳: これらの提案は、それぞれの前提に照らして評価すべきです。 \\
\addlinespace[\vocabrowspace]
\en{on top of}\newline{\ipa \ipastrut /ɑn tɑp əv/} & \pos{前置詞句}\newline …にさらに重ねて・…に加えて & すでにある負担や内容へ、別のものがさらに加わる。\newline 用例: The venue charges cleaning fees on top of the rent.\newline 訳: 会場は、使用料に加えて清掃料金を請求します。 \\
\addlinespace[\vocabrowspace]
\en{paper over}\newline{\ipa \ipastrut /p\ipastress{vkey}{ˈ}{eɪ}{0.339250}{0.192500}pər \ipastress{vkey}{ˈ}{oʊ}{0.462750}{0.192500}vər/} & \pos{動詞句（他動詞）}\newline （問題を）取り繕う・表面上だけ覆い隠す & 根本の問題を解決せず、表面だけを取り繕って目立たなくする。\newline 用例: A cheerful announcement cannot paper over serious staffing problems.\newline 訳: 明るい告知を出しても、深刻な人員不足を取り繕うことはできません。 \\
\addlinespace[\vocabrowspace]
\en{have no bearing on}\newline{\ipa \ipastrut /hæv noʊ b\ipastress{vkey}{ˈ}{ɛ}{0.115750}{0.193000}rɪŋ ɑn/} & \pos{動詞句（haveは他動詞）}\newline …に影響しない・…と関係がない & ある事情が、対象となる判断や結果を左右しない。\newline 用例: Your membership length has no bearing on the selection process.\newline 訳: 会員である期間は、選考に影響しません。 \\
\addlinespace[\vocabrowspace]
\en{pin down}\newline{\ipa \ipastrut /pɪn daʊn/} & \pos{動詞句（他動詞）}\newline …を正確に特定する & 曖昧だった原因や条件を調べ、はっきり特定する。\newline 用例: Repeated measurements helped us pin down the source of vibration.\newline 訳: 測定を繰り返したことで、振動の発生源を正確に特定できました。 \\
\addlinespace[\vocabrowspace]
\en{rest on}\newline{\ipa \ipastrut /rɛst ɑn/} & \pos{動詞句（自動詞＋前置詞）}\newline …に基づく・…を土台とする & 判断や説明が、特定の根拠や前提を土台としている。\newline 用例: Their recommendation rests on several years of field research.\newline 訳: その提言は、数年にわたる現地調査に基づいています。 \\
\addlinespace[\vocabrowspace]
\en{drill down}\newline{\ipa \ipastrut /drɪl daʊn/} & \pos{動詞句（自動詞）}\newline より細かい内訳へ掘り下げる & 全体の数値や情報から、さらに細かな内訳へ進んで調べる。\newline 用例: Readers can drill down into individual months on the dashboard.\newline 訳: 利用者は、ダッシュボードで月ごとの詳細まで掘り下げられます。 \\
\addlinespace[\vocabrowspace]
\en{dressed up as}\newline{\ipa \ipastrut /drɛst ʌp æz/} & \pos{過去分詞句（他動詞由来）}\newline …であるかのように装われた & 実際の性質とは異なる、もっともらしい外見や名目を与えられる。\newline 用例: The advertisement was dressed up as an independent product review.\newline 訳: その広告は、独立した製品評価の記事のように装われていました。 \\
\addlinespace[\vocabrowspace]
\en{on demand}\newline{\ipa \ipastrut /ɑn dɪm\ipastress{vkey}{ˈ}{æ}{0.341750}{0.192500}nd/} & \pos{副詞句}\newline 必要になったときに・求めに応じて & 必要な時点や依頼があった時点で、利用できること。\newline 用例: The library can print replacement membership cards on demand.\newline 訳: 図書館は、必要に応じて会員証を再発行できます。 \\
\addlinespace[\vocabrowspace]
\en{have nothing to do with}\newline{\ipa \ipastrut /hæv n\ipastress{vkey}{ˈ}{ʌ}{0.142750}{0.179500}θɪŋ tə du wɪð/} & \pos{動詞句（haveは他動詞）}\newline …とはまったく関係がない & 二つの事柄の間に、関わりや結び付きがまったくない。\newline 用例: These delays have nothing to do with the software update.\newline 訳: 今回の遅れは、ソフトウェア更新とはまったく関係ありません。 \\
\addlinespace[\vocabrowspace]
\en{break down}\newline{\ipa \ipastrut /breɪk daʊn/} & \pos{動詞句（自動詞）}\newline 成り立たなくなる・うまく機能しなくなる & 仕組みややり取りが、うまく機能しなくなる。\newline 用例: Communication broke down when the teams changed their reporting schedules.\newline 訳: 両チームが報告予定を変更してから、連絡がうまくいかなくなりました。 \\
\addlinespace[\vocabrowspace]
\en{in practice}\newline{\ipa \ipastrut /ɪn pr\ipastress{vkey}{ˈ}{æ}{0.341750}{0.192500}ktɪs/} & \pos{副詞句}\newline 実際には・実際に運用すると & 理論上の説明に対して、実際に使ったときの状態を述べる。\newline 用例: The idea sounds simple but requires careful preparation in practice.\newline 訳: 考え方は簡単そうですが、実際には入念な準備が必要です。 \\
\addlinespace[\vocabrowspace]
\en{tell apart}\newline{\ipa \ipastrut /tɛl əp\ipastress{vkey}{ˈ}{ɑ}{0.146500}{0.192500}rt/} & \pos{動詞句（他動詞）}\newline …を見分ける・…を区別する & 似ているものの違いを認識し、それぞれを区別する。\newline 用例: The two cables are difficult to tell apart in darkness.\newline 訳: 暗い場所では、その二本のケーブルを見分けるのは困難です。 \\
\addlinespace[\vocabrowspace]
\en{account for}\newline{\ipa \ipastrut /ək\ipastress{vkey}{ˈ}{aʊ}{0.465750}{0.192500}nt fɔr/} & \pos{動詞句（自動詞＋前置詞）}\newline …を考慮に入れる・…を踏まえる & 計画や計算をするとき、影響する事情を織り込む。\newline 用例: Our travel estimate accounts for slower traffic near the stadium.\newline 訳: 移動時間の見積もりには、競技場付近の渋滞を織り込んでいます。 \\
\addlinespace[\vocabrowspace]
\en{stand or fall on}\newline{\ipa \ipastrut /stænd ɔr fɔl ɑn/} & \pos{動詞句（自動詞＋前置詞）}\newline 妥当性・成否が…にかかっている & ある根拠や条件が、案の成否や妥当性を決める。\newline 用例: These arguments stand or fall on the reliability of the evidence.\newline 訳: これらの主張の妥当性は、証拠の信頼性にかかっています。 \\
\addlinespace[\vocabrowspace]
\en{by construction}\newline{\ipa \ipastrut /baɪ kənstr\ipastress{vkey}{ˈ}{ʌ}{0.142750}{0.179500}kʃən/} & \pos{副詞句}\newline 仕組みそのものによって・構造上 & 作り方や構成そのものから、一定の性質が成り立つ。\newline 用例: The generated list contains only positive integers by construction.\newline 訳: その一覧には、生成の仕組み上、正の整数だけが含まれます。 \\
\addlinespace[\vocabrowspace]
\en{at scale}\newline{\ipa \ipastrut /æt skeɪl/} & \pos{副詞句}\newline 大規模に・本格的な規模で & 少量の試行を越え、多数の対象を扱う規模で行う。\newline 用例: This method becomes expensive when we apply it at scale.\newline 訳: この方法は、大規模に適用すると費用がかかります。 \\
\addlinespace[\vocabrowspace]
\en{out of the box}\newline{\ipa \ipastrut /aʊt əv ðə bɑks/} & \pos{副詞句}\newline 初期状態のままで・特別な調整をせずに & 導入した状態のまま、追加の調整をせずに使える。\newline 用例: The keyboard works with our tablets out of the box.\newline 訳: そのキーボードは、追加設定なしで私たちのタブレットに使えます。 \\
\addlinespace[\vocabrowspace]
\en{crowd out}\newline{\ipa \ipastrut /kraʊd aʊt/} & \pos{動詞句（他動詞）}\newline 圧迫して押しのける・締め出す & 限られた場所や時間を占め、ほかのものを押しのける。\newline 用例: Frequent status meetings crowd out time for focused design work.\newline 訳: 頻繁な進捗会議が、設計に集中する時間を圧迫しています。 \\
\end{vtab}
\Needspace{14\baselineskip}

\section{5. エージェントの設計と実行}
\begin{vtab}
\en{delegation}\newline{\ipa \ipastrut /d\ipastress{vhead}{ˌ}{ɛ}{0.115750}{0.701500}ləɡ\ipastress{vkey}{ˈ}{eɪ}{0.339250}{0.192500}ʃən/} & \pos{名詞}\newline 作業を他の担当へ任せること & 仕事を別の担当へ任せること。人と自動化の役割を決める判断も含む。\newline 用例: Effective delegation gave the editor more time for planning.\newline 訳: 適切に仕事を任せることで、編集者は計画に時間を使えるようになりました。\newline 用例: Thoughtful delegation starts with deciding which tasks need human judgment.\newline 訳: 適切な仕事の分担は、どの作業に人の判断が必要かを決めることから始まります。 \\
\addlinespace[\vocabrowspace]
\en{synthesis}\newline{\ipa \ipastrut /s\ipastress{vkey}{ˈ}{ɪ}{0.032500}{0.179500}nθəsɪs/} & \pos{名詞}\newline 複数の結果を組み合わせてまとめること & 個別に得た情報を結び付け、一つのまとまった理解にする。\newline 用例: Her synthesis of the interviews revealed several shared concerns.\newline 訳: 彼女が面談内容をまとめたところ、共通する懸念がいくつか見えてきました。 \\
\addlinespace[\vocabrowspace]
\en{context isolation}\newline{\ipa \ipastrut /k\ipastress{vkey}{ˈ}{ɑ}{0.146500}{0.192500}ntɛkst \ipastress{vhead}{ˌ}{aɪ}{0.340750}{0.701000}səl\ipastress{vkey}{ˈ}{eɪ}{0.339250}{0.192500}ʃən/} & \pos{名詞句}\newline 各エージェントの会話情報が分かれていること & それぞれの処理が参照する会話や情報の範囲を分けること。\newline 用例: Context isolation keeps each writing assistant's draft notes separate.\newline 訳: コンテキストの分離により、各執筆支援エージェントの下書きメモを分けて保持できます。 \\
\addlinespace[\vocabrowspace]
\en{handoff}\newline{\ipa \ipastrut /h\ipastress{vkey}{ˈ}{æ}{0.341750}{0.192500}nd\ipastress{vhead}{ˌ}{ɔ}{0.125500}{0.701000}f/} & \pos{名詞}\newline 別の担当への引き継ぎ & 担当が交代するときに、作業と必要な情報を引き渡すこと。\newline 用例: A short checklist made the evening shift handoff smoother.\newline 訳: 短いチェックリストで、夜勤への引き継ぎが円滑になりました。 \\
\addlinespace[\vocabrowspace]
\en{parallel}\newline{\ipa \ipastrut /p\ipastress{vkey}{ˈ}{æ}{0.341750}{0.192500}rəl\ipastress{vhead}{ˌ}{ɛ}{0.115750}{0.701500}l/} & \pos{形容詞}\newline 並行して進む & 複数の作業や処理が、同じ期間に並んで進む。\newline 用例: Parallel experiments allowed the laboratory to compare several materials quickly.\newline 訳: 並行実験により、研究室は複数の材料を素早く比較できました。 \\
\addlinespace[\vocabrowspace]
\en{sequential}\newline{\ipa \ipastrut /sɪkw\ipastress{vkey}{ˈ}{ɛ}{0.115750}{0.193000}nʃəl/} & \pos{形容詞}\newline 順番に進む & 一つが終わってから次へ移る順序で、処理が進む。\newline 用例: The tutorial presents sequential steps for assembling the wooden frame.\newline 訳: その教材は、木枠を組み立てる手順を順番に示しています。 \\
\addlinespace[\vocabrowspace]
\en{deterministic}\newline{\ipa \ipastrut /dɪt\ipastress{vhead}{ˌ}{ɜ}{0.115750}{0.701500}rmən\ipastress{vkey}{ˈ}{ɪ}{0.032500}{0.179500}stɪk/} & \pos{形容詞}\newline 定めた規則に従って確定的に動く & 同じ条件を与えれば、決まった規則に従って同じ結果になる。\newline 用例: A deterministic simulation produces identical results from identical starting conditions.\newline 訳: 確定的なシミュレーションは、同じ開始条件から同じ結果を出します。 \\
\addlinespace[\vocabrowspace]
\en{probabilistic}\newline{\ipa \ipastrut /pr\ipastress{vhead}{ˌ}{ɑ}{0.146500}{0.701000}bəbəl\ipastress{vkey}{ˈ}{ɪ}{0.032500}{0.179500}stɪk/} & \pos{形容詞}\newline 確率的な・必ずそうなるとは限らない & 結果を一つに固定せず、起こりやすさを確率で扱う。\newline 用例: The researchers used a probabilistic model to estimate rainfall.\newline 訳: 研究者たちは、降雨量の推定に確率モデルを使いました。 \\
\addlinespace[\vocabrowspace]
\en{termination}\newline{\ipa \ipastrut /t\ipastress{vhead}{ˌ}{ɜ}{0.115750}{0.701500}rmən\ipastress{vkey}{ˈ}{eɪ}{0.339250}{0.192500}ʃən/} & \pos{名詞}\newline 処理を終了すること & 実行中の処理を、終了条件や指示に従って終わらせること。\newline 用例: Unexpected termination left the temporary folder on the disk.\newline 訳: 予期しない終了により、一時フォルダーがディスクに残りました。 \\
\addlinespace[\vocabrowspace]
\en{workflows}\newline{\ipa \ipastrut /w\ipastress{vkey}{ˈ}{ɜ}{0.115750}{0.193000}rkfl\ipastress{vhead}{ˌ}{oʊ}{0.462750}{0.701000}z/} & \pos{名詞（複数形）}\newline あらかじめ決めた手順で進む一連の処理 & 目的を達成するために、決められた順序で進める一連の作業。\newline 用例: The studio documented its workflows for recording and editing podcasts.\newline 訳: そのスタジオは、ポッドキャストの録音と編集の作業手順を文書化しました。 \\
\addlinespace[\vocabrowspace]
\en{routing}\newline{\ipa \ipastrut /r\ipastress{vkey}{ˈ}{u}{0.160750}{0.179500}tɪŋ/} & \pos{名詞（動名詞）}\newline 条件や種類に応じた処理先への振り分け & 入力の内容や条件を見て、適した処理先へ送り分けること。\newline 用例: Automated routing sends incoming requests to the appropriate regional office.\newline 訳: 自動振り分けにより、届いた依頼が適切な地域事務所へ送られます。 \\
\addlinespace[\vocabrowspace]
\en{chaining}\newline{\ipa \ipastrut /tʃ\ipastress{vkey}{ˈ}{eɪ}{0.339250}{0.192500}nɪŋ/} & \pos{名詞（動名詞）}\newline 処理を順につなぎ、結果を次へ渡すこと & 前の処理で得た結果を、次の処理の入力としてつなぐこと。\newline 用例: Chaining these image filters creates a distinctive visual effect.\newline 訳: これらの画像フィルターをつなぐと、独特の視覚効果が得られます。 \\
\addlinespace[\vocabrowspace]
\en{parallelization}\newline{\ipa \ipastrut /p\ipastress{vhead}{ˌ}{æ}{0.341750}{0.701000}rələlaɪz\ipastress{vkey}{ˈ}{eɪ}{0.339250}{0.192500}ʃən/} & \pos{名詞}\newline 独立した作業を並行して実行する構成にすること & 同時に進められる作業を分け、並行して処理できるようにすること。\newline 用例: Parallelization reduced the time needed to render the animation.\newline 訳: 並列化により、そのアニメーションの描画に必要な時間が短くなりました。 \\
\addlinespace[\vocabrowspace]
\en{environment inspection}\newline{\ipa \ipastrut /ɪnv\ipastress{vkey}{ˈ}{aɪ}{0.340750}{0.192500}rənmənt ɪnsp\ipastress{vkey}{ˈ}{ɛ}{0.115750}{0.193000}kʃən/} & \pos{名詞句}\newline 作業環境や操作の結果を観察・確認すること & 作業場所や実行環境の状態を、操作前後などに確かめること。\newline 用例: Environment inspection revealed that the recording software lacked microphone access.\newline 訳: 環境の確認で、録音ソフトにマイクへのアクセス権がないと分かりました。 \\
\addlinespace[\vocabrowspace]
\en{evaluator-optimizer pattern}\newline{\ipa \ipastrut /ɪv\ipastress{vkey}{ˈ}{æ}{0.341750}{0.192500}ljueɪtər \ipastress{vkey}{ˈ}{ɑ}{0.146500}{0.192500}ptəmaɪzər p\ipastress{vkey}{ˈ}{æ}{0.341750}{0.192500}tərn/} & \pos{名詞句}\newline 評価結果を生成・修正側へ返して改善を繰り返す構成 & 生成物への評価を修正に戻し、改善を繰り返す処理の構成。\newline 用例: The evaluator-optimizer pattern guided repeated improvements to the practice worksheet.\newline 訳: 評価と改善を繰り返す構成により、練習用ワークシートを繰り返し改良しました。 \\
\addlinespace[\vocabrowspace]
\en{task following}\newline{\ipa \ipastrut /tæsk f\ipastress{vkey}{ˈ}{ɑ}{0.146500}{0.192500}loʊɪŋ/} & \pos{名詞句}\newline 求められた課題に沿って答えること & 依頼で求められた作業内容や条件に沿って、応答を作ること。\newline 用例: Our rubric measures task following separately from writing style.\newline 訳: 私たちの評価基準では、課題への適合を文体とは別に測定します。 \\
\addlinespace[\vocabrowspace]
\en{autonomous}\newline{\ipa \ipastrut /ɑt\ipastress{vkey}{ˈ}{ɑ}{0.146500}{0.192500}nəməs/} & \pos{形容詞}\newline 逐一の指示を待たず、自ら判断して進める & 個々の操作について指示を待たず、判断しながら行動する。\newline 用例: The autonomous cart delivers supplies between rooms without constant supervision.\newline 訳: その自律走行カートは、常時監視なしで部屋の間を移動して物品を届けます。 \\
\addlinespace[\vocabrowspace]
\en{orchestration}\newline{\ipa \ipastrut /\ipastress{vhead}{ˌ}{ɔ}{0.125500}{0.701000}rkəstr\ipastress{vkey}{ˈ}{eɪ}{0.339250}{0.192500}ʃən/} & \pos{名詞}\newline 複数の処理や担当の動きを調整し、全体を進めること & 複数の担当や処理の順序・連携を調整し、全体を動かすこと。\newline 用例: Careful orchestration kept the exhibition setup running on schedule.\newline 訳: 綿密な調整により、展示会の設営は予定どおりに進みました。 \\
\addlinespace[\vocabrowspace]
\en{invoke}\newline{\ipa \ipastrut invoke /ɪnv\ipastress{vkey}{ˈ}{oʊ}{0.462750}{0.192500}k/  invocation /\ipastress{vhead}{ˌ}{ɪ}{0.032500}{0.688000}nvək\ipastress{vkey}{ˈ}{eɪ}{0.339250}{0.192500}ʃən/} & \pos{他動詞（名詞形：invocation）}\newline 機能やエージェントを呼び出して動かす & 必要な処理や機能を、呼び出して実行させる。\newline 用例: The shortcut invokes the dictionary whenever you select unfamiliar text.\newline 訳: そのショートカットは、知らない語句を選択すると辞書を呼び出します。\newline 用例: Each invocation of the converter creates a new output folder.\newline 訳: 変換機能を一回呼び出すごとに、新しい出力フォルダーが作られます。 \\
\addlinespace[\vocabrowspace]
\en{spawn}\newline{\ipa \ipastrut /spɔn/} & \pos{他動詞}\newline 新しい処理やエージェントを生成して起動する & 新しい処理を作り、独立した実行単位として動かし始める。\newline 用例: The renderer spawns additional workers when larger images arrive.\newline 訳: 描画プログラムは、大きな画像が届くと追加の作業プロセスを起動します。 \\
\addlinespace[\vocabrowspace]
\en{adaptive}\newline{\ipa \ipastrut /əd\ipastress{vkey}{ˈ}{æ}{0.341750}{0.192500}ptɪv/} & \pos{形容詞}\newline 途中で分かった状況に合わせて変化する & 状況の変化や新たな情報に応じて、動き方を調整する。\newline 用例: Adaptive lighting adjusts its brightness as natural daylight changes.\newline 訳: 適応型照明は、自然光の変化に合わせて明るさを調整します。 \\
\addlinespace[\vocabrowspace]
\en{open-ended}\newline{\ipa \ipastrut /\ipastress{vhead}{ˌ}{oʊ}{0.462750}{0.701000}pən\ipastress{vkey}{ˈ}{ɛ}{0.115750}{0.193000}ndɪd/} & \pos{形容詞}\newline 進め方や答えの範囲が最初から固定されていない & 答えや進め方を一つに絞らず、探索する余地がある。\newline 用例: The workshop ends with an open-ended discussion about future exhibitions.\newline 訳: その講習会は、今後の展示についての自由な議論で締めくくられます。 \\
\addlinespace[\vocabrowspace]
\en{predefined}\newline{\ipa \ipastrut /pr\ipastress{vhead}{ˌ}{i}{-0.013750}{0.901000}dɪf\ipastress{vkey}{ˈ}{aɪ}{0.340750}{0.192500}nd/} & \pos{形容詞}\newline 実行前にあらかじめ定められた & 処理を始める前に、使う値や条件がすでに決められている。\newline 用例: Visitors choose a badge color from a predefined set.\newline 訳: 来場者は、あらかじめ定められた候補から名札の色を選びます。 \\
\addlinespace[\vocabrowspace]
\en{synchronous}\newline{\ipa \ipastrut /s\ipastress{vkey}{ˈ}{ɪ}{0.032500}{0.179500}ŋkrənəs/} & \pos{形容詞}\newline 呼び出した処理の完了を待ってから先へ進む同期式の & 依頼した処理の結果が返るまで待ってから、次へ進む。\newline 用例: The synchronous download prevents the script from continuing immediately.\newline 訳: 同期式のダウンロードでは、スクリプトはすぐには先へ進めません。 \\
\addlinespace[\vocabrowspace]
\en{asynchronous}\newline{\ipa \ipastrut /eɪs\ipastress{vkey}{ˈ}{ɪ}{0.032500}{0.179500}ŋkrənəs/} & \pos{形容詞}\newline 呼び出した処理の完了をその場で待たずに進められる非同期式の & 依頼した処理を待つ間にも、別の処理を進められる。\newline 用例: Asynchronous uploads let users keep editing while files transfer.\newline 訳: 非同期アップロードにより、ファイル転送中も編集を続けられます。 \\
\addlinespace[\vocabrowspace]
\en{blocking}\newline{\ipa \ipastrut /bl\ipastress{vkey}{ˈ}{ɑ}{0.146500}{0.192500}kɪŋ/} & \pos{形容詞}\newline 完了するまで後続の処理や利用者を待たせる & 実行している間、呼び出し元などが先に進めなくなる。\newline 用例: A blocking operation kept the interface frozen for several seconds.\newline 訳: 待機を伴う処理により、画面操作が数秒間できなくなりました。 \\
\addlinespace[\vocabrowspace]
\en{bottlenecks}\newline{\ipa \ipastrut /b\ipastress{vkey}{ˈ}{ɑ}{0.146500}{0.192500}təln\ipastress{vhead}{ˌ}{ɛ}{0.115750}{0.701500}ks/} & \pos{名詞（複数形）}\newline 全体の進みを遅らせる処理上の詰まり & 作業の流れの中で、全体の処理速度を制限する箇所。\newline 用例: The measurements exposed bottlenecks in our image processing pipeline.\newline 訳: 測定により、画像処理の流れで進行を遅らせる箇所が明らかになりました。 \\
\addlinespace[\vocabrowspace]
\en{dependencies}\newline{\ipa \ipastrut /dɪp\ipastress{vkey}{ˈ}{ɛ}{0.115750}{0.193000}ndənsiz/} & \pos{名詞（複数形）}\newline ほかの部品や処理を必要とする関係・依存関係 & 実行や完成のために、別の部品や作業を必要とする関係。\newline 用例: The diagram shows dependencies between the project's build tasks.\newline 訳: その図は、プロジェクトのビルド作業間の依存関係を示しています。 \\
\addlinespace[\vocabrowspace]
\en{upstream / downstream}\newline{\ipa \ipastrut /\ipastress{vhead}{ˌ}{ʌ}{0.142750}{0.688000}pstr\ipastress{vkey}{ˈ}{i}{-0.013750}{0.392500}m/ /d\ipastress{vhead}{ˌ}{aʊ}{0.465750}{0.701000}nstr\ipastress{vkey}{ˈ}{i}{-0.013750}{0.392500}m/} & \pos{形容詞}\newline 処理の前段／後段 & 処理の流れで、前にある段階と後にある段階を示す。\newline 用例: An upstream filter removes empty rows before any calculations.\newline 訳: 前段のフィルターが、計算前に空の行を取り除きます。\newline 用例: Downstream components receive the cleaned records from this converter.\newline 訳: 後段の部品は、この変換器から整えられた記録を受け取ります。 \\
\addlinespace[\vocabrowspace]
\en{concurrently}\newline{\ipa \ipastrut /kənk\ipastress{vkey}{ˈ}{ɝ}{0.264750}{0.193000}əntli/} & \pos{副詞}\newline 同時に・並行して & 二つ以上の処理や活動が、同じ時間帯に進む。\newline 用例: The program downloads audio and generates thumbnails concurrently.\newline 訳: そのプログラムは、音声のダウンロードと縮小画像の生成を並行して行います。 \\
\end{vtab}
\Needspace{14\baselineskip}

\section{6. ツール・MCP・権限}
\begin{vtab}
\en{resources}\newline{\ipa \ipastrut /r\ipastress{vkey}{ˈ}{i}{-0.013750}{0.392500}sɔrsɪz/} & \pos{名詞（複数形）}\newline 読み取り用に公開されるデータ & アプリなどが読み取れる形で提供される、文書やデータ。\newline 用例: The server exposes reference tables as read-only resources.\newline 訳: そのサーバーは、参照用の表を読み取り専用のリソースとして公開します。 \\
\addlinespace[\vocabrowspace]
\en{transient}\newline{\ipa \ipastrut /tr\ipastress{vkey}{ˈ}{æ}{0.341750}{0.192500}nʃənt/} & \pos{形容詞}\newline 一時的な & 長く続く状態ではなく、一時的に発生している。\newline 用例: The monitoring chart showed a transient increase in temperature.\newline 訳: 監視グラフには、一時的な温度上昇が記録されていました。 \\
\addlinespace[\vocabrowspace]
\en{recoverable}\newline{\ipa \ipastrut /rɪk\ipastress{vkey}{ˈ}{ʌ}{0.142750}{0.179500}vərəbəl/} & \pos{形容詞}\newline 回復・復旧が可能な & 問題が起きても、正常な状態へ戻せる可能性がある。\newline 用例: The damaged document was recoverable from yesterday's backup copy.\newline 訳: 破損した文書は、昨日のバックアップから復旧できました。 \\
\addlinespace[\vocabrowspace]
\en{business errors}\newline{\ipa \ipastrut /b\ipastress{vkey}{ˈ}{ɪ}{0.032500}{0.179500}znəs \ipastress{vkey}{ˈ}{ɛ}{0.115750}{0.193000}rərz/} & \pos{名詞句}\newline 業務上の規則に反するために生じるエラー & 入力などが業務上の条件や取り決めを満たさず、処理できないこと。\newline 用例: The booking service reports business errors when seats are unavailable.\newline 訳: 座席を確保できないとき、予約サービスは業務上のエラーを返します。 \\
\addlinespace[\vocabrowspace]
\en{backoff}\newline{\ipa \ipastrut /b\ipastress{vkey}{ˈ}{æ}{0.341750}{0.192500}k\ipastress{vhead}{ˌ}{ɔ}{0.125500}{0.701000}f/} & \pos{名詞}\newline 再試行の前に待ち時間を設けること & 失敗後すぐに繰り返さず、次の試行まで間隔を置くこと。\newline 用例: A short backoff gives the overloaded service time to recover.\newline 訳: 少し待ってから再試行することで、過負荷のサービスに回復する時間を与えます。 \\
\addlinespace[\vocabrowspace]
\en{permission rules}\newline{\ipa \ipastrut /pərm\ipastress{vkey}{ˈ}{ɪ}{0.032500}{0.179500}ʃən rulz/} & \pos{名詞句}\newline 操作を許可・確認・拒否するための規則 & 誰がどの操作をできるか、確認が必要かなどを定める条件。\newline 用例: These permission rules limit editing access to the current organizers.\newline 訳: これらの権限規則は、編集できる人を現在の運営担当者に限定します。 \\
\addlinespace[\vocabrowspace]
\en{transport agnostic}\newline{\ipa \ipastrut /tr\ipastress{vkey}{ˈ}{æ}{0.341750}{0.192500}nspɔrt æɡn\ipastress{vkey}{ˈ}{ɑ}{0.146500}{0.192500}stɪk/} & \pos{形容詞句}\newline 特定の通信手段だけに依存しない & データを運ぶ通信方式が変わっても、特定の方式に縛られない。\newline 用例: The protocol remains transport agnostic across different deployment environments.\newline 訳: そのプロトコルは、異なる導入環境でも特定の通信方式に依存しません。 \\
\addlinespace[\vocabrowspace]
\en{standard input/output}\newline{\ipa \ipastrut /st\ipastress{vkey}{ˈ}{æ}{0.341750}{0.192500}ndərd \ipastress{vkey}{ˈ}{ɪ}{0.032500}{0.179500}np\ipastress{vhead}{ˌ}{ʊ}{0.156750}{0.688000}t \ipastress{vkey}{ˈ}{aʊ}{0.465750}{0.192500}tp\ipastress{vhead}{ˌ}{ʊ}{0.156750}{0.688000}t/} & \pos{名詞句}\newline 標準入力・標準出力 & プログラムが通常使う、入力を受け取る口と結果を出す口。\newline 用例: The utility exchanges text through standard input/output during automated processing.\newline 訳: そのユーティリティーは、自動処理中に標準入力・標準出力でテキストをやり取りします。 \\
\addlinespace[\vocabrowspace]
\en{server primitives}\newline{\ipa \ipastrut /s\ipastress{vkey}{ˈ}{ɜ}{0.115750}{0.193000}rvər pr\ipastress{vkey}{ˈ}{ɪ}{0.032500}{0.179500}mətɪvz/} & \pos{名詞句}\newline server が提供する基本的な機能の種類 & サーバーの提供機能を構成する、基本単位となる機能の種類。\newline 用例: The diagram groups server primitives according to their intended roles.\newline 訳: その図は、サーバーの基本機能を役割別にまとめています。 \\
\addlinespace[\vocabrowspace]
\en{Model-Controlled}\newline{\ipa \ipastrut /m\ipastress{vkey}{ˈ}{ɑ}{0.146500}{0.192500}dəl kəntr\ipastress{vkey}{ˈ}{oʊ}{0.462750}{0.192500}ld/} & \pos{形容詞}\newline 利用の判断をモデルが行う & いつ機能を利用するかという判断を、モデルが担う。\newline 用例: This model-controlled operation begins when the assistant identifies a need.\newline 訳: このモデル判断の操作は、アシスタントが必要性を認めると始まります。 \\
\addlinespace[\vocabrowspace]
\en{App-Controlled}\newline{\ipa \ipastrut /æp kəntr\ipastress{vkey}{ˈ}{oʊ}{0.462750}{0.192500}ld/} & \pos{形容詞}\newline 利用の判断をアプリケーションが行う & データなどをいつ使うかという判断を、アプリケーションが担う。\newline 用例: An app-controlled refresh loads reference data when the screen opens.\newline 訳: アプリが制御する更新処理で、画面を開くと参照データを読み込みます。 \\
\addlinespace[\vocabrowspace]
\en{User-Controlled}\newline{\ipa \ipastrut /j\ipastress{vkey}{ˈ}{u}{0.160750}{0.179500}zər kəntr\ipastress{vkey}{ˈ}{oʊ}{0.462750}{0.192500}ld/} & \pos{形容詞}\newline 利用の開始をユーザーが選ぶ & 動作を始めるかどうかを、利用者自身が選ぶ。\newline 用例: The user-controlled export starts only after someone presses the button.\newline 訳: 利用者が制御する書き出しは、ボタンが押されて初めて始まります。 \\
\addlinespace[\vocabrowspace]
\en{variable interpolation}\newline{\ipa \ipastrut /v\ipastress{vkey}{ˈ}{ɛ}{0.115750}{0.193000}riəbəl ɪnt\ipastress{vhead}{ˌ}{ɜ}{0.115750}{0.701500}rpəl\ipastress{vkey}{ˈ}{eɪ}{0.339250}{0.192500}ʃən/} & \pos{名詞句}\newline ひな型の変数へ値を埋め込むこと & 文章や設定のひな型にある変数を、実際の値に置き換えること。\newline 用例: Variable interpolation inserts the selected date into the invitation template.\newline 訳: 変数への値の埋め込みにより、招待状のひな型に選択した日付が入ります。 \\
\addlinespace[\vocabrowspace]
\en{keyword arguments}\newline{\ipa \ipastrut /k\ipastress{vkey}{ˈ}{i}{-0.013750}{0.392500}w\ipastress{vhead}{ˌ}{ɜ}{0.115750}{0.701500}rd \ipastress{vkey}{ˈ}{ɑ}{0.146500}{0.192500}rɡjəmənts/} & \pos{名詞句}\newline 名前を指定して渡す引数 & 値だけを並べず、引数の名前も指定して関数へ渡す値。\newline 用例: Keyword arguments make this function's optional settings easier to read.\newline 訳: キーワード引数を使うと、この関数の省略可能な設定が読みやすくなります。 \\
\addlinespace[\vocabrowspace]
\en{decorators}\newline{\ipa \ipastrut /d\ipastress{vkey}{ˈ}{ɛ}{0.115750}{0.193000}kər\ipastress{vhead}{ˌ}{eɪ}{0.339250}{0.701000}tərz/} & \pos{名詞（複数形）}\newline 関数などに指定を付け、振る舞いを加える仕組み & 関数などに付けて、追加の振る舞いや指定を与える仕組み。\newline 用例: These decorators record execution time without changing the function body.\newline 訳: これらのデコレーターは、関数本体を変えずに実行時間を記録します。 \\
\addlinespace[\vocabrowspace]
\en{request-response pattern}\newline{\ipa \ipastrut /rɪkw\ipastress{vkey}{ˈ}{ɛ}{0.115750}{0.193000}st rɪsp\ipastress{vkey}{ˈ}{ɑ}{0.146500}{0.192500}ns p\ipastress{vkey}{ˈ}{æ}{0.341750}{0.192500}tərn/} & \pos{名詞句}\newline 要求を送り、その応答を受け取る処理の形 & 一方が処理を要求し、相手が対応する結果を返すやり取り。\newline 用例: The weather client uses a request-response pattern to obtain forecasts.\newline 訳: その天気クライアントは、要求と応答の形式で予報を取得します。 \\
\addlinespace[\vocabrowspace]
\en{intercept}\newline{\ipa \ipastrut /\ipastress{vhead}{ˌ}{ɪ}{0.032500}{0.688000}ntərs\ipastress{vkey}{ˈ}{ɛ}{0.115750}{0.193000}pt/} & \pos{他動詞}\newline 処理が実行される前に途中で捉え、介入する & 流れている要求や処理を途中で捉え、介入する。\newline 用例: The extension intercepts outgoing links and checks their destinations.\newline 訳: その拡張機能は、外部へのリンクを途中で捉え、移動先を確認します。 \\
\addlinespace[\vocabrowspace]
\en{normalize}\newline{\ipa \ipastrut /n\ipastress{vkey}{ˈ}{ɔ}{0.125500}{0.192500}rməl\ipastress{vhead}{ˌ}{aɪ}{0.340750}{0.701000}z/} & \pos{他動詞}\newline 異なる表し方を共通の形式へそろえる & ばらつきのある表記や値を、共通の基準に合わせる。\newline 用例: The importer normalizes date formats before saving the records.\newline 訳: 取り込み機能は、記録の保存前に日付の書式を統一します。 \\
\addlinespace[\vocabrowspace]
\en{heterogeneous}\newline{\ipa \ipastrut /h\ipastress{vhead}{ˌ}{ɛ}{0.115750}{0.701500}təroʊdʒ\ipastress{vkey}{ˈ}{i}{-0.013750}{0.392500}niəs/} & \pos{形容詞}\newline 異なる種類や形式が混在する & 一つの集合の中に、性質や形式の違うものが含まれる。\newline 用例: The archive contains a heterogeneous collection of maps and recordings.\newline 訳: その資料庫には、地図や録音など異なる種類の資料が収められています。 \\
\addlinespace[\vocabrowspace]
\en{prerequisite}\newline{\ipa \ipastrut /pr\ipastress{vhead}{ˌ}{i}{-0.013750}{0.901000}r\ipastress{vkey}{ˈ}{ɛ}{0.115750}{0.193000}kwəzɪt/} & \pos{名詞}\newline 先に満たしておく必要がある条件 & 次の段階へ進む前に、満たしておく必要がある条件。\newline 用例: Basic swimming ability is a prerequisite for this sailing course.\newline 訳: このヨット講習には、基本的な泳力が前提条件として必要です。 \\
\addlinespace[\vocabrowspace]
\en{scoped}\newline{\ipa \ipastrut /skoʊpt/} & \pos{形容詞}\newline 対象や使える範囲を限定した & 権限や作業が及ぶ範囲を限定し、対象を明確にする。\newline 用例: The scoped token permits access to a single project folder.\newline 訳: 範囲を限定したトークンは、一つのプロジェクトフォルダーへのアクセスを許可します。\newline 用例: A well-scoped assignment helped the intern finish the prototype confidently.\newline 訳: 範囲が明確で適切な課題により、実習生は自信を持って試作品を完成できました。 \\
\addlinespace[\vocabrowspace]
\en{least privilege}\newline{\ipa \ipastrut /list pr\ipastress{vkey}{ˈ}{ɪ}{0.032500}{0.179500}vəlɪdʒ/} & \pos{名詞句}\newline 必要最小限の権限だけを与える考え方 & 役割を果たすために必要な範囲だけ、権限を与える考え方。\newline 用例: Least privilege limits the damage from a compromised account.\newline 訳: 最小権限の考え方は、不正利用されたアカウントによる被害を限定します。 \\
\addlinespace[\vocabrowspace]
\en{credentials}\newline{\ipa \ipastrut /krɪd\ipastress{vkey}{ˈ}{ɛ}{0.115750}{0.193000}nʃəlz/} & \pos{名詞（複数形）}\newline 接続する主体を確かめるための認証情報 & サービスに接続するときに、利用者などの身元を確かめる情報。\newline 用例: Store the credentials in the organization's approved password manager.\newline 訳: 認証情報は、組織が承認したパスワード管理ツールに保存してください。 \\
\addlinespace[\vocabrowspace]
\en{wrap}\newline{\ipa \ipastrut /ræp/} & \pos{他動詞}\newline 既存機能を包み、別の窓口から使えるようにする & 既存機能の外側に処理や窓口を設け、使い方をまとめる。\newline 用例: This helper wraps the old library with a simpler interface.\newline 訳: この補助機能は、古いライブラリをより簡単な窓口で使えるように包みます。 \\
\addlinespace[\vocabrowspace]
\en{backend}\newline{\ipa \ipastrut /b\ipastress{vkey}{ˈ}{æ}{0.341750}{0.192500}k\ipastress{vhead}{ˌ}{ɛ}{0.115750}{0.701500}nd/} & \pos{名詞・形容詞}\newline システムの裏側でデータを保存・処理する部分／その部分の & 画面の裏側で、保存や業務処理を担う部分や、その部分に関するもの。\newline 用例: The backend stores reservations and checks available meeting rooms.\newline 訳: バックエンドは予約を保存し、利用できる会議室を確認します。\newline 用例: A backend update improved how quickly the catalog loads.\newline 訳: バックエンド側の更新で、商品一覧の読み込みが速くなりました。 \\
\addlinespace[\vocabrowspace]
\en{secret}\newline{\ipa \ipastrut /s\ipastress{vkey}{ˈ}{i}{-0.013750}{0.392500}krət/} & \pos{名詞}\newline 秘密・外部に知られてはいけない情報 & アクセスを制限し、外部へ漏らさず扱う必要がある情報。\newline 用例: The deployment system retrieves each secret from a protected vault.\newline 訳: 配備システムは、保護された保管領域から各秘密情報を取得します。 \\
\end{vtab}
\Needspace{14\baselineskip}

\section{7. Claude Codeの設定と開発}
\begin{vtab}
\en{non-interactive}\newline{\ipa \ipastrut /n\ipastress{vhead}{ˌ}{ɑ}{0.146500}{0.701000}n\ipastress{vhead}{ˌ}{ɪ}{0.032500}{0.688000}ntər\ipastress{vkey}{ˈ}{æ}{0.341750}{0.192500}ktɪv/} & \pos{形容詞}\newline 対話の入力待ちをしない & 途中で利用者に入力を求めず、そのまま処理を進める。\newline 用例: The non-interactive installer runs overnight without waiting for keyboard input.\newline 訳: 非対話型インストーラーは、キー入力を待たずに夜間実行されます。 \\
\addlinespace[\vocabrowspace]
\en{refactoring}\newline{\ipa \ipastrut /rif\ipastress{vkey}{ˈ}{æ}{0.341750}{0.192500}ktərɪŋ/} & \pos{名詞}\newline コードの内部構造を整理・改善する作業 & 外から見える動作を保ちながら、コード内部の構成を整える作業。\newline 用例: The refactoring made the payment module easier to maintain.\newline 訳: リファクタリングにより、決済モジュールの保守が容易になりました。 \\
\addlinespace[\vocabrowspace]
\en{diff}\newline{\ipa \ipastrut /dɪf/} & \pos{名詞}\newline 変更前後の差分 & 二つの版を比べ、追加・削除・変更された部分を示したもの。\newline 用例: The diff highlights the lines changed in yesterday's update.\newline 訳: 差分には、昨日の更新で変わった行が強調表示されています。 \\
\addlinespace[\vocabrowspace]
\en{regression testing}\newline{\ipa \ipastrut /rɪɡr\ipastress{vkey}{ˈ}{ɛ}{0.115750}{0.193000}ʃən t\ipastress{vkey}{ˈ}{ɛ}{0.115750}{0.193000}stɪŋ/} & \pos{名詞句}\newline 変更で既存の動作が壊れていないか確認するテスト & 変更後も、それまで動いていた機能が正常に動くか確かめること。\newline 用例: Regression testing caught a sorting problem after the database upgrade.\newline 訳: 回帰テストで、データベース更新後の並べ替えの問題を検出しました。 \\
\addlinespace[\vocabrowspace]
\en{checkpoint}\newline{\ipa \ipastrut /tʃ\ipastress{vkey}{ˈ}{ɛ}{0.115750}{0.193000}kp\ipastress{vhead}{ˌ}{ɔɪ}{0.310000}{0.701000}nt/} & \pos{名詞}\newline 戻れるように保存した作業時点 & 後で戻れるように保存しておく、作業や処理の途中の状態。\newline 用例: We saved a checkpoint before changing the simulation settings.\newline 訳: シミュレーションの設定を変える前に、途中の状態を保存しました。 \\
\addlinespace[\vocabrowspace]
\en{revert}\newline{\ipa \ipastrut /rɪv\ipastress{vkey}{ˈ}{ɜ}{0.115750}{0.193000}rt/} & \pos{自動詞}\newline 以前の状態に戻る・戻す & 新しい状態から、以前の設定や状態へ戻る。\newline 用例: The display reverts to its previous layout after restarting.\newline 訳: 再起動すると、表示は以前の配置に戻ります。 \\
\addlinespace[\vocabrowspace]
\en{enforced configuration}\newline{\ipa \ipastrut /ɪnf\ipastress{vkey}{ˈ}{ɔ}{0.125500}{0.192500}rst kənf\ipastress{vhead}{ˌ}{ɪ}{0.032500}{0.688000}ɡjər\ipastress{vkey}{ˈ}{eɪ}{0.339250}{0.192500}ʃən/} & \pos{名詞句}\newline 仕組みによって強制的に適用される設定 & 利用者の任意の選択に任せず、仕組みで適用を確保する設定。\newline 用例: The enforced configuration prevents extensions from accessing private folders.\newline 訳: 強制適用される設定により、拡張機能は非公開フォルダーにアクセスできません。 \\
\addlinespace[\vocabrowspace]
\en{convention}\newline{\ipa \ipastrut /kənv\ipastress{vkey}{ˈ}{ɛ}{0.115750}{0.193000}nʃən/} & \pos{名詞}\newline 共通の取り決め・作法 & 共同作業で守る、共通の書き方や進め方の取り決め。\newline 用例: Our naming convention keeps archived photographs easy to find.\newline 訳: 名前の付け方を統一しているので、保存した写真を探しやすくなっています。 \\
\addlinespace[\vocabrowspace]
\en{sign-off}\newline{\ipa \ipastrut /s\ipastress{vkey}{ˈ}{aɪ}{0.340750}{0.192500}n\ipastress{vhead}{ˌ}{ɔ}{0.125500}{0.701000}f/} & \pos{名詞}\newline 先へ進むことを正式に認める承認 & 確認を終え、公開や実施などの次の段階を正式に認めること。\newline 用例: We need the curator's sign-off before printing the exhibition catalog.\newline 訳: 展覧会図録を印刷する前に、学芸員の正式な承認が必要です。 \\
\addlinespace[\vocabrowspace]
\en{unattended}\newline{\ipa \ipastrut /\ipastress{vhead}{ˌ}{ʌ}{0.142750}{0.688000}nət\ipastress{vkey}{ˈ}{ɛ}{0.115750}{0.193000}ndɪd/} & \pos{形容詞}\newline 人が付き添って監視していない & 人がそばで見守ったり、その都度操作したりしていない。\newline 用例: The unattended backup completed successfully while the office was closed.\newline 訳: オフィスの閉まっている間に、無人でのバックアップが正常に完了しました。 \\
\addlinespace[\vocabrowspace]
\en{classifier}\newline{\ipa \ipastrut /kl\ipastress{vkey}{ˈ}{æ}{0.341750}{0.192500}səf\ipastress{vhead}{ˌ}{aɪ}{0.340750}{0.701000}ər/} & \pos{名詞}\newline 分類に基づいて判定する仕組み・モデル & 入力を調べ、定められた種類や区分へ振り分ける仕組み。\newline 用例: The classifier sorts recordings by instrument rather than file size.\newline 訳: その分類器は、ファイルの大きさではなく楽器別に録音を分けます。 \\
\addlinespace[\vocabrowspace]
\en{correctness}\newline{\ipa \ipastrut /kər\ipastress{vkey}{ˈ}{ɛ}{0.115750}{0.193000}ktnəs/} & \pos{名詞}\newline 正しく動くこと・正しさ & 結果や動作が、求められる条件や事実に合っていること。\newline 用例: The proof establishes the correctness of the sorting algorithm.\newline 訳: その証明により、並べ替えアルゴリズムの正しさが示されます。 \\
\addlinespace[\vocabrowspace]
\en{severity}\newline{\ipa \ipastrut /səv\ipastress{vkey}{ˈ}{ɛ}{0.115750}{0.193000}rəti/} & \pos{名詞}\newline 問題の深刻さ・重大さ & 問題の影響がどの程度大きいかという、深刻さの度合い。\newline 用例: The report ranks defects by severity and repair cost.\newline 訳: 報告書は、不具合を重大さと修理費用で順位付けしています。 \\
\addlinespace[\vocabrowspace]
\en{privileges}\newline{\ipa \ipastrut /pr\ipastress{vkey}{ˈ}{ɪ}{0.032500}{0.179500}vəlɪdʒɪz/} & \pos{名詞（複数形）}\newline 操作やアクセスを行うための権限 & システム上で、特定の操作や情報利用を行うための権限。\newline 用例: Temporary privileges allow volunteers to update the event calendar.\newline 訳: 一時的な権限により、ボランティアは催し物の予定表を更新できます。 \\
\addlinespace[\vocabrowspace]
\en{namespaced}\newline{\ipa \ipastrut /n\ipastress{vkey}{ˈ}{eɪ}{0.339250}{0.192500}msp\ipastress{vhead}{ˌ}{eɪ}{0.339250}{0.701000}st/} & \pos{過去分詞（受け身）}\newline 名前が衝突しないよう区分名を付けられた & 同じ名前がぶつからないよう、所属する区分を名前に付ける。\newline 用例: Each extension's settings are namespaced to avoid accidental conflicts.\newline 訳: 意図しない衝突を防ぐため、各拡張機能の設定名には区分が付いています。 \\
\addlinespace[\vocabrowspace]
\en{manifest}\newline{\ipa \ipastrut /m\ipastress{vkey}{ˈ}{æ}{0.341750}{0.192500}nəf\ipastress{vhead}{ˌ}{ɛ}{0.115750}{0.701500}st/} & \pos{名詞}\newline 構成や属性を宣言する情報ファイル & プログラムや配布物の構成要素・属性などを記載するファイル。\newline 用例: The package manifest lists supported platforms and required libraries.\newline 訳: パッケージのマニフェストには、対応環境と必要なライブラリが記載されています。 \\
\addlinespace[\vocabrowspace]
\en{test suite}\newline{\ipa \ipastrut /tɛst swit/} & \pos{名詞句}\newline 関連するテストをまとめた一式 & 関連する動作をまとめて確かめるために用意した、テストの一式。\newline 用例: The test suite checks keyboard navigation across every dialog.\newline 訳: そのテスト一式は、すべてのダイアログでキーボード操作を確認します。 \\
\addlinespace[\vocabrowspace]
\en{modular}\newline{\ipa \ipastrut /m\ipastress{vkey}{ˈ}{ɑ}{0.146500}{0.192500}dʒələr/} & \pos{形容詞}\newline 役割ごとに分けた部品から成る & 独立した役割を持つ部品を組み合わせて、全体を構成する。\newline 用例: The modular design lets us replace the sensor independently.\newline 訳: モジュール式の設計により、センサーだけを交換できます。 \\
\addlinespace[\vocabrowspace]
\en{monolithic}\newline{\ipa \ipastrut /m\ipastress{vhead}{ˌ}{ɑ}{0.146500}{0.701000}nəl\ipastress{vkey}{ˈ}{ɪ}{0.032500}{0.179500}θɪk/} & \pos{形容詞}\newline 多数の機能が一つの大きなまとまりになった & 多くの機能や部品が、分割されず一つの大きな単位になっている。\newline 用例: The monolithic application handles scheduling, storage, and billing together.\newline 訳: その一体型アプリケーションは、日程管理・保存・請求をまとめて処理します。 \\
\addlinespace[\vocabrowspace]
\en{technical debt}\newline{\ipa \ipastrut /t\ipastress{vkey}{ˈ}{ɛ}{0.115750}{0.193000}knɪkəl dɛt/} & \pos{名詞句}\newline 将来の修正や保守の負担として残る技術上の問題 & 目先の実装上の都合が、後の変更や保守を難しくしている状態。\newline 用例: Repeated shortcuts accumulated technical debt in the reporting module.\newline 訳: その場しのぎの対応を繰り返し、報告モジュールに技術的負債が積み重なりました。 \\
\addlinespace[\vocabrowspace]
\en{dependency audit}\newline{\ipa \ipastrut /dɪp\ipastress{vkey}{ˈ}{ɛ}{0.115750}{0.193000}ndənsi \ipastress{vkey}{ˈ}{ɔ}{0.125500}{0.192500}dɪt/} & \pos{名詞句}\newline 利用している依存ライブラリなどの状態を点検すること & 使用しているライブラリなどの依存先を調べ、状態や問題を確認すること。\newline 用例: A dependency audit identified several libraries that needed updates.\newline 訳: 依存先の点検により、更新が必要なライブラリがいくつか見つかりました。 \\
\addlinespace[\vocabrowspace]
\en{resumption}\newline{\ipa \ipastrut /rɪz\ipastress{vkey}{ˈ}{ʌ}{0.142750}{0.179500}mpʃən/} & \pos{名詞}\newline 中断していた会話や処理を再開すること & 途中で止めた作業や会話を、再び動かし始めること。\newline 用例: The application supports resumption after an interrupted file transfer.\newline 訳: そのアプリケーションは、中断したファイル転送の再開に対応しています。 \\
\addlinespace[\vocabrowspace]
\en{codebase}\newline{\ipa \ipastrut /k\ipastress{vkey}{ˈ}{oʊ}{0.462750}{0.192500}db\ipastress{vhead}{ˌ}{eɪ}{0.339250}{0.701000}s/} & \pos{名詞}\newline 一つのソフトウェアを構成するコード全体 & 一つの製品やプロジェクトを動かしている、コードの集合全体。\newline 用例: The new maintainer spent a week exploring the codebase.\newline 訳: 新しい保守担当者は、一週間かけてコード全体を調べました。 \\
\addlinespace[\vocabrowspace]
\en{boilerplate}\newline{\ipa \ipastrut /b\ipastress{vkey}{ˈ}{ɔɪ}{0.310000}{0.192500}lərpleɪt/} & \pos{名詞}\newline 繰り返し使う定型的なコードや文章 & 複数の場所で似た形のまま繰り返し使う、定型的な部分。\newline 用例: The generator removes much of the boilerplate from new modules.\newline 訳: その生成ツールにより、新しいモジュールで定型コードの多くを省けます。 \\
\addlinespace[\vocabrowspace]
\en{production}\newline{\ipa \ipastrut /prəd\ipastress{vkey}{ˈ}{ʌ}{0.142750}{0.179500}kʃən/} & \pos{名詞}\newline 本番環境・本番運用 & 試験用の環境ではなく、利用者が実際の業務などに使う環境。\newline 用例: The team tested the change before deploying it to production.\newline 訳: チームは、その変更を本番環境へ配備する前にテストしました。 \\
\addlinespace[\vocabrowspace]
\en{deprecation}\newline{\ipa \ipastrut /d\ipastress{vhead}{ˌ}{ɛ}{0.115750}{0.701500}prək\ipastress{vkey}{ˈ}{eɪ}{0.339250}{0.192500}ʃən/} & \pos{名詞}\newline 非推奨化・使用を勧めなくすること & 古い機能などの利用を勧めず、代わりの方法へ移行を促すこと。\newline 用例: The deprecation notice recommends a newer image conversion method.\newline 訳: その非推奨化の通知は、より新しい画像変換方式を勧めています。 \\
\addlinespace[\vocabrowspace]
\en{abstraction}\newline{\ipa \ipastrut /æbstr\ipastress{vkey}{ˈ}{æ}{0.341750}{0.192500}kʃən/} & \pos{名詞}\newline 抽象化・細部を隠して扱いやすくした設計 & 共通の性質を取り出し、内部の複雑さを意識せず扱えるようにすること。\newline 用例: This abstraction hides the differences between several storage providers.\newline 訳: この抽象化により、複数の保存先サービスの違いを隠せます。 \\
\addlinespace[\vocabrowspace]
\en{lint}\newline{\ipa \ipastrut /lɪnt/} & \pos{名詞（修飾用法）}\newline コードの問題や書き方を規則に照らして調べること & コードを規則に照らし、問題になりそうな記述や書き方を調べること。\newline 用例: The lint rules flag unused variables before code reaches review.\newline 訳: lintの規則は、レビューに進む前に未使用の変数を指摘します。 \\
\addlinespace[\vocabrowspace]
\en{stack trace}\newline{\ipa \ipastrut /stæk treɪs/} & \pos{名詞句}\newline 関数呼び出しのつながりを示す一覧 & エラーなどが起きた時点で、関数呼び出しの経路を示す記録。\newline 用例: The stack trace led us to the failing image decoder.\newline 訳: スタックトレースをたどり、失敗している画像デコーダーに行き着きました。 \\
\addlinespace[\vocabrowspace]
\en{off-by-one error}\newline{\ipa \ipastrut /ɑf baɪ wʌn \ipastress{vkey}{ˈ}{ɛ}{0.115750}{0.193000}rər/} & \pos{名詞句}\newline 数え方や範囲の境界が1つずれる誤り & 開始位置や終了条件などの数え方が、一つずれる誤り。\newline 用例: An off-by-one error omitted the final character from each label.\newline 訳: 一つ分のずれがある誤りにより、各ラベルの最後の文字が欠けました。 \\
\addlinespace[\vocabrowspace]
\en{auto-modernize}\newline{\ipa \ipastrut /\ipastress{vhead}{ˌ}{ɑ}{0.146500}{0.701000}toʊm\ipastress{vkey}{ˈ}{ɑ}{0.146500}{0.192500}dərnaɪz/} & \pos{他動詞}\newline …を自動的に新しい方式へ改める & 古い書き方や仕組みを、ツールによって新しい方式へ変える。\newline 用例: The utility can auto-modernize outdated import statements across the package.\newline 訳: そのユーティリティーは、パッケージ内の古いimport文を自動で新しい書き方に変更できます。 \\
\addlinespace[\vocabrowspace]
\en{cross-file}\newline{\ipa \ipastrut /krɑs faɪl/} & \pos{形容詞}\newline 複数ファイルにまたがる & 一つのファイル内に収まらず、複数のファイルへ広がっている。\newline 用例: Cross-file references connect the interface definition to its implementations.\newline 訳: ファイルをまたぐ参照が、インターフェース定義とその実装を結び付けています。 \\
\addlinespace[\vocabrowspace]
\en{fixture library}\newline{\ipa \ipastrut /f\ipastress{vkey}{ˈ}{ɪ}{0.032500}{0.179500}kstʃər l\ipastress{vkey}{ˈ}{aɪ}{0.340750}{0.192500}br\ipastress{vhead}{ˌ}{ɛ}{0.115750}{0.701500}ri/} & \pos{名詞句}\newline テスト用データや準備処理を再利用するためのライブラリ & テストで共通利用するデータや、事前準備の処理をまとめたもの。\newline 用例: Our fixture library provides sample calendars for the scheduling tests.\newline 訳: 私たちのテスト用ライブラリは、日程管理のテストにサンプルの予定表を提供します。 \\
\addlinespace[\vocabrowspace]
\en{migration}\newline{\ipa \ipastrut /maɪɡr\ipastress{vkey}{ˈ}{eɪ}{0.339250}{0.192500}ʃən/} & \pos{名詞}\newline 移行 & データや機能などを、新しい環境や構成へ移す作業。\newline 用例: The database migration finished before the library opened Monday.\newline 訳: データベースの移行は、月曜日に図書館が開く前に終わりました。 \\
\end{vtab}
\Needspace{14\baselineskip}

\section{8. プロンプト・構造化出力・評価}
\begin{vtab}
\en{required fields}\newline{\ipa \ipastrut /rɪkw\ipastress{vkey}{ˈ}{aɪ}{0.340750}{0.192500}ərd fildz/} & \pos{名詞句}\newline 必須項目・必ず存在する必要がある項目 & 書式やデータの条件として、存在することが求められる項目。\newline 用例: The registration form marks required fields with a red asterisk.\newline 訳: 登録フォームは、必須項目に赤いアスタリスクを付けています。 \\
\addlinespace[\vocabrowspace]
\en{optional}\newline{\ipa \ipastrut /\ipastress{vkey}{ˈ}{ɑ}{0.146500}{0.192500}pʃənəl/} & \pos{形容詞}\newline 必須でない・省略可能な & 条件として必須ではなく、使わない・省くこともできる。\newline 用例: Adding a profile photograph is optional during account setup.\newline 訳: アカウント設定時に顔写真を追加するかどうかは任意です。 \\
\addlinespace[\vocabrowspace]
\en{validation}\newline{\ipa \ipastrut /v\ipastress{vhead}{ˌ}{æ}{0.341750}{0.701000}ləd\ipastress{vkey}{ˈ}{eɪ}{0.339250}{0.192500}ʃən/} & \pos{名詞}\newline 定めた条件に合っているかを検証すること & 入力や出力が、あらかじめ定めた条件を満たすか確かめること。\newline 用例: Validation rejected the uploaded file because its dates were invalid.\newline 訳: 日付が不正だったため、検証処理はアップロードされたファイルを受け付けませんでした。 \\
\addlinespace[\vocabrowspace]
\en{semantic}\newline{\ipa \ipastrut /sɪm\ipastress{vkey}{ˈ}{æ}{0.341750}{0.192500}ntɪk/} & \pos{形容詞}\newline 意味・内容に関する & 文字の並び方だけでなく、表している意味や内容に関係する。\newline 用例: The two phrases have similar spelling but different semantic roles.\newline 訳: その二つの表現はつづりが似ていますが、意味上の役割は異なります。 \\
\addlinespace[\vocabrowspace]
\en{syntactically valid}\newline{\ipa \ipastrut /sɪnt\ipastress{vkey}{ˈ}{æ}{0.341750}{0.192500}ktɪkli v\ipastress{vkey}{ˈ}{æ}{0.341750}{0.192500}lɪd/} & \pos{形容詞句}\newline 構文・書式上は正しい & 言語やデータ形式の構文規則に合い、読み取れる形になっている。\newline 用例: The query is syntactically valid and passes the parser's checks.\newline 訳: そのクエリは構文上正しく、構文解析器の検査を通過します。 \\
\addlinespace[\vocabrowspace]
\en{ground truth}\newline{\ipa \ipastrut /ɡraʊnd truθ/} & \pos{名詞句}\newline 比較・評価の基準となる正解情報 & 予測や抽出結果を評価するときに、正解として照合する情報。\newline 用例: Hand-labeled photographs provided ground truth for the image recognition experiment.\newline 訳: 人がラベルを付けた写真を、画像認識実験の正解情報として使いました。 \\
\addlinespace[\vocabrowspace]
\en{grader}\newline{\ipa \ipastrut /ɡr\ipastress{vkey}{ˈ}{eɪ}{0.339250}{0.192500}dər/} & \pos{名詞}\newline 出力を基準に照らして評価する仕組み・担当 & 回答などを基準に照らして調べ、点数や評価を付ける担当・仕組み。\newline 用例: The grader awards points for both accuracy and clear explanations.\newline 訳: 採点者は、正確さと説明の分かりやすさの両方に点を付けます。 \\
\addlinespace[\vocabrowspace]
\en{baseline score}\newline{\ipa \ipastrut /b\ipastress{vkey}{ˈ}{eɪ}{0.339250}{0.192500}sl\ipastress{vhead}{ˌ}{aɪ}{0.340750}{0.701000}n skɔr/} & \pos{名詞句}\newline 改善前の比較基準になる評価点 & 変更による効果を比べるために、最初に測っておく点数。\newline 用例: We recorded a baseline score before changing the training materials.\newline 訳: 研修資料を変更する前に、比較基準となる点数を記録しました。 \\
\addlinespace[\vocabrowspace]
\en{format compliance}\newline{\ipa \ipastrut /f\ipastress{vkey}{ˈ}{ɔ}{0.125500}{0.192500}rmæt kəmpl\ipastress{vkey}{ˈ}{aɪ}{0.340750}{0.192500}əns/} & \pos{名詞句}\newline 指定された出力形式に従っていること & 指定された書式や構成に、出力が従っていること。\newline 用例: The checker measures format compliance using the publisher's submission rules.\newline 訳: 検査ツールは、出版社の投稿規則を使って書式への適合を測定します。 \\
\addlinespace[\vocabrowspace]
\en{parsing}\newline{\ipa \ipastrut /p\ipastress{vkey}{ˈ}{ɑ}{0.146500}{0.192500}rsɪŋ/} & \pos{名詞（動名詞）・他動詞（parse）}\newline 文章・コード等を構文規則に従って読み取る処理 & 決まった規則に従って文字列を読み、構造や各部分を取り出す処理。\newline 用例: Parsing converts the raw timetable into labeled departure records.\newline 訳: 構文解析により、未加工の時刻表を項目の付いた出発記録へ変換します。\newline 用例: The script parses configuration files before launching the simulation.\newline 訳: そのスクリプトは、シミュレーション開始前に設定ファイルを構文解析します。 \\
\addlinespace[\vocabrowspace]
\en{reconcile}\newline{\ipa \ipastrut /r\ipastress{vkey}{ˈ}{ɛ}{0.115750}{0.193000}kənsaɪl/} & \pos{他動詞}\newline 食い違う内容を照合し、どう扱うか整える & 異なる記録などを突き合わせ、不一致を解消して整合を取る。\newline 用例: We reconciled the inventory records before ordering additional supplies.\newline 訳: 追加の備品を注文する前に、在庫記録を照合して整合させました。 \\
\addlinespace[\vocabrowspace]
\en{heuristic}\newline{\ipa \ipastrut /hjur\ipastress{vkey}{ˈ}{ɪ}{0.032500}{0.179500}stɪk/} & \pos{形容詞}\newline 経験則や簡便な基準に頼る & 厳密な全探索などをせず、経験則や手軽な目安で判断する。\newline 用例: A heuristic search found a useful route within seconds.\newline 訳: 経験則を使った探索で、数秒以内に使える経路が見つかりました。 \\
\addlinespace[\vocabrowspace]
\en{discrepancy}\newline{\ipa \ipastrut /dɪskr\ipastress{vkey}{ˈ}{ɛ}{0.115750}{0.193000}pənsi/} & \pos{名詞}\newline 一致するはずの値や内容の食い違い & 本来一致する値や説明の間に、ずれが見つかること。\newline 用例: We noticed a discrepancy between the diagram and the measurements.\newline 訳: 図と測定値の間に、食い違いがあることに気付きました。 \\
\addlinespace[\vocabrowspace]
\en{enum} & \pos{識別子}\newline 値として選べる候補を列挙する型・指定 & 取り得る値を、あらかじめ列挙した候補の中に限定する指定。\newline 用例: The enum restricts the direction value to four named options.\newline 訳: そのenum指定は、方向の値を名前の付いた四つの候補に限定します。 \\
\addlinespace[\vocabrowspace]
\en{edge cases}\newline{\ipa \ipastrut /\ipastress{vkey}{ˈ}{ɛ}{0.115750}{0.193000}dʒ k\ipastress{vhead}{ˌ}{eɪ}{0.339250}{0.701000}sɪz/} & \pos{名詞句（複数形）}\newline 通常の想定から外れる特殊なケース & 通常の入力では起きにくい、境界や特殊な条件に当たるケース。\newline 用例: The tests include edge cases involving empty files and unusual characters.\newline 訳: テストには、空のファイルや特殊な文字を扱う例外的なケースが含まれます。 \\
\addlinespace[\vocabrowspace]
\en{rationale}\newline{\ipa \ipastrut /r\ipastress{vhead}{ˌ}{æ}{0.341750}{0.701000}ʃən\ipastress{vkey}{ˈ}{æ}{0.341750}{0.192500}l/} & \pos{名詞}\newline 判断や行動を選ぶ理由・論拠 & なぜその判断や行動を選ぶのかを支える、理由や考え方。\newline 用例: The proposal explains the rationale for extending the opening hours.\newline 訳: 提案書は、開館時間を延長する理由を説明しています。 \\
\addlinespace[\vocabrowspace]
\en{threshold}\newline{\ipa \ipastrut /θr\ipastress{vkey}{ˈ}{ɛ}{0.115750}{0.193000}ʃhoʊld/} & \pos{名詞}\newline 判断や処理を切り替える境目の値・しきい値 & 数値がそこを越えるかどうかで、判断や動作を切り替える境目。\newline 用例: An alarm sounds when the temperature crosses the safety threshold.\newline 訳: 温度が安全上のしきい値を越えると、警報が鳴ります。 \\
\addlinespace[\vocabrowspace]
\en{reconciliation}\newline{\ipa \ipastrut /r\ipastress{vhead}{ˌ}{ɛ}{0.115750}{0.701500}kəns\ipastress{vhead}{ˌ}{ɪ}{0.032500}{0.688000}li\ipastress{vkey}{ˈ}{eɪ}{0.339250}{0.192500}ʃən/} & \pos{名詞}\newline 照合して食い違いを確かめ、整合させること & 複数の記録を照らし合わせ、差異を調べて一致させる作業。\newline 用例: Monthly reconciliation keeps the equipment register aligned with actual inventory.\newline 訳: 毎月の照合作業により、備品台帳と実際の在庫を一致させています。 \\
\addlinespace[\vocabrowspace]
\en{party}\newline{\ipa \ipastrut /p\ipastress{vkey}{ˈ}{ɑ}{0.146500}{0.192500}rti/} & \pos{名詞}\newline 当事者・関係者 & 取引や取り決めなどに参加し、権利や責任を持つ人・組織。\newline 用例: Each party must receive a copy of the signed agreement.\newline 訳: 各当事者は、署名済みの合意書の写しを受け取る必要があります。 \\
\end{vtab}
\Needspace{14\baselineskip}

\section{9. コンテキストと情報の信頼性}
\begin{vtab}
\en{scratchpad file}\newline{\ipa \ipastrut /skr\ipastress{vkey}{ˈ}{æ}{0.341750}{0.192500}tʃp\ipastress{vhead}{ˌ}{æ}{0.341750}{0.701000}d faɪl/} & \pos{名詞句}\newline 作業中の発見や進捗を残すメモファイル & 作業の途中で考えたことや調べたことを、一時的に書き留めるファイル。\newline 用例: I recorded possible titles in a scratchpad file during brainstorming.\newline 訳: 案出しの間、思い付いた題名を作業用メモファイルに記録しました。 \\
\addlinespace[\vocabrowspace]
\en{provenance}\newline{\ipa \ipastrut /pr\ipastress{vkey}{ˈ}{ɑ}{0.146500}{0.192500}vənəns/} & \pos{名詞}\newline 情報がどこから来たかという来歴・出所 & 情報や資料がどこで生まれ、どのように伝わってきたかという来歴。\newline 用例: The catalog records the provenance of every donated photograph.\newline 訳: 目録には、寄贈された写真一枚ごとの来歴が記録されています。 \\
\addlinespace[\vocabrowspace]
\en{attribution}\newline{\ipa \ipastrut /\ipastress{vhead}{ˌ}{æ}{0.341750}{0.701000}trəbj\ipastress{vkey}{ˈ}{u}{0.160750}{0.179500}ʃən/} & \pos{名詞}\newline 主張をその出典に結び付けること & 使った情報や主張を、その元になった資料へ結び付けて示すこと。\newline 用例: Clear attribution lets readers locate the source of each statistic.\newline 訳: 出典を明示すると、読者は各統計値の元資料を探せます。 \\
\addlinespace[\vocabrowspace]
\en{citations}\newline{\ipa \ipastrut /saɪt\ipastress{vkey}{ˈ}{eɪ}{0.339250}{0.192500}ʃənz/} & \pos{名詞（複数形）}\newline 引用・出典を示す記載 & 文章中の記述がどの資料に基づくかを示す、引用・出典の表記。\newline 用例: The article includes citations for all historical population estimates.\newline 訳: その記事には、過去の人口推計値すべての出典が記載されています。 \\
\addlinespace[\vocabrowspace]
\en{retrieval}\newline{\ipa \ipastrut /rɪtr\ipastress{vkey}{ˈ}{i}{-0.013750}{0.392500}vəl/} & \pos{名詞}\newline 必要な資料や情報を取り出すこと & 保存先や外部の資料から、必要な情報を探して取り出すこと。\newline 用例: An index speeds up retrieval of older newspaper articles.\newline 訳: 索引により、古い新聞記事を取り出す処理が速くなります。 \\
\addlinespace[\vocabrowspace]
\en{chunking}\newline{\ipa \ipastrut /tʃ\ipastress{vkey}{ˈ}{ʌ}{0.142750}{0.179500}ŋkɪŋ/} & \pos{名詞（動名詞）}\newline 大きな内容を小さなまとまりに分けること & 大きな文章やデータを、扱いやすい小さな単位に分割すること。\newline 用例: Chunking the audio made transcription easier on the small device.\newline 訳: 音声を小分けにすると、小型端末で文字起こししやすくなりました。 \\
\addlinespace[\vocabrowspace]
\en{context budget}\newline{\ipa \ipastrut /k\ipastress{vkey}{ˈ}{ɑ}{0.146500}{0.192500}ntɛkst b\ipastress{vkey}{ˈ}{ʌ}{0.142750}{0.179500}dʒɪt/} & \pos{名詞句}\newline 参照情報に使える容量の割り当て & モデルが参照する指示・資料・履歴などに使える容量の配分。\newline 用例: The retrieved chapters used most of the available context budget.\newline 訳: 取得した章が、利用できるコンテキスト容量の大半を使いました。 \\
\addlinespace[\vocabrowspace]
\en{truncation}\newline{\ipa \ipastrut /trʌŋk\ipastress{vkey}{ˈ}{eɪ}{0.339250}{0.192500}ʃən/} & \pos{名詞}\newline 上限などによる途中での切り落とし & 長さや容量の上限により、続きが途中で切り落とされること。\newline 用例: Truncation removed the last paragraph from the displayed message.\newline 訳: 途中での切り落としにより、表示されたメッセージの最終段落が消えました。 \\
\addlinespace[\vocabrowspace]
\en{stale}\newline{\ipa \ipastrut /steɪl/} & \pos{形容詞}\newline 古くなり、現在の状況に合わない & 以前の情報が更新されず、今の状態と合わなくなっている。\newline 用例: The dashboard showed stale figures until we refreshed its data.\newline 訳: データを更新するまで、ダッシュボードには古い数値が表示されていました。 \\
\addlinespace[\vocabrowspace]
\en{recency}\newline{\ipa \ipastrut /r\ipastress{vkey}{ˈ}{i}{-0.013750}{0.392500}sənsi/} & \pos{名詞}\newline どれだけ最近のものであるか & 情報や出来事が、現在からどれほど近い時期のものかという新しさ。\newline 用例: The search interface lets readers sort articles by recency.\newline 訳: その検索画面では、読者は記事を新しい順に並べ替えられます。 \\
\addlinespace[\vocabrowspace]
\en{lost-in-the-middle effect}\newline{\ipa \ipastrut /lɔst ɪn ðə m\ipastress{vkey}{ˈ}{ɪ}{0.032500}{0.179500}dəl ɪf\ipastress{vkey}{ˈ}{ɛ}{0.115750}{0.193000}kt/} & \pos{名詞句}\newline 長い入力の中ほどの情報が拾われにくくなる現象 & 長い入力の中央付近にある情報を、モデルが十分に利用しにくい現象。\newline 用例: The experiment measured the lost-in-the-middle effect across different passage lengths.\newline 訳: その実験は、文章の長さを変えて中央付近の情報が拾われにくい現象を測定しました。 \\
\addlinespace[\vocabrowspace]
\en{escalation}\newline{\ipa \ipastrut /\ipastress{vhead}{ˌ}{ɛ}{0.115750}{0.701500}skəl\ipastress{vkey}{ˈ}{eɪ}{0.339250}{0.192500}ʃən/} & \pos{名詞}\newline より適切な判断者・担当へ対応を引き上げること & 自分の担当範囲で扱いきれない問題を、適切な担当や上位判断者へ回すこと。\newline 用例: The handbook explains the escalation process for urgent equipment faults.\newline 訳: 手引きには、緊急の機器故障を適切な担当へ引き上げる手順が書かれています。 \\
\addlinespace[\vocabrowspace]
\en{source-grounded}\newline{\ipa \ipastrut /sɔrs ɡr\ipastress{vkey}{ˈ}{aʊ}{0.465750}{0.192500}ndɪd/} & \pos{形容詞}\newline 元の資料を根拠にした & 元の資料に書かれた内容を根拠として、説明や判断を行う。\newline 用例: A source-grounded explanation links its main points to original records.\newline 訳: 元資料に基づく説明は、主要な点を原記録に結び付けます。 \\
\addlinespace[\vocabrowspace]
\en{inherit}\newline{\ipa \ipastrut /ɪnh\ipastress{vkey}{ˈ}{ɛ}{0.115750}{0.193000}rɪt/} & \pos{他動詞}\newline 親や先行するものから情報・設定を引き継ぐ & 上位や前段にあるものの性質・情報・設定を受け継ぐ。\newline 用例: New folders inherit the access settings of their parent directory.\newline 訳: 新しいフォルダーは、親ディレクトリーのアクセス設定を引き継ぎます。 \\
\addlinespace[\vocabrowspace]
\en{verbose}\newline{\ipa \ipastrut /vərb\ipastress{vkey}{ˈ}{oʊ}{0.462750}{0.192500}s/} & \pos{形容詞}\newline 出力や説明の量が多い・冗長な & 伝える内容に対して、説明や出力の量が多い。\newline 用例: The verbose status message overwhelmed the small notification panel.\newline 訳: 長すぎる状態通知が、小さな通知欄を埋め尽くしました。 \\
\addlinespace[\vocabrowspace]
\en{persistent}\newline{\ipa \ipastrut /pərs\ipastress{vkey}{ˈ}{ɪ}{0.032500}{0.179500}stənt/} & \pos{形容詞}\newline 後の処理やセッションでも使えるよう保持される & 処理や利用を終えても消えず、後で使えるよう保持される。\newline 用例: Persistent storage preserves your preferences after the application closes.\newline 訳: 永続的な保存領域は、アプリケーションを閉じた後も設定を保持します。 \\
\addlinespace[\vocabrowspace]
\en{transcript}\newline{\ipa \ipastrut /tr\ipastress{vkey}{ˈ}{æ}{0.341750}{0.192500}nskrɪpt/} & \pos{名詞（可算）}\newline 会話・やり取りを文字で残した記録 & 音声や文字で行われた会話を、文字の形で残した記録。\newline 用例: The researcher checked the interview transcript against the recording.\newline 訳: 研究者は、面談の文字起こしを録音と照合しました。\newline 用例: The chat transcript helped us recover the missing workshop address.\newline 訳: チャットの記録のおかげで、分からなくなった講習会場の住所を見つけられました。 \\
\addlinespace[\vocabrowspace]
\en{crawl}\newline{\ipa \ipastrut /krɔl/} & \pos{名詞}\newline Webページを自動で巡回・取得する処理 & ウェブサイトのページを自動的にたどり、内容を取得する処理。\newline 用例: The overnight crawl collected updated pages from the public archive.\newline 訳: 夜間の巡回取得で、公開資料庫から更新されたページを集めました。 \\
\addlinespace[\vocabrowspace]
\en{author-year}\newline{\ipa \ipastrut /\ipastress{vkey}{ˈ}{ɑ}{0.146500}{0.192500}θər jɪr/} & \pos{複合形容詞}\newline 著者名と出版年で出典を示す形式の & 著者の名前と出版年を使い、参照した資料を示す形式。\newline 用例: The journal requires author-year references throughout the main text.\newline 訳: その雑誌は、本文全体で著者名と出版年による出典表記を求めています。 \\
\addlinespace[\vocabrowspace]
\en{bibliography}\newline{\ipa \ipastrut /b\ipastress{vhead}{ˌ}{ɪ}{0.032500}{0.688000}bli\ipastress{vkey}{ˈ}{ɑ}{0.146500}{0.192500}ɡrəfi/} & \pos{名詞}\newline 参考文献一覧・文献目録 & 文書の作成時に参照した本や論文などをまとめた一覧。\newline 用例: The bibliography includes books, journal articles, and several archival letters.\newline 訳: 参考文献一覧には、本、学術論文、保存されていた手紙が含まれます。 \\
\addlinespace[\vocabrowspace]
\en{cache}\newline{\ipa \ipastrut /kæʃ/} & \pos{名詞}\newline 再利用のための一時的な保存領域・保存データ & 一度取得したデータを保存し、次回の取得や計算を減らす仕組み。\newline 用例: The browser cache stores images for quicker loading later.\newline 訳: ブラウザーのキャッシュは、後で素早く読み込むために画像を保存します。 \\
\end{vtab}
\Needspace{14\baselineskip}

\section{10. APIの通信と処理}
\begin{vtab}
\en{stateless}\newline{\ipa \ipastrut /st\ipastress{vkey}{ˈ}{eɪ}{0.339250}{0.192500}tləs/} & \pos{形容詞}\newline 前の処理の状態を自動で持ち越さない & 前の要求や処理の情報を、次の要求へ自動では持ち越さない。\newline 用例: The stateless service treats each incoming request as independent.\newline 訳: そのステートレスなサービスは、届く各要求を独立したものとして扱います。 \\
\addlinespace[\vocabrowspace]
\en{response streaming}\newline{\ipa \ipastrut /rɪsp\ipastress{vkey}{ˈ}{ɑ}{0.146500}{0.192500}ns str\ipastress{vkey}{ˈ}{i}{-0.013750}{0.392500}mɪŋ/} & \pos{名詞句}\newline 応答を生成途中から少しずつ送ること & 応答の完成を待たず、できた部分から順に送り出す方式。\newline 用例: Response streaming displays the opening sentence before generation finishes.\newline 訳: 応答のストリーミングにより、生成が終わる前に最初の文が表示されます。 \\
\addlinespace[\vocabrowspace]
\en{cache breakpoints}\newline{\ipa \ipastrut /kæʃ br\ipastress{vkey}{ˈ}{eɪ}{0.339250}{0.192500}kp\ipastress{vhead}{ˌ}{ɔɪ}{0.310000}{0.701000}nts/} & \pos{名詞句}\newline キャッシュする入力の範囲を区切る指定位置 & 入力のどこまでを再利用する範囲にするかを示す位置。\newline 用例: We placed cache breakpoints after the reusable instruction blocks.\newline 訳: 再利用する指示ブロックの後ろに、キャッシュ範囲の区切りを置きました。 \\
\addlinespace[\vocabrowspace]
\en{latency}\newline{\ipa \ipastrut /l\ipastress{vkey}{ˈ}{eɪ}{0.339250}{0.192500}tənsi/} & \pos{名詞}\newline 処理を待つ時間・応答までの遅延 & 処理を依頼してから、応答や必要な結果を得るまでの待ち時間。\newline 用例: The closer server reduced latency during remote drawing sessions.\newline 訳: 近くのサーバーを使うことで、遠隔描画の作業中の遅延が減りました。 \\
\addlinespace[\vocabrowspace]
\en{serialized}\newline{\ipa \ipastrut /s\ipastress{vkey}{ˈ}{ɪ}{0.032500}{0.179500}riəl\ipastress{vhead}{ˌ}{aɪ}{0.340750}{0.701000}zd/} & \pos{過去分詞}\newline 保存・受け渡しできる表現へ変換された & データを保存や送信に使える文字列・バイト列などへ変換した状態。\newline 用例: The serialized settings travel with the exported project file.\newline 訳: 変換済みの設定データは、書き出したプロジェクトファイルと一緒に渡されます。 \\
\addlinespace[\vocabrowspace]
\en{propagate}\newline{\ipa \ipastrut /pr\ipastress{vkey}{ˈ}{ɑ}{0.146500}{0.192500}pəɡ\ipastress{vhead}{ˌ}{eɪ}{0.339250}{0.701000}t/} & \pos{自動詞・他動詞}\newline エラーや情報が別の処理へ伝わる・伝える & 変更や情報が別の部分へ広がること。また、それらを伝えること。\newline 用例: Updates propagate through the shared catalog within a few minutes.\newline 訳: 更新は、数分以内に共有の商品一覧全体へ伝わります。\newline 用例: The service propagates schedule changes to all connected displays.\newline 訳: そのサービスは、予定の変更を接続されたすべての表示装置へ伝えます。 \\
\addlinespace[\vocabrowspace]
\en{error propagation}\newline{\ipa \ipastrut /\ipastress{vkey}{ˈ}{ɛ}{0.115750}{0.193000}rər pr\ipastress{vhead}{ˌ}{ɑ}{0.146500}{0.701000}pəɡ\ipastress{vkey}{ˈ}{eɪ}{0.339250}{0.192500}ʃən/} & \pos{名詞句}\newline 発生したエラーを関連する処理や担当へ伝えること & ある場所で発生したエラーが、関連する処理や呼び出し元へ伝わること。\newline 用例: The diagram shows error propagation from storage to the interface.\newline 訳: その図は、保存処理から画面側へエラーが伝わる流れを示しています。 \\
\addlinespace[\vocabrowspace]
\en{exponential backoff}\newline{\ipa \ipastrut /\ipastress{vhead}{ˌ}{ɛ}{0.115750}{0.701500}kspoʊn\ipastress{vkey}{ˈ}{ɛ}{0.115750}{0.193000}nʃəl b\ipastress{vkey}{ˈ}{æ}{0.341750}{0.192500}k\ipastress{vhead}{ˌ}{ɔ}{0.125500}{0.701000}f/} & \pos{名詞句}\newline 再試行までの待ち時間を一定の倍率で増やす方式 & 再試行を重ねるたびに、待つ時間を一定の倍率で長くする方式。\newline 用例: Exponential backoff spaces out repeated connection attempts during an outage.\newline 訳: 指数的な待機時間の増加により、障害中の接続再試行の間隔を広げます。 \\
\addlinespace[\vocabrowspace]
\en{observability}\newline{\ipa \ipastrut /əbz\ipastress{vhead}{ˌ}{ɜ}{0.115750}{0.701500}rvəb\ipastress{vkey}{ˈ}{ɪ}{0.032500}{0.179500}lɪti/} & \pos{名詞}\newline 記録や通知から内部の動きや状態を把握できること & ログなどの外から得られる情報で、内部の状態や動きを把握できること。\newline 用例: Better observability helped engineers diagnose the intermittent processing delays.\newline 訳: 内部状態を把握しやすくしたことで、技術者は断続的な処理遅延を調べられました。 \\
\addlinespace[\vocabrowspace]
\en{round trips}\newline{\ipa \ipastrut /raʊnd trɪps/} & \pos{名詞句}\newline 要求を送り、結果が戻るまでの往復 & 要求を送ってから応答を受け取るまでを、一往復として数えたもの。\newline 用例: Combining these requests reduces round trips to the database.\newline 訳: これらの要求をまとめると、データベースとの往復が減ります。 \\
\addlinespace[\vocabrowspace]
\en{root cause}\newline{\ipa \ipastrut /rut kɔz/} & \pos{名詞句}\newline 表面の症状を生んでいる根本的な原因 & 目に見える症状の背後で、問題を引き起こしている根本の原因。\newline 用例: A loose connector was the root cause of the flickering.\newline 訳: ちらつきの根本原因は、接続部の緩みでした。 \\
\addlinespace[\vocabrowspace]
\en{fallback}\newline{\ipa \ipastrut /f\ipastress{vkey}{ˈ}{ɔ}{0.125500}{0.192500}lb\ipastress{vhead}{ˌ}{æ}{0.341750}{0.701000}k/} & \pos{名詞}\newline 通常の方法が使えない場合に用意する代替手段 & 通常の方法が使えないときに備えた、代わりの手段。\newline 用例: The printed directory serves as a fallback during network failures.\newline 訳: ネットワーク障害中は、紙の名簿を代替手段として使います。 \\
\addlinespace[\vocabrowspace]
\en{cadence}\newline{\ipa \ipastrut /k\ipastress{vkey}{ˈ}{eɪ}{0.339250}{0.192500}dəns/} & \pos{名詞}\newline 定期的に繰り返すペース・間隔 & 同じ活動や処理を、どのくらいの間隔で繰り返すかというペース。\newline 用例: The team agreed on a weekly cadence for design reviews.\newline 訳: チームは、設計レビューを毎週行うペースに決めました。 \\
\addlinespace[\vocabrowspace]
\en{call}\newline{\ipa \ipastrut /kɔl/} & \pos{名詞・他動詞}\newline 呼び出し／…を呼び出す & プログラムなどから処理を依頼することと、その依頼の一回分。\newline 用例: Each call to the service returns the current temperature.\newline 訳: そのサービスを一回呼び出すごとに、現在の温度が返されます。\newline 用例: The dashboard calls the service whenever the user refreshes it.\newline 訳: ダッシュボードは、利用者が更新するたびにそのサービスを呼び出します。 \\
\end{vtab}

\end{document}
